% concatenated from arxiv_final.tex
\documentclass[aps,pra,11pt,onecolumn,nofootinbib,tightenlines]{revtex4-2}

\usepackage{mathrsfs}
\usepackage{graphicx}
\usepackage{comment}
\usepackage{dsfont}
\usepackage{amsmath}
\usepackage{amsfonts}
\usepackage{amssymb}
\usepackage{amsthm}
\usepackage{float}
\usepackage[colorlinks=true,urlcolor=blue,citecolor=blue,linkcolor=blue]{hyperref}
\usepackage{caption}
\captionsetup{justification=raggedright}
\usepackage{lipsum}
\usepackage{subcaption}
\usepackage{mwe}
\usepackage{enumerate}
\usepackage{sidecap}
\usepackage{natbib}
\usepackage{appendix}
\usepackage{tikz}
\usepackage{booktabs}
\usepackage{array}
\usepackage{multirow}
\usepackage{forest}
\usepackage{bm}
\usepackage{empheq}
\usepackage[shortlabels]{enumitem}

\renewcommand\captionfont{\footnotesize} 
\renewcommand\captionlabelfont{\bfseries}

\setlength{\jot}{9pt}

\newtheorem{theorem}{Theorem}[section]
\newtheorem{lemma}[theorem]{Lemma}
\newtheorem{proposition}[theorem]{Proposition}
\newtheorem{corollary}[theorem]{Corollary}
\newtheorem{conjecture}[theorem]{Conjecture}

\theoremstyle{definition}
\newtheorem{definition}[theorem]{Definition}
\newtheorem{assumption}{Assumption}
\newtheorem{example}[theorem]{Example}
\newtheorem{remark}[theorem]{Remark}

\let\oldremark\remark
\renewcommand{\remark}{\oldremark\upshape}

\newcommand{\fo}{\rightarrowtail}
\newcommand{\setFO}{\mathcal{FO}}
\newcommand{\setCO}{\mathcal{CO}}
\newcommand{\setCCO}{\mathcal{CCO}}
\newcommand{\mono}{\mathfrak{R}}
\newcommand{\F}{\mathcal{F}}
\newcommand{\bigket}[1]{\Big| #1 \Big\rangle}
\newcommand{\bigbra}[1]{\Big \langle #1 \Big |}

\newcommand{\mfF}{\mathfrak{F}}
\def \d {\mathrm{d}} 
\newcommand{\D}{\mathcal{D}}
\newcommand{\N}{\mathcal{N}}
\newcommand{\M}{\mathcal{M}}
\newcommand{\qstate}{\mathcal{S}(A)}

\renewcommand{\H}{\mathbb{H}}
\newcommand{\cE}{\mathcal{E}}
\newcommand{\bbN}{\mathbb{N}}
\newcommand{\cS}{\mathcal{S}} 
\newcommand{\W}{\mathcal{W}} 
\newcommand{\X}{\mathcal{X}} 
\newcommand{\Prob}{\mathcal{P}} 
\newcommand{\V}{\mathcal{V}} 
\newcommand{\Y}{\mathcal{Y}} 
\newcommand{\cH}{\mathbb{H}}
\newcommand{\cF}{\mathcal{F}} 

\newcommand{\id}{I}

\newcommand{\R}{\mathbb{R}}
\newcommand{\T}{\mathbb{T}}


\newcommand{\RR}[1]{\textcolor{red}{(RR: #1)}}

\DeclareMathOperator*{\argmin}{arg\,min}
\DeclareMathOperator*{\argmax}{arg\,max}
\DeclareMathOperator*{\opt}{opt}
\DeclareMathOperator*{\Trm}{Tr}
\DeclareMathOperator*{\Op}{Op \!\!\,}

\def\cE{\mathcal E}

\definecolor{darker}{RGB}{2.0,100.0,20.0}

%\newcommand{\RR}[1]{\textcolor{red}{(RR: #1)}}
\newcommand{\MT}[1]{\textcolor{purple}{(MT: #1)}}
\newcommand{\EH}[1]{\textcolor{darker}{(EH: #1)}}

\renewcommand{\thesection}{\arabic{section}}

%%%%%Erkka's definitions
%fonts
\newcommand{\mc}[1]{\mathcal{#1}}
\newcommand{\ms}[1]{\mathsf{#1}}
\newcommand{\mf}[1]{\mathfrak{#1}}
\newcommand{\mb}[1]{\mathbb{#1}}
%\newcommand{\<}{\langle}
\newcommand{\ketbra}[2]{\left|#1\middle\rangle\middle\langle #2\right|}
\newcommand{\ptr}[2]{{\rm Tr}_{#1}\left[#2\right]}

% Semiring preorders
\newcommand{\rleq}{\trianglelefteq}
\newcommand{\rgeq}{\trianglerighteq}


\newenvironment{coloredtext}[1]{\begingroup\color{#1}}{\par\endgroup}

\begin{document}

\title{\Large A complete characterisation of conditional entropies}

\begin{abstract}
Entropies are fundamental measures of uncertainty with central importance in information theory and statistics and applications across all the quantitative sciences. Under a natural set of operational axioms, the most general form of entropy is captured by the family of Rényi entropies, parameterized by a real number $\alpha$. 
%
Conditional entropy extends the notion of entropy by quantifying uncertainty from the viewpoint of an observer with access to potentially correlated side information. However, despite their significance and the emergence of various useful definitions, a complete characterization of measures of conditional entropy that satisfy a natural set of operational axioms has remained elusive.
%
In this work, we provide a complete characterization of conditional entropy, defined through a set of axioms that are essential for any operationally meaningful definition: additivity for independent random variables, invariance under relabeling, and monotonicity under conditional mixing channels. 
%
We prove that the most general form of conditional entropy is captured by a family of measures that are exponential averages of Rényi entropies of the conditioned distribution and parameterized by a real parameter and a probability measure on the positive reals. Finally, we show that these quantities determine the rate of transformation under conditional mixing and provide a set of second laws of quantum thermodynamics with side information for states diagonal in the energy eigenbasis.


\end{abstract}



\author{Roberto Rubboli$^{*,\dagger,\ddagger}$}
\author{Erkka Haapasalo$^{*}$}
\author{Marco Tomamichel$^{*,\S}$}

\maketitle


% Footnotes with matching symbols
\begingroup
\renewcommand{\thefootnote}{\fnsymbol{footnote}}

\footnotetext[3]{Email: \texttt{ror@math.ku.dk}}       % \dagger
\footnotetext[1]{Centre for Quantum Technologies, National University of Singapore, Singapore 117543, Singapore} % *
\footnotetext[2]{Department of Mathematical Sciences, University of Copenhagen, Universitetsparken 5, 2100 Denmark}                         % \ddagger
\footnotetext[4]{Department of Electrical and Computer Engineering,
National University of Singapore, Singapore 117583, Singapore}                         % \S

\endgroup

\tableofcontents 

\section{Introduction}

Entropies are measures of uncertainty about the outcome of a random experiment. Given a discrete random variable $X$ on $\mc{X}$ governed by a probability mass function (pmf) $P(x)$, its surprisal is the random variable $S(X) = \log \frac{1}{P(X)}$. The Shannon entropy~\cite{shannon1948mathematical} is its expectation\footnote{Unless stated otherwise, all sums are understood to run over the full domain.},
\begin{equation}
    H(X) := \mb{E}[S] = \sum_{x} P(x) \log \frac{1}{P(x)} \,.
\end{equation}

While Shannon entropy plays a central role in information processing by characterizing the ultimate performance limits for a range of tasks, most notably data compression~\cite{shannon1948mathematical}, it is evident that a more fine-grained specification of the surprisal is necessary to understand tail events. Notable examples are the min-entropy, $H_{\min}(X) = \min_{x \in \mc{X}} S(x)$, which plays a pivotal role in cryptography, and the Hartley entropy~\cite{hartley28}, $H_{\max}(X) = \log \sum_{x \in \mc{X}} \mathbf{1}\{ P(x) > 0\}$. To fully characterize the distribution we consider the family of exponential expectations,
\begin{align}
  \mathbb{R} \ni  t \mapsto \mb{E}_t[X] := \frac{1}{t} \log \mb{E}\big[ \exp(t X) \big], 
\end{align}
and observe that $\lim_{t \to 0} \mb{E}_t[X] = \mb{E}[X]$ recovers the usual expectation and the limits $t \to \pm \infty$ recover the maximal and minimal value in $X$ can take, respectively. This function is closely related to the cumulant generation function (cgf) given by $K_t[X] = t\,\mb{E}_t[X]$.

Given this, the distribution of the surprisal $S$ is conveniently characterised in terms of the family of Rényi entropies~\cite{renyi61} with parameter $\alpha \in [0, \infty]$ as
\begin{align}
    H_{\alpha}(X) := \frac{1}{1-\alpha} \log \left( \sum_{x} P(x)^{\alpha} \right) = \mb{E}_{1-\alpha}[S] \,,
\end{align}
The Hartley, Shannon, and min-entropy emerge as point-wise limits when $\alpha \to \{0, 1, \infty\}$, respectively. 
The cgf and Rényi entropies arise naturally whenever tails of the distribution of the surprisal are relevant, e.g., in the large deviation regime of source coding, when investigating error or strong converse exponents, and in one-shot information theory.

Their ubiquity in information theory leads to the question of whether Rényi entropies are complete under some set of well-motivated information-theoretic axioms. Axiomatic approaches to entropy have a rich history, surveyed in the books by Acz\'el-Dar\'oczy~\cite{aczel75} and Ebanks-Sahoo-Sander~\cite{ebanks98}, as well as the review by Csisz\'ar~\cite{csiszar08}. All these approaches have in common that at least one of the axioms is chosen for mathematical convenience and does not have an operational justification. As such, they fail to fully explain the role Rényi entropies play in information theory. 

In contrast, recent work~\cite{gour21_axiomatic} proposes a set of axioms that 
any operationally meaningful measure of entropy should satisfy: monotonicity under doubly stochastic channels, invariance under embedding, additivity for independent variables, and normalization. We shall soon explain the meaning of these concepts.
%
%Monotonicity under doubly stochastic channels---also known as mixing channels---ensures that entropy does not decrease when additional randomness is introduced through a random permutation. Invariance under embedding guarantees that entropy is independent of the size of the alphabet of the underlying random variable. Additivity captures the desire that the total uncertainty of two independent random variables equals the sum of their individual uncertainties. Finally, normalization ensures that a fair coin toss is assigned one unit of entropy.
%
Any quantity $\H(X)$ satisfying these axioms can be expressed as a convex combination of Rényi entropies~\cite{mu21_renyi,gour21_axiomatic}.\footnote{More precisely, a one-to-one relationship between entropy and relative entropy is established in~\cite{gour21_axiomatic} and \cite{mu21_renyi} gives a complete characterization of relative entropy. The statement is also a special case of our main theorem.} Namely,
\begin{align}\label{eq:EntropyBarycentre}
    \H(X) = \int_{[0,\infty]} \d\mu(\alpha) \, H_{\alpha}(X)
\end{align}
for some probability measure $\mu$ on the extended positive reals. Our goal here is to derive a similar complete characterization for conditional entropy.

Conditional entropies extend the notion of entropy to settings where additional information, known as side information, is available, quantifying uncertainty from the perspective of an observer with access to it. 
Consider thus two discrete random variables $X$ on $\mc{X}$ and $Y$ on $\mc{Y}$ governed by a joint pmf $P(x,y) = P(x|y)P(y)$ that we decomposed in terms of the conditional pmf $P(x|y)$ and marginal $P(y)$. We define the conditional surprisal given an event $Y = y$ was observed as $S_{y}(X) = \log \frac{1}{P(X|y)}$ for each $y \in \mc{Y}$. Noting that $H(X|Y\!=\!y) := \mb{E}[S_y]$ is a random variable once we vary $y$, the conditional Shannon entropy of $X$ given $Y$ is then expressed as a total expectation, 
\begin{align}
    H(X|Y) := \sum_{x, y} P(x,y) \log \frac{1}{P(x|y)} = \mb{E}\big[\mb{E}[S_Y] \big] \,,
\end{align}

Several definitions of conditional Rényi entropies have been proposed and found operational significance (see, e.g., \cite{teixeira2012conditional} and \cite{fehr14conditional} for reviews that cover some aspects of conditional entropy). We introduce some of the most widely used definitions here, and will later show that they are all special cases of a general class. Two prominent examples are $H_{\alpha}^H$, due to Hayashi and Skorić \emph{et al.}~\cite{hayashi2011exponential,vskoric2011sharp}, and $H_{\alpha}^A$, due to Arimoto~\cite{arimoto1977information}. They can be expressed as exponential expectations of the entropies $H_{\alpha}(X|Y\!=\!y) := \mb{E}_{1-\alpha}[S_y]$, i.e., as a total exponential expectation
\begin{align}
\label{Hayashi conditional entropy}
H_{\alpha}^H(X|Y) &:= \frac{1}{1-\alpha}\log \bigg(\sum_{x,y} P(y) P(x|y)^\alpha\bigg) = \mb{E}_{1-\alpha}\big[ \mb{E}_{1-\alpha}[S_Y] \big] \,,\\
\label{Arimoto conditional entropy}
H_{\alpha}^A(X|Y) &:= \frac{\alpha}{1-\alpha}\log{\bigg(\sum_y P(y)\bigg(\sum_x P(x|y)^\alpha\bigg)^\frac{1}{\alpha}\bigg)} = \mb{E}_{\frac{1-\alpha}{\alpha}} \big[ \mb{E}_{1-\alpha}[S_Y] \big] \,.
\end{align}
Convex combinations of these definitions have been considered in~\cite{hayashi2016uniform}.
A two-parameter family interpolating between them has been introduced by Hayashi and Tan~\cite{hayashi2016equivocations}:
$H_{\alpha,\beta}(X|Y) := \mb{E}_{\beta (1-\alpha)/\alpha}\big[ \mb{E}_{1-\alpha}[S_Y] \big]$.
As part of a study of quantum conditional entropy, its properties were studied in detail by some of the present authors~\cite%[Eq. (309)]
{rubboli2024quantum}.
This quantity possesses desirable properties if $(\alpha,\beta)$ are either in $(0,1)^{\times 2}$ or $(1,\infty)^{\times 2}$ and reduce to $H_{\alpha}^H$ and $H_{\alpha}^A$ for $\beta = \alpha$ and $\beta = 1$, respectively. The fact that these quantities have desirable properties for a large range of parameters shows that $\alpha$ and $\beta$ can be chosen largely independently.

In another line of work, Cachin~\cite{cachin1997entropy} and Renner and Wolf~\cite{renner2005simple} introduced the conditional entropies, which can see as limiting cases of the above family:
\begin{align}
    H^C_\alpha(X|Y) & = \sum_yP(y)H_\alpha(X|Y=y) = \mb{E}[\mb{E}_{1-\alpha}[S_Y]]\, \quad \textrm{and} \\
    H^R_\alpha(X|Y) &= \max_{y} H_\alpha(X|Y=y) = \mb{E}_{\infty}[\mb{E}_{1-\alpha}]]\,.    
\end{align} 
for $\alpha < 1$ and with $\max$ replaced by $\min$ for $\alpha > 1$.

Another two-parameter conditional entropy was proposed by Tan and Hayashi~\cite%[Eq.~(20)]
{tan2018conditional}. It can be written as 
$H_{\alpha,\beta}^T(X|Y) =  \mb{E}_{1-\alpha, (1-\alpha) (1-\beta)/\beta}[\mb{E}_{1-\alpha}[S_{Y}],\, \mb{E}_{1-\beta}[S_{Y}]]$, where we introduced a two-variable exponential expectation
\begin{align}
    \mb{E}_{t_1,t_2}[X_1, X_2] := \frac{1}{t_1+t_2} \log \mb{E}\big[ \exp( t_1 X_1 + t_2 X_2) \big] \,.
\end{align}
This indicates an even larger parameter space of operationally significant conditional entropies. 

%Another two-parameter conditional entropy was proposed by Hayashi and Watanabe~\cite%[Eq.~(15)]
%{hayashi2016uniform}, namely $\frac{1}{\alpha} H^T_{\alpha,\beta}(X|Y) + \frac{\alpha-1}{\alpha}H^A_\alpha(X|Y)$.

Indeed, in this work, we consider a very general form of conditional entropy (that we later prove to be the most general form) that is defined in terms of an exponential average of a continuous set of random variables. For any probability measure $\tau$ and $t \in \mathbb{R}$, we define
\begin{align}
    \mb{E}_{t, \tau}[ \alpha \mapsto X_\alpha] = \frac{1}{t} \log \mb{E}\Bigg[ \exp \Bigg( t \int_{[0,\infty]} X_\alpha\, \d \tau(\alpha) \Bigg) \Bigg] \, ,
\end{align}
which allows us to finally introduce our quantity of interest,
\begin{align}
 & H_{t, \tau}(X|Y) = \mb{E}_{t, \tau} \big[ \alpha \mapsto \mb{E}_{1-\alpha}[S_Y] \big] =\frac{1}{t} \log \sum_{y} P(y)\exp \bigg(t \int_{[0,\infty]}H_\alpha(X|Y\!=\!y)\, \d\tau(\alpha)\bigg).\label{eq:General}
\end{align}
All the entropies listed above arise as special cases of this general conditional entropy for suitable choices of the parameters (see Section~\ref{sec: literature} for more details). %In particular, the conditional Shannon entropy is recovered when the measure $\tau$ is concentrated at the point $\alpha=1$, and we consider the limit $t \rightarrow 0$. Moreover, the Arimoto and Hayashi R\'{e}nyi conditional entropies are recovered in the case of a point measure $\tau$ concentrated at $\alpha$ and $t= (1-\alpha)/\alpha$ and $t= 1-\alpha$, respectively.

% Consider two discrete random variables $X$ on $\mc{X}$ and $Y$ on $\mc{Y}$.
% Without loss of generality (see our first axiom below), we can fix $\mc{X} = \{1, 2, \ldots, m\}$ and $\mc{Y} = \{1, 2, \ldots, n\}$. In the following, we represent the bipartite probability distribution $P(x,y)$ on $\mc{X} \times \mc{Y}$ as a matrix $P$ with $(P)_{x,y} = P(x,y)$. We define the set of such matrices by $\mathcal{P}_{m,n}$, i.e., $m \times n$ matrices with entries in $[0,1]$ for which the sum of all entries is $1$. A generic conditional entropy is then a functional
% \begin{align}
%     \H: \bigcup_{m , n \in \mathbb{N}} \mathcal{P}_{m,n} \to \mathbb{R}_{0}^{+} . \label{eq:general-condent}
% \end{align}
% For a joint distribution \(P\) on \(X\) (target variable) and \(Y\) (side information), a generic conditional entropy is written in the standard form \(\H(X|Y)_P\).

\paragraph*{An axiomatic approach.}
We consider the following four axioms (see also~\cite{gour2018conditional,gour2024inevitability}), which are motivated by our desire to interpret conditional entropy as a measure of uncertainty from the perspective of an observer with side information. Below, we view any joint probabilities $P$ on $\mc{X}\times\mc Y$ as matrices with rows and columns labeled by $x$ and $y$, respectively, i.e.,
$(P)_{x,y} = P(x,y)$.
\begin{enumerate}
    \item \textbf{Invariance under relabeling.} We require that relabeling outcomes or embedding them into a larger alphabet will not change the conditional entropy. 

    \smallskip
    
    Formally, consider two injective functions $f: \mc{X} \to \mc{X}'$ and $g: \mc{Y} \to \mc{Y}'$ and $Q$ with $Q\big(f(x),g(y)\big) = P(x,y)$ for all $x \in \mc{X}, y \in \mc{Y}$ and zero elsewhere. We require that
    \begin{align}
        \H(X|Y)_P = \H(X'|Y')_Q \,.
    \end{align}

\label{it: relabeling and embedding}
    \item \textbf{Monotonicity under conditional mixing.} We expect that the entropy is non-decreasing when we apply a mixing (doubly stochastic) channel acting on $X$. For example, if $X$ is the value of the top card of a deck, shuffling (i.e., applying a more or less random permutation) should never decrease the uncertainty of $X$ from the perspective of an observer.
    %
    Moreover, we also want to account for conditional operations where the mixing channel depends on the side information. This accounts for the possibility that an observer is correlated with the randomness used in the mixing channel, e.g., the observer might know which shuffling technique is used, reducing the randomness introduced in the process compared to an observer without such information. Such conditional mixing channels can signal from $Y$ to $X$ but do not leak information about $X$ to the observer, and thus ought not reduce uncertainty.
    Finally, we require monotonicity under arbitrary channels applied to the side information $Y$ since local processing of the side information can never reduce our uncertainty of $X$.
    The closure under composition of such operations gives us conditionally mixing channels.

    \smallskip

   When representing a joint probability mass function as a matrix, doubly stochastic channels acting on the system $X$ can be described by right multiplication with a doubly stochastic matrix $S$. In addition, channels acting on the side information correspond to right multiplication by a row-stochastic matrix $D$. Finally, conditional operations can be represented by a family of row sub-stochastic matrices $D_i$ acting on the right, whose sum $\sum_i D_i$ is row-stochastic, together with a corresponding family of doubly stochastic matrices $S_i$ acting on the left (see~\cite{gour2018conditional} for more details).
    Explicitly, conditionally mixing channels from $m \times n$ matrices $P$ to $m \times n'$ matrices $Q$ act as
    \begin{align}
        P \mapsto Q = \sum_{i=1}^k S_i P D_i \,, \label{eq:intro/mixing}
    \end{align}
    where $S_i$ are doubly stochastic $m \times m$ matrices for all $i$ and $D_i$ are sub-stochastic $n \times n'$ matrices such that their sum $\sum_{i=1}^k D_i$ is row-stochastic.   We require that for any conditionally mixing channel, any $P$, and $Q$ as defined in~\eqref{eq:intro/mixing}, it holds that
    \begin{align}
        \H(X|Y)_P \leq \H(X|Y')_Q \,.
    \end{align}
\label{it: conditionally mixing channels}
    \item \textbf{Additivity.} We require that entropies for independent random variables are additive. This axiom expresses the principle that the uncertainty of two independent random variables is the sum of their individual uncertainties, and it admits a natural extension to scenarios with side information.

    \smallskip
    
    Formally, for any joint pmf $P$ and $Q$, we require that their Kronecker product, $P \otimes Q$, satisfies
    \begin{align}
        \H(XX'|YY')_{P\otimes Q} = \H(X|Y)_P + \H(X'|Y')_Q \,.
    \end{align}

    \item \textbf{Normalization}: We fix the entropy of a fair coin toss to be $1$ bit, i.e., we set $\H(X)_P = 1$ for $P = (1/2,1/2)^T$. 

\end{enumerate}

% \bigskip

% It is useful to combine the methods of comparing joint distributions presented in Conditions~\ref{it: relabeling and embedding}) and~\ref{it: conditionally mixing channels}) above. 
% For joint probability distributions $P$ and $Q$, we write $P \succeq Q$ if there exist $P'$ and $Q'$, obtained from $P$ and $Q$, respectively, by relabeling and embedding, such that $Q'$ can then be obtained from $P'$ through conditional mixing.
% In this case, we say that $P$ conditionally majorizes $Q$.

% A particular consequence of~\ref{it: relabeling and embedding}) and~\ref{it: conditionally mixing channels}) is that $\H(XY|Z)_P \geq \H(X|Z)_P$ for any joint distribution $P$ over three random variables, i.e., the property that entropy is non-increasing under coarse-graining. This can be verified since the channel $W_{Y|XZ}$ satisfying $P_{XYZ} = W_{Y|XZ} P_{XZ}$ can be implemented using such channels.

In this work, we prove that any conditional entropy that satisfies these four properties is given by a convex combination of conditional entropies $H_{t,\tau}$ defined in~\eqref{eq:General} and their pointwise limits (see Theorem~\ref{th:barycenter conditional entropies} for more details). 
While this result determines the most general form of a conditional entropy, some choices of the parameters \(t\) and \(\tau\) should be excluded if the resulting \(H_{t,\tau}\) does not satisfy all required axioms.
Specifically, we derive a set of sufficient conditions on the parameters $t$ and $\tau$ under which $H_{t,\tau}$ satisfies all axioms. 
These conditions are as follows:
\begin{enumerate}
    \item $t> 0$, $\;\text{supp}(\tau) \subseteq [0,1]$, and
    \begin{equation}\label{eq:integralcond}
    \int_{[0,+\infty]}\frac{\alpha}{1-\alpha}\,\mathrm{d}\tau(\alpha)\leq \frac{1}{t}
    \end{equation}
    %$\;\displaystyle\int_{[0,1)} \frac{\alpha}{1 - \alpha} \, \mathrm{d}\tau(\alpha) \leq 1/t$.
    \item $t \leq 0$ and
    \begin{itemize}
    \item[(a)] $\;\text{supp}(\tau) \subseteq [0,1]$ or
    \item[(b)] $\;\operatorname{supp}(\tau) \subseteq [0,1] \cup \{\alpha_*\}$ for some $\;\alpha_* \in (1, +\infty]$ so that also \eqref{eq:integralcond} is satisfied.
    %\item $t \leq 0$, $\;\operatorname{supp}(\tau) \subseteq [0,1] \cup \{\alpha_k\}$, $\;\alpha_k \in (1, +\infty]$, $\;\displaystyle\int_{[0,+\infty]} \frac{\alpha}{1 - \alpha} \, \mathrm{d}\tau(\alpha) \leq 1/t$.
    \end{itemize}
\end{enumerate}
In addition, we show that these conditions are also necessary whenever the measure \(\tau\) is discrete and supported on a finite number of points, or more generally when it satisfies a regularity assumption (see Section~\ref{sec: set of conditional entropies} for details).
Our analysis further leads us to conjecture that the same conditions remain necessary for a general measure $\tau$. 


\paragraph*{Applications.}

Majorization is a well-studied concept with applications across many fields, including probability theory and statistics~\cite{marshall1979inequalities}, mathematics~\cite{bhatia1997graduate}, and quantum information~\cite{nielsen1999conditions,horodecki2013fundamental}. Conditional majorization provides a natural extension of classical majorization to scenarios involving side information. Recalling the definitions introduced above, given two joint probability mass functions $P$ and $Q$, we say that $P$ conditionally majorizes $Q$ if there exists a conditionally mixing channel that maps $P$ to $Q$. When the underlying dimensions of the distributions do not coincide, suitable embeddings are implicitly understood (see Definition~\ref{def:cond-maj-cond-mix}).
Conditional majorization has found applications in areas such as conditional uncertainty relations in quantum information~\cite{gour2018conditional} and the study of games of chance~\cite{brandsen2022entropy}. 
In this work, we provide an operational interpretation of the newly defined conditional entropies $H_{t,\tau}$ as the quantities that determine the optimal rate $R$ for the multi-copy transformation $ P_{XY}^{\otimes n} \;\longrightarrow\; Q_{X'Y'}^{\otimes R n}$
under conditional majorization, for a sufficiently large number of copies $n$. Specifically, we show that
\begin{equation}
R(P_{XY}\rightarrow Q_{X'Y'}) = \inf\left\{ \frac{H_{t,\tau}(X'|Y')_Q}{H_{t,\tau}(X|Y)_P}\Bigg| (t,\tau) \in \D_{\rm{bulk}}\right\} \cup \left\{ \frac{H_{-\infty,\tau}(X'|Y')_Q}{H_{-\infty,\tau}(X|Y)_P}\Bigg| \tau \in \D_{-\infty}\right\}.
\end{equation}
Here, \(\mc D_{\rm bulk}\) denotes the set of parameters for which \(H_{t,\tau}\) is monotone under conditionally mixing channels, while \(\mc D_{-\infty}\) denotes the corresponding set of parameters for which the limiting entropy \(H_{-\infty,\tau}\), obtained as \(t \to -\infty\), satisfies the same monotonicity property.
This establishes a direct link between conditional entropies and asymptotic transformation rates.

\medskip

Our final result addresses transformations with a conditionally preserved probability distribution. We derive sufficient conditions for the existence of such a conditionally preserving map that, given an input distribution together with another probability distribution acting as a catalyst, produces an output distribution arbitrarily close to the desired target, while leaving the catalyst unchanged at the end of the transformation (see Theorem~\ref{thm:ConservedClassical}).
We further apply this result to derive a set of second laws of thermodynamics in the presence of side information, a setting that has been investigated in several works~\cite{sagawa2008second,morris2019assisted,rio2011thermodynamic,bera2017generalized,ji2025fundamental,narasimhachar2017resource}.
In the absence of side information, a complete set of necessary and sufficient conditions for transforming states that are diagonal in the energy eigenbasis under thermal operations—assisted by a catalyst that is returned unchanged at the end of the process—was established in~\cite{Brandao}.
 These conditions are commonly referred to as the second laws of quantum thermodynamics. In the absence of side information, these laws can be expressed in terms of certain free energies, which are connected to the Rényi relative entropies between a state and its Gibbs state. 
 In our work, we find that for states that are diagonal in the energy eigenbasis, certain quantities, denoted by $F_{t,\tau}$,  can be used to provide a set of sufficient conditions for transformations under conditioned thermal operations assisted by a catalyst. In particular, 
 \(F_{t,\tau}\) can be constructed from \(H_{t,\tau}\) in~\eqref{eq:General} by replacing the entropies in the exponential expectation with the R\'enyi relative entropies between the state and the Gibbs state. These quantities provide a set of sufficient conditions for state transformations under conditioned thermal operations assisted by a catalyst (see Corollary~\ref{cor:secondlaws} for more details).
This partially solves the open problem posed in \cite[Section V.B]{ji2025fundamental}.
This result shows that even in the classical case, the presence of side information significantly complicates the second laws. Nevertheless, for the specific case we consider, the problem admits a partial solution, whereas the fully quantum case remains an open challenge, even in the absence of side information.


\paragraph*{Methods.} 
To derive the general form of conditional entropy, we establish sufficient conditions under which, in the limit of a large number of copies, one joint distribution can be transformed into another joint distribution using conditionally mixing channels. This is because, by known arguments~\cite{mu21_renyi} (see also~\cite{haapasalo_inprep} for a general statement), this problem can be shown to be equivalent to the axiomatic characterization.
To derive these conditions, we apply the general results on preordered semirings developed in~\cite{fritz2023abstract,fritz2023abstractII}. 
In particular, we construct a specific semiring where the conditional entropies \( H_{t,\tau} \) arise as the relevant functions of the theory.   

\bigskip

The remainder of this manuscript is structured as follows: In Section~\ref{sec: formal definitions}, we introduce the formal definitions adopted in the manuscript. In Subsection~\ref{sec: set of entropies}, we derive the general form of entropies as convex combinations of R\'enyi entropies an important result often subsequently needed. In Section~\ref{sec: semiring of conditional entropies}, we define the semiring of conditional majorization and determine the general form of the extremal conditional entropies $H_{t,\tau}$. In Sections~\ref{sec: sufficient conditions measure} and~\ref{sec: necessary conditions measures}, we establish sufficient and necessary conditions on the parameters $t$ and $\tau$ for which the conditional entropies $H_{t,\tau}$ satisfy the required axioms. In Section~\ref{sec: large sample}, we derive large-sample and catalytic results that form the basis for our applications. In Section~\ref{sec: set of conditional entropies}, we prove the general form of conditional entropies, showing that they are convex combinations of the (extremal) entropies $H_{t,\tau}$. In Section~\ref{sec: Applications}, we discuss applications to rates and the second laws of quantum thermodynamics with side information.

\section{Definitions and preliminary results}
\label{sec: formal definitions}
This section provides the formal definitions that serve as the foundation for the remainder of the manuscript. We also present and prove an exhaustive characterization of unconditional entropies serving as a stepping stone towards the characterization of conditional entropies.

\subsection{Definitions relevant for entropies}

We denote the (always finite) set of possible values of a random variable $X$ by $\mc X$. The probability mass function of $X$ is denoted by $P_X$ and often if there is no risk of confusion, we may drop the subscript. To avoid confusion, we also use other names for the probability mass functions like $Q_X$, $R_X$, and the like. We use similar notational conventions for other random variables like $Y$, $Z$, and so forth. We also usually call probability mass functions as probability distributions for brevity.

\begin{definition}[Embedding and relabeling of probability distributions]
Given two probability distributions \(P_{X}\) and \(\widetilde{P}_{\widetilde{X}}\), we say that \(\widetilde{P}_{\widetilde{X}}\) is obtained from \(P_{X}\) through embedding into a larger space and relabeling of outcomes if there exists an injective function \( f : \mathcal{X} \to \widetilde{\mathcal{X}} \) such that  $\widetilde{P}_{\widetilde{X}}(f(x)) = P_X(x)$
for all $x \in \mathcal{X}$
and zero elsewhere.
\end{definition}

\begin{definition}[Doubly stochastic channel]
A channel is called \emph{doubly stochastic} (or \emph{mixing}) if, when acting on a probability vector, it is represented by a matrix with nonnegative entries whose rows and columns each sum to one.
\end{definition}

\begin{definition}[Majorization]
Given two probability distributions $P_X$ and $Q_{X'}$, we say that $P_X$ majorizes $Q_{X'}$, written as
\begin{equation}
    P_X \succeq Q_{X'}
\end{equation}
if there exist probability distributions \(\widetilde{P}_{\widetilde{X}}\) and \(\widetilde{Q}_{\widetilde{X}'}\), obtained from \(P_X\) and \(Q_{X'}\) via embedding and relabeling, and a doubly stochastic map \(D\) such that $D \widetilde{P}_{\widetilde{X}} = \widetilde{Q}_{\widetilde{X}'}$.
\end{definition}

\begin{remark}
    It is known that it suffices to consider embeddings into a space whose dimension equals the maximum of the supports of the input and output probability distributions. Accordingly, in the following, we assume that the probability distributions $P_X$ and $Q_X$ are defined on a common underlying space. 
\end{remark}
It is well known that the existence of a doubly stochastic map is equivalent to the following condition on the entries of the probability distributions, often referred to as the initial sum condition~\cite{birkhoff1946tres} (see also~\cite{marshall1979inequalities}).
\begin{lemma}[Initial sum condition]
\label{lem: Initial sum condition}
Let \( P_X \) and \( Q_X \) be two $d$-dimensional vectors with nonnegative entries. Moreover, let \( P_X^{\downarrow} \) and \( Q_X^{\downarrow} \) denote the vectors obtained by rearranging the components of \( P_X \) and \( Q_X \), respectively, in non-increasing order. Then \( P_X \succeq Q_X \) if and only if
\begin{equation}
    \sum_{i=1}^{k} P_X^{\downarrow}(i) \geq \sum_{i=1}^{k} Q_X^{\downarrow}(i) \qquad \text{for all } k = 1, \ldots, d-1,
\end{equation}
and equality holds for \( k = d \), that is, $\sum_{i=1}^{d} P_X^{\downarrow}(i) = \sum_{i=1}^{d} Q_X^{\downarrow}(i)$.
\end{lemma}



\begin{definition}[Entropy]
\label{def: Entropy}
Let us denote, for any finite random variable $X$, the set of probability distributions on $X$ by $\mc P_X$. A function
\begin{equation}
\H : \bigcup_{X} \mc P_{X} \rightarrow \mathbb{R}_+
\label{eq:ent-gen-def-1st}
\end{equation}
is called entropy if it satisfies the following postulates for all systems $X$, $X'$, and probability vectors $P_{X}$ and $Q_{X'}$:
\begin{enumerate}
    \item Invariance under relabeling and embedding: If $P_X$ can be relabeled and embedded into $Q_{X'}$, then $\H(X)_P=\H(X')_Q$. 
     \label{eq:E-2nd-postulate-iso}
    \item Monotonicity under mixing channels: if $Q_{X'}$ is obtained from $P_{X}$ using a mixing channel, then $ \H(X)_{P} \leq \H(X')_{Q}$.
    \label{eq:E-1st-postulate-mono} 
    \item Additivity: $  \H(XX')_{P \otimes Q} = \H(X)_{P} + \H(X')_{Q}$
     \label{eq:E-2nd-postulate-add}
    \item Normalization: $ \H(X)_{(1/2,1/2)} = 1$
    \label{eq:E-4-postulate-normalization}
\end{enumerate}
\end{definition}
Note that the first two properties imply that if \(P_{X} \succeq Q_{X'}\) (i.e., \(P_X\) majorizes \(Q_{X'}\)), then \(\H(X)_P \leq \H(X')_Q\). In other words, any entropy is monotone under majorization. In the following, we will 
also often use the notation $\H(X)_P=\H(P_X)$. Taking into account condition~\ref{eq:E-2nd-postulate-iso} above, we may effectively consider an entropy $\H$ as a map on the union of $\mc P_n$ over $n\in\mathbb{N}$ where $\mc P_n$ is the set of probability distributions over $\{1,\ldots,n\}$. 

\subsection{Definitions relevant for conditional entropies}

In the sequel, we shall discuss joint probability distributions $P_{XY}:\X\times\Y\to[0,1]$. We may often again drop the subscripts. 
The marginals of \(P_{XY}\) are denoted by \(P_X\) and \(P_Y\), and are defined as
\begin{equation}
P_X(x) = \sum_{y \in \mathcal{Y}} P_{XY}(x,y), \quad \forall x \in \mathcal{X}, \qquad
P_Y(y) = \sum_{x \in \mathcal{X}} P_{XY}(x,y), \quad \forall y \in \mathcal{Y}.
\end{equation}
We will always identify a joint probability distribution $P_{XY}$ with the matrix
\begin{equation}
P_{XY}=\big(P_{XY}(x,y)\big)_{(x,y)\in\X\times\Y}.
\end{equation}
This makes many of our future notation and definitions, including the next one, simpler.

\begin{definition}[Conditionally mixing channel]
A \emph{conditionally mixing channel} is a channel of the form
\begin{equation}
P_{XY} \longmapsto Q_{XY'} = \sum_{i=1}^k S^{(i)} P_{XY} D^{(i)}, \footnote{We chose the letters $S$ and $D$ to reflect the Latin terms for left and right, namely sinister and dexter, respectively.}
\end{equation}
where the following conditions hold:
\begin{enumerate}
    \item The matrices \( S^{(i)} \) are doubly stochastic for all $i=1,\dots,k$.
    
    \item The matrices \( D^{(i)} \) are sub-stochastic for all \( i = 1, \ldots, k \).
    \item The sum $\sum_{i=1}^k D^{(i)}$ is row-stochastic.
\end{enumerate}
\end{definition}
For clarity of notation and to facilitate the understanding of matrix dimensions in the arguments that follow, we note that if \( Q_{XY'} \) is a \( d \times n' \) matrix and \( P_{XY} \) is a \( d \times n \) matrix, then the matrices \( S^{(i)} \) have dimensions \( d \times d \), while the selection matrices \( D^{(i)} \) have dimensions \( n \times n' \).


\begin{definition}[Embedding and relabeling of joint probability distributions]
Given two joint probability distributions \(P_{XY}\) and \(\widetilde{P}_{\widetilde{X}\widetilde{Y}}\), we say that \(\widetilde{P}_{\widetilde{X}\widetilde{Y}}\) is obtained from \(P_{XY}\) through embedding into a larger space and relabeling of outcomes if there exist two injective functions
$f:\X\to\widetilde{\X}$ and $g:\Y\to\widetilde{\Y}$ such that $\widetilde{P}_{\widetilde{X}\widetilde{Y}}(f(x),g(y))=P_{XY}\big(x,y\big)$ for all $x\in\X$ and $y\in\Y$ and zero elsewhere.
\end{definition}



\begin{definition}[Conditional majorization]
\label{def:cond-maj-cond-mix}
Given two probability distributions $P=P_{XY}$ and $Q=Q_{X'Y'}$, we say that $P_{XY}$ conditionally majorizes $Q _{X'Y'}$, and write 
\begin{equation}
    P_{XY}\succeq Q_{X'Y'}
        \label{eq:cond-maj-def-1}
\end{equation}
if there exist probability distributions $\widetilde{P}_{\widetilde{X}\widetilde{Y}}$ and $\widetilde{Q}_{\widetilde{X}'\widetilde{Y}'}$, obtained from \(P_{XY}\) and \(Q_{X'Y'}\) via embedding and relabeling, and a conditionally mixing channel that maps $\widetilde{P}_{\widetilde{X}\widetilde{Y}}$ into $\widetilde{Q}_{\widetilde{X}'\widetilde{Y}'}$.
This means that, up to relabeling and embedding, $P_{XY}$ can be transformed into $Q_{X'Y'}$ using a conditional mixing channel.
\end{definition}


\begin{remark}
\label{remark: no isometries}
We observe that when considering transformations between two probability distributions \(P_{XY}\) and \(Q_{XY'}\) that have the same number of rows, no relabeling or embedding beyond the conditionally mixing channel is required. In other words, if \(P_{XY} \succeq Q_{XY'}\) according to Definition~\ref{def:cond-maj-cond-mix}, then there exists a conditionally mixing channel that maps \(P_{XY}\) to \(Q_{XY'}\). 

To see this, consider embeddings \(P_{\widetilde{X}Y}\) and \(Q_{\widetilde{X}Y'}\) obtained by padding \(P_{XY}\) and \(Q_{XY'}\) with additional rows of zeros. Since the sum of positive elements can vanish only if all elements are zero, it follows that the rows with index \(j > |\mathcal{X}|\) of $S^{(i)} P_{\widetilde{X}Y} R^{(i)}$
must be zero. Moreover, \(P_{\widetilde{X}Y} R^{(i)}\) acts solely on the columns, leaving the additional rows \(j > |\mathcal{X}|\) unchanged at zero. By Birkhoff's theorem, each \(S^{(i)}\) can be expressed as a convex combination of permutations. Any permutation that acts only on the zero elements can be replaced by the identity, while permutations transferring mass from \(j \leq |\mathcal{X}|\) to \(j > |\mathcal{X}|\) are forbidden since the latter rows must remain zero. Therefore, it suffices to consider only permutations acting within the subspace spanned by rows \(j \leq |\mathcal{X}|\).
\end{remark}

\begin{definition}[Conditional entropy]
\label{def: conditional entropy}
For any finite random variables $X$ and $Y$, we denote the set of probability distributions $P_{XY}$ by $\mc P_{XY}$. A function
\begin{equation}
\H : \bigcup_{X, Y} \mc P_{XY} \rightarrow \mathbb{R}_+
\label{eq:cond-ent-gen-def-1st}
\end{equation}
is called a conditional entropy if it satisfies the following postulates for all systems $X$, $Y$, $X'$, $Y'$, and joint probability distributions $P_{XY}$ and $Q_{X'Y'}$:
\begin{enumerate}
    \item Invariance under relabeling and embedding: if $Q_{X'Y'}$ is obtained from $P_{XY}$ through relabeling and embedding, then $\H(X|Y)_P=\H(X'|Y')_Q$. 
 \label{eq:QCE-2nd-postulate-iso}
    \item Monotonicity under conditionally mixing channels: if $Q_{X'Y'}$ is obtained from $P_{XY}$ using a conditionally mixing channel, then $\H(X|Y)_P\leq\H(X'|Y')_Q.$ 
\label{eq:QCE-1st-postulate-mono} 
    \item Additivity: $\H(XX'|YY')_{P \otimes Q} = \H(X|Y)_{P} + \H(X'|Y')_{Q}$.
  \label{eq:QCE-3rd-postulate-add}
    \item Normalization: $\H(X|Y)_{(1/2,1/2) \otimes Q_Y} = 1$
     \label{eq:QCE-4-postulate-normalization}
\end{enumerate}
\end{definition}
Note that the first two properties imply that if \(P_{XY} \succeq Q_{X'Y'}\) (i.e., \(P_{XY}\) conditionally majorizes \(Q_{X'Y'}\)), then \(\H(X|Y)_P \leq \H(X'|Y')_Q\). In other words, any conditional entropy is monotone under conditional majorization. In the following, we will 
also often use the notation $\H(X|Y)_P=\H(P_{XY})$.

\subsection{The set of entropies}
\label{sec: set of entropies}
In this section, we derive the general form of an entropy, showing that it can be written as a convex combination of Rényi entropies. We also refer to this representation as the barycentric decomposition of entropy.
This result builds on~\cite[Proposition 3.7]{Jensen_Kjaerulf_2019}, which establishes sufficient conditions for transforming multiple copies of one probability distribution into another under doubly stochastic maps. We refer to this as a large-sample result for entropy. We then 
follow the proof strategy of~\cite[Theorem 2]{mu21_renyi}, which uses a compact set of conditions characterizing large-sample transformations involving a set of monotones to establish barycentric decompositions based on those same monotones. This approach is quite general, as emphasized in~\cite{haapasalo_inprep}. In the present setting, the monotones coincide with the Rényi entropies.

\begin{lemma}[{\cite[Proposition 3.7]{Jensen_Kjaerulf_2019}}]\label{lemma:Jensen}
Let $P_X,Q_X$ be two probability distributions and assume that 
\begin{equation}
    H_\alpha(P_X) < H_\alpha(Q_X)  \quad \forall \alpha\in[0,+\infty] \,.
\end{equation}
Then, for sufficiently large $n \in \mathbb{N}$ there exists a doubly stochastic matrix $S_n$ such that
$S_nP_X^{\otimes n} = Q_X^{\otimes n}$.
\end{lemma}
The latter result implies a barycentric decomposition for entropies as we state next.
In particular, this yields a barycentric decomposition with respect to a positive measure, which must be normalized if the additional normalization axiom is imposed (see Proposition~\ref{proposition:barycenter entropies} below).
 Notably, this result will subsequently serve as a key ingredient in deriving the general form of a conditional entropy.  Note that, as we discuss in Remark~\ref{remark: negative alpha dilemma}, only the Rényi entropies with 
$\alpha \geq 0$ appear in the decomposition. Indeed, the entropies with $\alpha<0$ fail to satisfy the invariance under relabeling and embedding axiom.



\begin{proposition}[General form of entropies]\label{proposition:barycenter entropies}
Let $\H$ be an entropy according to Definition~\ref{def: Entropy}. Then, there exists a probability measure $\mu:\mc B\big([0,+\infty]\big)\to\mb R_+$ such that for all probability distributions $P_X$
\begin{equation}\label{eq:barycenter}
\H(P_X)=\int_{[0,+\infty]}H_\alpha(P_X)\d\mu(\alpha) \,.
\end{equation}
Furthermore, if the normalization axiom~\ref{eq:E-4-postulate-normalization} is not assumed, the conclusion remains valid upon replacing the probability measure with a finite positive measure.
\end{proposition}
\begin{proof}
    The result follows from Lemma \ref{lemma:Jensen}; the proof is essentially the same as that of \cite[Theorem 2]{mu21_renyi} and this result can be seen as a special case of Theorem 7 of \cite{haapasalo_inprep}. We give an outline of this proof for completeness. This same proof also works for Theorem~\ref{th:barycenter conditional entropies} later in this work with slight adaptations.

    The first step is to show that $H_\alpha(P_X)\leq H_\alpha(Q_X)$ for all $\alpha\in[0,\infty]$ implies $\mb H(P_X)\leq\mb H(Q_X)$. First, let us assume that $H_\alpha(P_X)< H_\alpha(Q_X)$ for all $\alpha\in[0,\infty]$. According to Lemma \ref{lemma:Jensen}, this means that, for any $n\in\mb N$ sufficiently large, there is a doubly stochastic matrix $S_n$ such that $S_nP_X^{\otimes n}=Q_X^{\otimes n}$. Thus, since $\mb H$ is monotone and additive,
    \begin{equation}
    n\mb H(P_X)=\mb H(P_X^{\otimes n})\leq\mb H(Q_X^{\otimes n})=n\mb H(Q_X),
    \end{equation}
    i.e., $\mb H(P_X)\leq\mb H(Q_X)$. Let us then assume that $H_\alpha(P_X)\leq H_\alpha(Q_X)$ for all $\alpha\in[0,\infty]$. Denoting by $R$ the two-outcome uniform distribution $(1/2,1/2)$, it follows that
    \begin{equation}
    H_\alpha(P_X^{\otimes m})=mH_\alpha(P_X)\leq mH_\alpha(Q_X)<mH_\alpha(Q_X)+1=H_\alpha(Q_X^{\otimes m}\otimes R\big)
    \end{equation}
    for all $m\in\mb N$ and $\alpha\in[0,\infty]$. Thus, according our initial step, we also have
    \begin{equation}
    m\mb H(P_X)=\mb H(P_X^{\otimes m})\leq\mb H(Q_X^{\otimes m}\otimes R)=m\mb H(Q_X)+1
    \end{equation}
    for all $m\in\mb N$. This means $\mb H(P_X)\leq\mb H(Q_X)+1/m$ for all $m\in\mb N$, so that $\mb H(P_X)\leq\mb H(Q_X)$, and the initial result is proven.

    For the details of the rest of the proof, we refer to the proof of Theorem 7 of \cite{haapasalo_inprep} and that of Theorem 2 of \cite{mu21_renyi}. The strategy is to define a positive linear functional on the set $C\big([0,\infty]\big)$ of continuous real functions on the compact set $[0,\infty]$ using the initial step proven above and then use the representation result for such functionals.
    
    Consider the evaluation functions ${\rm ev}_P\in C\big([0,\infty]\big)$, ${\rm ev}_P(\alpha)=H_\alpha(P)$ at any finite probability distribution $P$ and all $\alpha\in[0,\infty]$. By `evaluation' we essentially mean that we actually identify $[0,\infty]$ with the set of R\'{e}nyi entropies $H_\alpha$, $\alpha\in[0,\infty]$. This kind of identification has to be done in the general case of Theorem 7 of \cite{haapasalo_inprep} where the situation is such that large-sample ordering conditions are given by a compact set of functions which is Hausdorff. This is also the case of Theorem \ref{th:barycenter conditional entropies} later on. It can be shown very easily, using our initial step, that we may set up a well-defined positive functional $F$ on the positive cone generated by the evaluation functions in $C\big([0,\infty]\big)$ through $F({\rm ev}_P)=\mb H(P)$. For the well-definedness, note that if $P$ and $Q$ are such that ${\rm ev}_P={\rm ev}_Q$, i.e., $H_\alpha(P)=H_\alpha(Q)$ for all $\alpha\in[0,\infty]$, then, according to our initial result, also $\mb H(P)=\mb H(Q)$. The positivity of $F$ follows similarly. With some extra work, we may show that $F$ extends into a positive functional $G$ on the vector space $V$ spanned by the evaluation functions.
    
    Suppose that $f\in C\big([0,\infty]\big)$ and fix any $n\in\mb N$ such that $n\geq\|f\|_\infty$. Denoting by $1$ the constant function 1 on $[0,\infty]$ and by $R$ the two-outcome uniform distribution, we have
    \begin{equation}
    f\leq \|f\|_\infty\,1\leq n\,1=n\,{\rm ev}_R={\rm ev}_{R^{\otimes n}},
    \end{equation}
    so that all functions in $C\big([0,\infty]\big)$ are upper bounded by functions from $V$. Now a theorem due to Kantorovich \cite{Kantorovich} related to the Hahn-Banach theorem implies that $G$ extends into a positive linear functional $I:C\big([0,\infty]\big)\to\R$. Due to the Riesz-Markov-Kakutani theorem, there is a finite positive measure $\mu:\mc B\big([0,\infty]\big)\to\R_+$ (the Borel $\sigma$-algebra associated with the one-point compactification of the natural topology of $[0,\infty)$) such that
    \begin{equation}
    I(f)=\int_{[0,\infty]}f\,\mathrm{d}\mu,\qquad\forall f\in C\big([0,\infty]\big).
    \end{equation}
    This means
    \begin{equation}
    \mb H(P)=F({\rm ev}_P)=I({\rm ev}_P)=\int_{[0,\infty]}{\rm ev}_P(\alpha)\,\mathrm{d}\mu(\alpha)=\int_{[0,\infty]}H_\alpha(P)\,\mathrm{d}\mu(\alpha).
    \end{equation}
    Denoting again by $R$ the two-outcome uniform distribution, we have
    \begin{equation}
    1=\mb H(R)=\int_{[0,\infty]}\underbrace{H_\alpha(R)}_{=1}\,\mathrm{d}\mu(\alpha)=\mu\big([0,\infty]\big),
    \end{equation}
    i.e., $\mu$ is a probability measure.
\end{proof}

\begin{remark}
\label{remark: negative alpha dilemma}
  One might wonder why Rényi entropies for $\alpha < 0$ are excluded from Proposition~\ref{proposition:barycenter entropies}. The reason lies in the requirement of invariance under isometries stated in Definition~\ref{def: Entropy}. Intuitively, when a full-rank probability distribution is embedded into a larger space—effectively appending zeros to the distribution—the Rényi entropy for $\alpha < 0$ diverges, since raising zero to a negative power yields infinity. As a result, the entropy changes from a finite to an infinite value under such an embedding, violating the isometric invariance axiom~\ref{eq:E-2nd-postulate-iso}.

To formalize this argument, suppose there exists an entropy function that coincides with the Rényi entropy of order $\alpha < 0$ on full-rank probability distributions, is invariant under embeddings, and satisfies monotonicity under doubly stochastic maps—namely, $H_\alpha(X')_P \geq H_\alpha(X')_Q$ whenever $Q_{X'} = S P_{X'}$ for some doubly stochastic matrix $S$. For $\alpha < 0$, such monotonicity is expected due to the Schur concavity of the function $\left( \sum_i x_i^\alpha \right)^{1/\alpha}$. In particular, for negative orders $\alpha<0$, the entropies $H_\alpha$ are expected to be monotone non-increasing, rather than non-decreasing, under doubly stochastic maps.
 Let us define
    \begin{align}
    P_X=\begin{pmatrix}
        \frac{1}{2} \\
        \frac{1}{2}
    \end{pmatrix} \,, \quad
    P_{X'}=\begin{pmatrix}
        \frac{1}{2} \\
        \frac{1}{2}  \\
        0
    \end{pmatrix} 
\,,\quad S=\begin{pmatrix}
            1-\varepsilon & 0 & \varepsilon \\
            0& 1 & 0\\
            \varepsilon & 0 & 1-\varepsilon \\
        \end{pmatrix} \,, \quad 
       Q_{X'} = SP_{X'} = \frac{1}{2}
       \begin{pmatrix}
           1-\varepsilon\\
           1\\
           \varepsilon
       \end{pmatrix} \,.
    \end{align}
Here, $P_{X'}$ represents $P_X$ embedded into a larger space, and $Q_{X'}$ is the output of the doubly stochastic map $S$ applied to $P_X'$. Then, invariance under isometries and monotonicity under a doubly stochastic map would imply the chain of inequalities
    \begin{align}
        1= H_\alpha(X)_P = H_\alpha(X')_P \geq H_\alpha(X')_Q = \frac{1}{1-\alpha}\log{\left(\frac{1}{2^\alpha}\left((1-\varepsilon)^{\alpha}+1+\varepsilon^\alpha\right)\right)} \xrightarrow[\varepsilon \to 0]{} +\infty.
    \end{align}
However, since $1 \neq  +\infty$, this leads to a contradiction. Therefore, no entropy function can simultaneously coincide with the Rényi entropy of order $\alpha < 0$ on full-rank probability distributions, be invariant under embeddings, and satisfy monotonicity under doubly stochastic maps.
\end{remark}

\section{The semiring of conditional majorization}
\label{sec: semiring of conditional entropies}
In this section, we introduce the main technical tool used to establish sufficient conditions for multiple-copy transformation (large sample result) that lead to the general form of conditional entropies (barycentric decomposition). Our approach relies on the real-algebraic theory of preordered semirings developed in \cite{fritz2023abstract,fritz2023abstractII}. This framework is particularly powerful because it provides a large sample theorem (Theorem \ref{thm:Fritz2022}) whenever one can endow the underlying problem with a suitable semiring structure.
The large sample comparison theorem provides a characterization of the sufficient conditions under which multiple copies of one element of the semiring can be transformed into multiple copies of another. 
In our setting, the semiring is essentially the set of joint probability distributions $P_{XY}$, and the large-sample theorem yields necessary and sufficient conditions under which multiple copies of a joint distribution $P_{XY}$ can be transformed into multiple copies of another joint distribution $Q_{XY}$ under conditionally mixing operations in terms of the so-called monotone homomorphisms (see later for the definition) of the theory. In the case of conditional entropies, the monotone homomorphisms and derivations are exactly the extremal conditional entropies (up to a constant and a logarithm). These are the conditional entropies in terms of which any general conditional entropy can be expressed as a convex combination. From here, a standard argument \cite{mu21_renyi,haapasalo_inprep} allows us to deduce from the large sample theorem that any conditional entropy can be expressed as a convex combination of these extremal conditional entropies.

We begin by introducing the fundamental concepts underlying the theory of semirings and the general results that will be employed in our manuscript in Section~\ref{sec: background semiring}. We then construct the semiring associated with conditional entropies and derive the general form of monotone homomorphisms and derivations for this particular setting in Sections~\ref{sec: all monotone homorphisms}, and~\ref{sec: derivations}.

\subsection{Background on preordered semirings and a \textit{Vergleichsstellensatz}}
\label{sec: background semiring}
We define a preordered semiring as a tuple $S = (\mathcal{S}, +, \cdot, 0, 1, \rleq)$, where $\mathcal{S}$ is a set equipped with binary operations of addition $+$ and multiplication $\cdot$, a zero element $0 \in \mathcal{S}$, a multiplicative unit $1 \in \mathcal{S}$, and a preorder relation $\rleq$  (a reflexive and transitive binary relation) defined on $\mathcal{S}$  satisfying
\begin{equation}
    x \rleq y \ \Rightarrow \ 
\begin{cases}
x + a \ \rleq\ y + a, \\[0.3em]
x a \ \rleq\ y a,
\end{cases}
\end{equation}
for all $a \in \mathcal{S}$.
Moreover, 
$(\mathcal{S}, +, 0)$ and $(\mathcal{S}, \cdot, 1)$ are commutative semigroups, and the multiplication distributes over the addition.
For compactness, we omit the multiplication dot when writing products in the semiring.
We write $x \sim y$ for the equivalence relation induced by $\rleq$, meaning that $x \sim y$ if and only if there exist $z_1, \ldots, z_n \in \mathcal{S}$ such that
\begin{equation}
    x \rleq z_1 \rgeq z_2 \rleq \cdots \rgeq z_n \rleq y.
\end{equation}
A preordered semiring $S$ has polynomial growth if it admits a power universal element $u \in \mathcal{S}$, that is,
\begin{equation}
    x \rleq y \quad \Rightarrow \quad \exists\, k \in \mathbb{N}:\ y \rleq x u^k.
\end{equation}
We say that a preordered semiring $S$ is a preordered semidomain if
\begin{equation}
    \begin{aligned}
&xy = 0 \ \Rightarrow\ x = 0 \ \text{or} \ y = 0, \\
&0 \rleq x \rleq 0 \ \Rightarrow\ x = 0.
\end{aligned}
\end{equation}
Finally, $S$ is zerosumfree if $x + y = 0$ implies $x = 0 = y$.

Given preordered semirings $S$ and $T$, we say that a map $\Phi:\cS\to T$ is a {\it monotone homomorphism} if it satisfies the properties
\begin{enumerate}
\item Additivity: $\Phi(x+y)=\Phi(x)+\Phi(y)$ for all $x,y\in \cS$,
\item Multiplicativity: $\Phi(xy)=\Phi(x)\Phi(y)$ for all $x,y\in \cS$,
\item Monotonicity: $x\rleq y$ $\Rightarrow$ $\Phi(x)\rleq\Phi(y)$, and
\item $\Phi(0)=0$ and $\Phi(1)=1$.
\end{enumerate}
We call a monotone homomorphism \textit{degenerate} if $ x \rleq y \quad \Rightarrow \quad \Phi(x) = \Phi(y)$. Otherwise, it is called \textit{nondegenerate}.
We need to consider monotone homomorphisms with values in certain special semirings, namely:
\begin{enumerate}
    \item $\mathbb{R}_+$: the half-line $[0,+\infty)$ with the usual addition, multiplication, and total order.
    \item $\mathbb{R}_+^{\mathrm{op}}$: the same set with reversed order. 
    \item $\mathbb{T}\mathbb{R}_+$: the half-line $[0,+\infty)$ with the usual multiplication, order, and the tropical sum $x + y := \max\{x,y\}$.
    \item $\mathbb{T}\mathbb{R}_+^{\mathrm{op}}$: the same as for $\mathbb{T}\mathbb{R}_+$, but with reversed order.
\end{enumerate}
The term \emph{temperate reals} refers to the pair $(\mathbb{R}_+, \mathbb{R}_+^{\mathrm{op}})$, while \emph{tropical reals} denotes the pair $(\mathbb{T}\mathbb{R}_+, \mathbb{T}\mathbb{R}_+^{\mathrm{op}})$.
Given a monotone homomorphism $\Phi: \cS \to \mathbb{R}_+$, an additive map $\Delta: \cS \to \mathbb{R}$ is called a derivation at $\Phi$ (or a $\Phi$-derivation) if it satisfies the Leibniz rule
\begin{equation}
    \Delta(xy) = \Delta(x)\,\Phi(y) + \Phi(x)\,\Delta(y)
\end{equation}
for all $x, y \in S$.
In this work, we are interested in $\Phi$-derivations at degenerate homomorphisms that are also monotone, i.e.,
\begin{equation}
    x \rleq y \quad \Rightarrow \quad \Delta(x) \leq \Delta(y).
\end{equation}
The following theorem provides both a large-sample and a catalytic result. It belongs to a series of results collectively known as the ``Vergleichsstellens\"{a}tze'' \cite{fritz2023abstractII}.
\begin{theorem}[Based on Theorem 8.6 in \cite{fritz2023abstractII}]\label{thm:Fritz2022}
Let $S$ be a zerosumfree preordered semidomain with a power universal element $u$. Assume that for some $d\in\mb N$ there is a surjective homomorphism $\|\cdot\|:\cS\to\R_{>0}^d\cup\{(0,\ldots,0)\}$ with trivial kernel and such that
\begin{equation}\label{eq:surjectivehomomorphismproperties}
a\rgeq b\ \Rightarrow\ \|a\|=\|b\|\quad {\rm and} \quad \|a\|=\|b\|\ \Rightarrow\ a\sim b.
\end{equation}
Denote the component homomorphisms of $\|\cdot\|$ by $\|\cdot\|_{(j)}$, $j=1,\ldots,d$. Let $x,y\in S\setminus\{0\}$ with $\|x\|=\|y\|$. If 
\begin{itemize}
\item[(i)] for every $\mb K\in\{\R_+,\R_+^{\rm op},\T\R_+,\T\R_+^{\rm op}\}$ and every nondegenerate monotone homomorphism $\Phi : \cS \to \mb K$ with trivial kernel, we have $\Phi(x) > \Phi(y)$, and
\item[(ii)] $\Delta(x) > \Delta(y)$ for every monotone $\|\cdot\|_{(j)}$-derivation $\Delta : \cS \to \R$ with $\Delta(u) = 1$ for all component indices $j = 1,\ldots,d$,
\end{itemize}
then
\begin{enumerate}[label=(\alph*),ref=2(\alph*)]
 \item there exists a nonzero $c\in \cS$ such that $cx\rgeq cy$, and 
\label{it: catalytic}
\item if additionally $x$ is power universal, then $x^n \rgeq y^n$ for all sufficiently large $n\in\mb N$. 
\label{it: asymptotic}
\end{enumerate}
Conversely, if either of these properties holds for at least an integer $n$ or a catalyst $c$, then the inequalities in items (i) and (ii) above hold non-strictly. 
\end{theorem}
Large-sample ordering as described in item~\ref{it: asymptotic} of Theorem~\ref{thm:Fritz2022} implies catalytic ordering as in item~\ref{it: catalytic} of Theorem~\ref{thm:Fritz2022}, where the catalyst can be chosen as
\begin{equation}
    c = \sum_{\ell=0}^{n-1} x^\ell y^{n-1-\ell}
\end{equation}
for sufficiently large $n \in \mathbb{N}$. This implication was originally established in \cite{duan2005multiple} for the case where $x, y$ are probability vectors, but the argument extends naturally to the more general framework considered here.
The converse, however, generally fails to hold. In particular,~\cite[Theorem 3]{feng2006relation} presents a method for constructing explicit counterexamples within a specific context.

\subsection{Definition of the conditional majorization semiring}
In this section, we define the semiring of conditional majorization. In particular, we need to specify the elements of the tuple $S=(\mathcal{S}, +, \cdot, 0, 1, \rgeq)$.

\begin{enumerate}

\item \textbf{The set of elements ``$\cS$"}. The elements of the semiring are, essentially, all (not normalized) joint probability distributions $P_{XY}$.
However, since distributions that differ only by embedding into a larger outcome space and/or permutation of outcomes are considered equivalent, the set $\cS$ is formally defined as the set of equivalence classes of joint probability distributions under this similarity relation. 
We make this precise in the following.

We can represent (unnormalized) joint distributions as matrices, with rows corresponding to the outcomes of the variable $X$ and columns to those of $Y$. For example, if $X$ has $d$ outcomes and $Y$ has $n$ outcomes, then $P_{XY}$ can be represented as an $d \times n$ matrix
\begin{equation}
P_{XY} = \begin{pmatrix}
P(x_1,y_1) & \hdots & P(x_1,y_n) \\
\vdots &  & \vdots \\
P(x_d,y_1) & \hdots & P(x_d,y_n) 
\end{pmatrix} \,.
\end{equation}
We also denote the set of outcomes $\X=\{x_1,...,x_d\}$ and $\Y=\{y_1,...,y_n\}$.
We denote with $\mathcal{V}_{d \times n}$ the set of all $d \times n$ matrices with non-negative entries. Furthermore, define
\begin{equation}
    \mathcal{V} := \bigcup_{d,n \in \mathbb{N}} \mathcal{V}_{d \times n}
\end{equation}
as the set of all matrices with non-negative entries of arbitrary finite dimension.

We say that two joint probability distributions $P_{XY}$ and $Q_{X'Y'}$ are similar, denoted
\begin{equation}
    P_{XY} \approx Q_{X'Y'}  \,,
\end{equation}
if one can be obtained from the other by permuting rows or columns, and/or inserting or deleting rows or columns consisting entirely of zeros.
In other words, $P_{XY}$ and $Q_{X'Y'}$ represent the same distribution up to reordering of outcomes and inclusion or exclusion of impossible events. In yet other words, there is $R_{X''Y''}$ such that both $P_{XY}$ and $Q_{X'Y'}$ can be relabeled and embedded into $R_{X''Y''}$ when we lift the notion of relabeling and embedding to non-normalized distributions.
For $P \in \mathcal{V}$, we define 
    $[P] := \{ P' \in \mathcal{V} \mid P \approx P' \}$
as the equivalence class of $P$ under the relation $\approx$. The elements of the semiring are then given by the quotient set
\begin{equation}
    \mathcal{S} := \mathcal{V} / \approx,
\end{equation}
which consists of all equivalence classes of $\mathcal{V}$ under the similarity relation $\approx$.

\item \textbf{The addition ``$+$"}.
Addition in $\mathcal{S} = \mathcal{V}/\!\approx$ is defined via the direct sum of matrices. For two joint distributions $P_{XY}$ and $Q_{X'Y'}$, their direct sum is given by
\begin{equation}
    P_{XY} \oplus Q_{X'Y'} := \begin{pmatrix}
P_{XY} & 0 \\
0 & Q_{X'Y'}
\end{pmatrix}.
\end{equation}
Hence, the sum of two equivalence classes $[P_{XY}]$ and $[Q_{X'Y'}]$ in $\mathcal{S}$ is defined as the equivalence class of the direct sum of any representatives of the respective classes
\begin{equation}
    [P_{XY}] + [Q_{X'Y'}] := [P_{XY} \oplus Q_{X'Y'}] \,.
\end{equation}

\item \textbf{The multiplication ``$\cdot$"}. 
Multiplication in $\mathcal{S} = \mathcal{V}/\!\approx$ is defined via the tensor product (or Kronecker product) of two matrices. For $P_{XY} \in \V_{d \times n}$, $Q_{X'Y'} \in \V_{d' \times n'}$ 
\begin{equation}
P_{XY} \otimes Q_{X'Y'} = \begin{pmatrix}
P(x_1,y_1)Q_{X'Y'} & \hdots & P(x_1,y_n) Q_{X'Y'} \\
\vdots &  & \vdots \\
P(x_d,y_1) Q_{X'Y'} & \hdots & P(x_d,y_n) Q_{X'Y'} 
\end{pmatrix} \,,
\end{equation}
where the product of a scalar and a matrix is defined as entry-wise (i.e., element-wise) multiplication. For example,
\begin{equation}
P(x_1,y_1)Q_{X'Y'} = \begin{pmatrix}
P(x_1,y_1)Q(x'_1,y'_1) & \hdots & P(x_1,y_1) Q(x'_1,y'_{n'}) \\
\vdots &  & \vdots \\
P(x_1,y_1) Q(x'_{d'},y'_1) & \hdots & P(x_1,y_1) Q(x'_{d'},y'_{n'}) 
\end{pmatrix} \,.
\end{equation}
Accordingly, the multiplication of two equivalence classes $[P_{XY}]$ and $[Q_{X'Y'}]$ in $\mathcal{S}$ is defined as the equivalence class of the tensor product of any representatives of the respective classes
\begin{equation}
    [P_{XY}] \cdot [Q_{X'Y'}] := [P_{XY} \otimes Q_{X'Y'}] \,.
\end{equation}

\item \textbf{The zero element ``$0$"}. We define a zero element $0$ in $\V/\!\approx$ as the equivalence class of the $1\times 1$ matrix $0$.

\item \textbf{The unit element ``$1$"}.  We define the unit element $1$ as the equivalence class of the $1\times 1$ matrix $1$.

\item \textbf{The preorder ``$\rgeq$"}. 
\label{item: preorder}
We write
\begin{equation}
\label{eq: relation preorder}
\quad[Q_{X'Y'}]\rgeq[P_{XY}].
\end{equation}
if $P_{XY}$ conditionally majorizes $Q_{X'Y'}$, that is, $P_{XY} \succeq Q_{X'Y'}$ in the sense of Definition~\ref{def:cond-maj-cond-mix} when we lift this definition also for non-normalized distributions. Note that in the latter relation, the roles of $P$ and $Q$ are reversed compared to those in equation~\eqref{eq: relation preorder}.
\end{enumerate}

\subsection{Invariance of the homomorphisms under embedding and shifting}
\label{sec: invariance embedding and shifting}
In this section, we prove that monotone homomorphisms are invariant under embeddings of the $X$ register. This invariance is consistent with the requirement that homomorphisms assign a unique value to each element of an equivalent class. Moreover, we show that they are also invariant under column-wise shifts.
This property will be crucial in characterizing the general form of such homomorphisms, which in turn determines the general form of conditional entropies. 

\bigskip
We begin by introducing some notation. Let $0_{n \times 1}$ denote the column vector of zeros of dimension $n$, $0_{1 \times n}$ the corresponding row vector, and $0_{n \times n}$ the $n \times n$ zero matrix. We denote by $E^{j \times j}_{n \times n}$ the matrix of size $n \times n$ with a 1 in the $(j, j)$-th entry and zeros elsewhere. Finally, $I_{n \times n}$ stands for the $n \times n$ identity matrix.

The next lemma states that there exists a transformation from $P_{XY}\oplus 0_{1\times n} $ and
\begin{equation}
    P_{X'Y} = \begin{pmatrix}
        P_{XY}\\
        0_{1 \times n} 
    \end{pmatrix}
\end{equation}
and vice versa. Hence, since the homomorphisms are additive under direct sums and assign zero to the null element, it follows that they assign the same value to both $P_{XY}$ and $P_{X'Y}$.

\begin{lemma}[Embedding]\label{lemma:embedding}
Let $P_{XY}\in\V_{d\times n}$. Denote the disjoint union of $X$ and $\{1\}$ by $X'$. We have
    \begin{equation}\label{eq:bothways}
P_{XY}\oplus 0_{1\times n}\succeq  \begin{pmatrix}
        P_{XY}\\
        0_{1 \times n}
    \end{pmatrix}\succeq P_{XY}\oplus 0_{1 \times n}.
\end{equation} 
\end{lemma}
\begin{proof}
The proof follows directly by considering a channel acting on the conditioning system that appends and removes outcomes associated with zero entries. Let us construct these channels explicitly. 
Let us denote the matrix in the middle of \eqref{eq:bothways} by $P_{X'Y}$, where $X'$ is the extension of the variable $X$ with one additional outcome. Similarly, let $Y'$ denote the extension of $Y$ by $n$ additional outcomes. Then, $P_{XY} \oplus 0_{1 \times n} \in \mathcal{V}_{X'Y'}$.
We prove the two directions by explicitly writing down the channel that transforms one into the other.

Let us first prove that $P_{XY} \oplus 0_{1 \times n} \succeq P_{X'Y}$. Let us define the matrix that discards $n$ outcomes from $Y'$, namely  
\begin{equation}
    D=\begin{pmatrix}
        I_{n \times n} \\
        I_{n \times n} 
    \end{pmatrix} \,.
\end{equation}
Then, we have that 
\begin{equation}
   (P_{XY} \oplus 0_{1\times n} ) D = \begin{pmatrix}
        P_{XY} & 0_{d \times n}\\
        0_{1 \times n} & 0_{1\times n}
    \end{pmatrix} \begin{pmatrix}
        I_{n \times n} \\
        I_{n \times n} 
    \end{pmatrix} = \begin{pmatrix}
        P_{XY}\\
        0_{1 \times n}
    \end{pmatrix}  = P_{X'Y} \,.
\end{equation}
To prove that $P_{X'Y} \succeq P_{XY} \oplus 0_{1 \times n}$, one can add outcomes that never occur by defining the matrix
\begin{equation}
    D=\begin{pmatrix}
        I_{n \times n} & 0_{n \times n}
    \end{pmatrix} \,.
\end{equation}
Then, we have that 
\begin{equation}
   P_{X'Y} D = \begin{pmatrix}
        P_{XY} \\
        0_{1 \times n} 
    \end{pmatrix} \begin{pmatrix}
        I_{n \times n} & 0_{n \times n}
    \end{pmatrix} = \begin{pmatrix}
        P_{XY} & 0_{d \times n}\\
        0_{1 \times n} & 0_{1 \times n}
    \end{pmatrix}  = P_{XY} \oplus 0_{1\times n} \,.
\end{equation} 
\end{proof}

The next lemma states that the columns can be shifted without affecting the value of any homomorphism, as there exists a channel implementing the transformation in both directions.
\begin{lemma}[Shifting]\label{lemma:shifting}
    Let $P_{XY} \in \V_{d\times n}$. Moreover,  for each $i=1,...,n$, let $P_{X,Y=y_i}$ be the $i$-th column of $P_{XY}$, that is, $P_{XY} = (P_{X,Y=y_1},...,P_{X,Y=y_n})$. Define
    \begin{equation}
\widetilde{P}_{\widetilde{X}Y}=\bigoplus_{i=1}^n P_{X,Y=y_i}= \begin{pmatrix}
P_{X,Y=y_1} & 0_{d \times 1}  & \hdots & 0_{d \times 1} \\
0_{d \times 1} & P_{X,Y=y_2} & \hdots & 0_{d \times 1} \\
\vdots &  & \vdots & \vdots\\
0_{d \times 1} & \hdots & 0_{d \times 1} & P_{X,Y=y_n}\,
\end{pmatrix} \,,
\end{equation} 
where $\widetilde{X}=X\sqcup\cdots\sqcup X$ is the $n$-fold separate union of $X$. We have
\begin{equation}
 \begin{pmatrix}
        P_{XY}\\
        0_{dn \times n}
    \end{pmatrix}\succeq\widetilde{P}_{\widetilde{X}Y}\succeq  \begin{pmatrix}
        P_{XY}\\
        0_{dn \times n}
    \end{pmatrix}.
\end{equation}
\end{lemma}

\begin{proof}
Let us first prove that
\begin{equation}
    \begin{pmatrix}
        P_{XY}\\
        0_{dn \times n}
    \end{pmatrix}\succeq\widetilde{P}_{\widetilde{X}Y}\,.
\end{equation}
For this case, the proof follows by first measuring the $Y$ register and, depending on the outcome, performing an appropriate relabeling of the outcomes in the $X$ register.
 Explicitly, let us define $D^{(j)}= E^{j \times j}_{n \times n}$ which selects the outcome $y_j$ of the $Y$ register. For example, for $j=2$, we have
\begin{equation}
    D^{(2)}=\begin{pmatrix}
        0 & 0 & 0_{1 \times (n-2)}  \\
       0 & 1 & 0_{1 \times (n-2)} \\
       0_{(n-2) \times 1} & 0_{(n-2) \times 1} & 0_{(n-2) \times (n-2)} 
    \end{pmatrix} \,.
\end{equation}
The matrix $S^{(j)}$ relabels the outcomes on the $X$ register. For example, for $j=2$, we have 
\begin{equation}
    S^{(2)}=\begin{pmatrix}
        0_{d \times d} & I_{d \times d} & 0_{d \times d(n-2)} \\
        I_{d \times d} & 0_{d \times d} & 0_{d \times d(n-2)} \\
        0_{d(n-2) \times d} & 0_{d(n-2) \times d} & I_{d(n-2) \times d(n-2)}
    \end{pmatrix} \,.
\end{equation}
To prove the other direction, we use the same matrices as before to select the appropriate outcome of $Y$ and then apply the same permutation to restore the original ordering of the entries.
\end{proof}

\begin{remark}\label{rem:SomOfColumns}
Lemma \ref{lemma:shifting} essentially tells us that, for a matrix $P_{XY}\in\V_{d \times n}$ with columns $P_{X,Y=y_i}$, we may write
\begin{equation}
[P_{XY}]\rgeq[P_{X,Y=y_1}]+\cdots+[P_{X,Y=y_n}]\rgeq [P_{XY}] \,.
\end{equation}
This will subsequently aid us in finding all the monotone homomorphisms of $S$ into $\mb K\in\{\mb R_+,\mb R_+^{\rm op},\mb{TR}_+,\mb{TR}_+^{\rm op}\}$.
\end{remark}

\subsection{Some important properties of the semiring $S$}
To use the large-sample results in Theorem~\ref{thm:Fritz2022}, we must first show that the semiring is of polynomial growth, i.e., it contains a power-universal element. Since the large sample results apply only when the input element is power-universal, we also need to determine exactly which elements have this property. Lastly, we must construct a surjective degenerate homomorphism $\|\cdot\|:\cS\to\R_{>0}^d\cup\{0\}$ with the properties required in Theorem~\ref{thm:Fritz2022}. We begin with the last point. In the following lemma, we show that such a homomorphism can be chosen as the total weight, given by the sum of the matrix entries associated with elements of the semiring. 
Finally, note that the total weight is a scalar quantity rather than a vector, and therefore there exists only a single component homomorphism of \(\|\cdot\|\) in Theorem~\ref{thm:Fritz2022}, namely \(j = 1\).
\begin{lemma}
\label{lem: degenerate and zig-zag}
In the preordered semiring of conditional majorization $S$, the map
\begin{equation}
\|[P_{XY}]\| = \sum_{x\in\mc X}\sum_{y\in\mc Y} P_{XY}(x,y) \,,
\end{equation}
where $P_{XY}$ can be chosen to be any representative of the equivalence class $[P_{XY}]$,
defines a surjective homomorphism $\|\cdot\| : \cS \mapsto \mathbb{R}_+$ with trivial kernel. Moreover, it satisfies
\begin{equation}
[P] \rgeq [Q] \quad \Rightarrow \quad \|[P]\| = \|[Q]\| \quad \Rightarrow\quad  [P] \sim [Q]\,.
\end{equation}
\end{lemma}
\begin{proof}
Clearly, the total weight is a homomorphism since $\|[P]\cdot [Q]\|=\|[P\otimes Q]\|=\|[P]\|\|[Q]\|$ and $\|[P]+[Q]\|=\|[P\oplus Q]\|=\|[P]\|+\|[Q]\|$ for all $[P],[Q]\in\cS$. 

The first direction, namely that $[P] \rgeq [Q] \Rightarrow \|[P]\| = \|[Q]\|$, follows immediately, since conditionally mixing operations preserve the total weight $\|P\|$ of any element $P$. This also shows that the total weight is a degenerate homomorphism since it is invariant under the preorder $\rgeq$.

Let us not turn to the implication $\|[P]\|=\|[Q]\|\Rightarrow [P]\sim [Q]$.  We show that if we define $\|P\|=\|Q\| = c$, then $P\preceq c\succeq Q$ 
and hence $[P]\sim [Q]$. In particular, we show that from the element $[c]$ we can generate both $P$ and $Q$. This is straightforward by considering the matrix with a single non-zero entry $c$ in the top-left corner. Indeed, starting from this matrix, we can construct both $P$ and $Q$. For example, let us focus on $P_{XY} \in \V_{d\times n}$.
Then the steps are the following. Let us denote with $P_{Y}$ the row vectors with components $P_Y(y_i) = \sum_jP_{XY}(x_j,y_i)$. 
First, we create for each outcome $Y$ a column whose first entry is $P_Y(y_i)$ and all other entries are zero. This is done by multiplying on the r.h.s by a row stochastic matrix.
\begin{align}
\begin{pmatrix}
     P_{Y} \\
    0_{(d-1) \times n}   
\end{pmatrix}=
\begin{pmatrix}
    \|P_{XY}\| & 0_{1\times (n-1)} \\
    0_{(d-1)\times 1} & 0_{(d-1)\times (n-1)}
\end{pmatrix} 
\begin{pmatrix}
    \|P_{XY}\|^{-1} P_{Y} \\
    \vdots \\
    \|P_{XY}\|^{-1} P_{Y}    
\end{pmatrix}\,.
\end{align}
Next, we measure the register $Y$ and, conditioned on the outcome $y_i$, we transform the deterministic column (with a single non-zero entry) into a more distributed column $P_{X,Y=y_i}$. Let us denote with $P_{X,Y=y_1}$ the column with components $P_{X,Y=y_1}(x_j)=P_{X,Y}(x_j,y_1)$ and with $P_{X|Y=y_1}$ the column with components $P_{X|Y=y_1}(x_j)=P_{X,Y}(x_j,y_1)/P_Y(y_1)$,
For the first outcome $y_1$, we obtain
\begin{align}
\begin{pmatrix}
     P_{X,Y=y_1} & 0_{d \times (n-1)}  
\end{pmatrix}=
\begin{pmatrix}
    P_{X|Y=y_1} & \star_{1\times (d-1)} 
\end{pmatrix} 
\begin{pmatrix}
     P_{Y} \\
    0_{(d-1) \times n}   
\end{pmatrix}
\begin{pmatrix}
    1 & 0_{1\times (n-1)}\\
    0_{(n-1) \times 1} & 0_{1\times (n-1)} 
\end{pmatrix}
\end{align}
Here, $\star_{1\times (d-1)}$ is any entry that completes the matrix to a doubly stochastic one. By summing over all the outcomes $y_i$ for all $i=1,..,n$, we obtain the desired probability $P_{XY}$. This shows that overall, we have
\begin{equation}
    P_{XY} = \sum_{i=1}^n S^{(i)} P DD^{(i)} \,, 
\end{equation}
where 
\begin{align}
P=\begin{pmatrix}
    \|P_{XY}\| & 0_{1\times (n-1)} \\
    0_{(d-1)\times 1} & 0_{(d-1)\times (n-1)}
\end{pmatrix} \,, \quad 
    D=\begin{pmatrix}
    \|P_{XY}\|^{-1} P_{Y} \\
    \vdots \\
    \|P_{XY}\|^{-1} P_{Y}    
\end{pmatrix} \,,
\end{align}
and for example, for $i=1$, we have
\begin{align}
D^{(1)}=\begin{pmatrix}
    1 & 0_{1\times (n-1)}\\
    0_{(n-1) \times 1} & 0_{1\times (n-1)} 
\end{pmatrix} \,,\quad 
S^{(1)} = \begin{pmatrix}
    P_{X|Y=y_1} & \star_{1\times (d-1)} 
\end{pmatrix}  \,.
\end{align}

\end{proof}

As we mentioned above, to apply the large sample and catalytic results of Theorem~\ref{thm:Fritz2022}, we must also identify a power universal element. By definition, this implies that the semiring is of polynomial growth. As defined in Section~\ref{sec: background semiring}, a power universal element is one that, when provided in sufficient copies, enables the realization of any transformation between elements of the semiring. Within the semiring of conditional majorization, such an element corresponds to any probability distribution exhibiting non-zero randomness in the $X$ register—that is, a probability distribution that is not fully deterministic in $X$. Furthermore, this condition must be satisfied for every possible outcome of the $Y$ register.

\begin{proposition}
\label{prop: power universals}
The following statements are equivalent for $[R_{XY}] \in \cS$:
\begin{enumerate}
    \item \( [R_{XY}] \in \cS \) is power universal.
    \label{item: power universal}
     \item  \( \|[R_{XY}]\| = 1 \), and every non-zero column \( R_{X,Y=y} \) of any representative \( R_{XY} \) contains at least two non-zero entries.
    \label{item: two non-zero entries}
    \item There exists \( r \in (0,1) \) such that
    \begin{equation}\label{eq:pucond}
    \begin{pmatrix}
        r \\
        1 - r
    \end{pmatrix}
    \succeq R_{XY}.
    \end{equation}
    \label{item: majorization R_r}
\end{enumerate}
Consequently, the semiring $S$ is of polynomial growth.
\end{proposition}

\begin{proof}
We prove separately the implications $\ref{item: power universal} \implies \ref{item: two non-zero entries}$, $\ref{item: two non-zero entries} \implies \ref{item: majorization R_r}$ and $\ref{item: majorization R_r} \implies \ref{item: power universal}$. This shows that the statements are all equivalent.
In the following, for brevity, we simply write $P_{XY}$ to refer to a representative of the equivalence class $[P_{XY}]$, without always stating it explicitly. Throughout this proof, we assume that $R_{XY}\in\mc V_{d\times n}$ (i.e.\ $|\mc X|=d$ and $|\mc Y|=n$) and we number $\mc X$ and $\mc Y$ explicitly to aid our matrix calculations.

\bigskip

Let us start with the implication $\ref{item: power universal} \implies \ref{item: two non-zero entries}$. The fact that \( \|[R_{XY}]\| = 1 \) is immediate, since conditional majorization preserves the total weight of elements. If instead \( \|[R_{XY}]\| \neq 1 \), the transformation would not be possible. Next, we show that if $R_{XY}$ contains a column with only one non-zero entry, then it cannot be a power universal. The key idea is that if \( R_{XY} \) contains a column with only one non-zero entry, then any tensor power \( R_{XY}^{\otimes m} \) will also contain a column of the same form. However, starting from a mixed probability distribution, it is impossible to generate a deterministic distribution conditioned on the event that \( Y \) corresponds to such a column. This follows from the fact that doubly stochastic maps cannot produce a deterministic distribution from a non-deterministic one; instead, they only increase the level of mixing.
Explicitly, let us define \( P_{XY} = (1/2, 1/2) \) and take \( Q_{XY} = 1 \) to be the single unit element. Note that we use the subscript $XY$ for both, as the smaller matrix can always be embedded into the space of the larger one by padding with zeros.
Observe that \( Q_{XY} \succeq P_{XY} \).  Assume now that \( R_{XY} \) is power universal. Then, there exists a sufficiently large \( m \in \mathbb{N} \) such that
\begin{equation}
    P_{XY} \succeq Q_{XY} \otimes R_{XY}^{\otimes m} = R_{XY}^{\otimes m}\,.
\end{equation}
This implies that, for a matrix \( R_{XY} \in \V_{d \times n} \), there exist an integer \( k \in \mathbb{N} \), doubly stochastic matrices \( S^{(1)}, \ldots, S^{(k)}\), and (without loss of generality—since the zero entries in the embedding will eliminate the effect of the remaining components) single-row matrices \( D^{(i)} = (d^{(i)}_1, \ldots, d^{(i)}_{n^m}) \),
such that an embedding of the vector \( P_{XY} \) can be transformed into \( R_{XY}^{\otimes m} \) 
\begin{align}\label{eq:apu3}
R_{XY}^{\otimes m}&=\sum_{i=1}^k S^{(i)}
\begin{pmatrix}
    \frac{1}{2}\\
    \frac{1}{2} \\
    0_{(d^m-2)\times 1}
\end{pmatrix}
\big(d^{(i)}_{1}\ \cdots\ d^{(i)}_{n^m}\big)=\sum_{i=1}^k S^{(i)} 
\begin{pmatrix}
    \frac{1}{2}d^{(i)}_{1}\ \cdots\ \frac{1}{2}d^{(i)}_{n^m} \\
   \frac{1}{2} d^{(i)}_{1}\ \cdots\ \frac{1}{2}d^{(i)}_{n^m} \\
   0_{(d^m-2) \times 1}
\end{pmatrix} \,.
\end{align}
We now distinguish two cases: either \( d_j^{(i)} = 0 \) or \( d_j^{(i)} > 0 \). In both situations, the column of \( R_{XY}^{\otimes m} \) with a single non-zero entry cannot be generated by any column of the matrix \( P_{XY} D^{(i)} \). Indeed, if \( d_j^{(i)} = 0 \), the corresponding column is identically zero; if \( d_j^{(i)} > 0 \), then the resulting column contains at least two non-zero entries of the form \( \frac{1}{2} d_j^{(i)} \), and a doubly stochastic matrix cannot reduce this uncertainty to produce a deterministic outcome. This follows from the fact that both the rows and columns of a doubly stochastic matrix must sum to one, thus preserving entropy rather than eliminating it.

\bigskip
Let us now turn to the implication $\ref{item: two non-zero entries} \implies \ref{item: majorization R_r}$
Let us first assume that $\|R_{XY}\|=1$ and $R_{XY}$ is such that its columns have supports with at least two points. We need to find $r\in(0,1)$ such that the condition in item (3) holds. 
The key idea is that each column of \( R_{XY} \) containing at least two non-zero entries can be generated from a vector of the form \( (r, 1 - r) \), for some sufficiently large \( r \in (0,1) \). Such a vector represents an almost deterministic distribution, which—after suitable embedding—can be transformed, via mixing, into any non-deterministic distribution with the same support. Furthermore, the same vector \( (r, 1 - r) \) can be used to generate all columns of \( R_{XY} \), since we are free to enlarge the system \( Y \) to accommodate additional outcomes if needed. Explicitly, 
because of the restriction on $R_{XY}$, we have $R_{XY}(x_j,y_i)<R_Y(y_i)$ for all $x_j\in X$ and $y_i\in Y$ such that the column $R_{X,Y=y_i}$ is non-zero. Therefore, there exists a $r\in[1/2,1)$ such that
\begin{equation}
R_Y(y_i)r\geq R_{XY}(x_j,y_i)\qquad\forall(x_j,y_i)\in X\times Y:\ R_{X,Y=y_i}\neq0.
\end{equation}
By Lemma~\ref{lem: Initial sum condition}, it follows that for each \( y_i \in Y \) with \( R_{X,Y=y_i} \neq 0 \), the two-dimensional embedded vector $\widetilde{R}_Y(y_i) := \big(r R_Y(y_i), (1 - r) R_Y(y_i)\big)$
majorizes the column \( R_{X,Y=y_i} \).
Hence, for each \( y_i \in Y \) with \( R_{X,Y=y_i} \neq 0 \), there exists a doubly stochastic matrix \( S^{(i)} \) and an embedding of the vector \( \widetilde{R}_Y(y_i) \), denoted by \( \widehat{R}_Y(y_i) \), such that $S^{(i)} \widehat{R}_Y(y_i) = R_{X,Y=y_i}$.
To isolate the outcome \( y_i \in Y \), we define the matrices \( D^{(i)} \) with a $1$ in the \( (y_i, y_i) \)-th position and zeros elsewhere. Treating \( R_Y \) as a row-stochastic vector and indexing the outcomes of $Y$ as \( i = 1,..,n \), we obtain
\begin{align}
\left(\begin{array}{c}
r\\1-r
\end{array}\right)&\succeq\left(\begin{array}{c}
r\\1-r
\end{array}\right)R_Y\\
&=\left(\begin{array}{ccc}
rR_Y(y_1)&\cdots&rR_Y(y_n)\\
(1-r)R_Y(y_1)&\cdots&(1-r)R_Y(y_n)
\end{array}\right)\\
&\succeq\sum_{i=1}^n S^{(i)}\left(\begin{array}{ccc}
\widehat{R}_Y(y_1)\cdots \widehat{R}_Y(y_n)
\end{array}\right)R^{(i)}\\
&=\big(S^{(y_1)}\widehat{R}_Y(y_1)\cdots S^{(y_n)}\widehat{R}_Y(y_n)\big)\\
&=\big(R_{X,Y=y_1}\cdots R_{X,Y=y_n}\big)\\
&=R_{XY}.
\end{align}
Hence, we obtain what we wanted to prove.

\bigskip

Finally, let us prove the implication $\ref{item: majorization R_r} \implies \ref{item: power universal}$. We need to show that if there exists $r\in(0,1)$ such that $R_r\succeq R_{XY}$, where $R_r$ is the matrix (column) appearing on the left-hand side of \eqref{eq:pucond},
then $[R_{XY}]$ is power universal. To establish this, we first show that $[R_r]$ is a power universal element; that is, for any elements $[P_{XY}]$ and $[Q_{XY}]$ such that $\|[P_{XY}]\| = \|[Q_{XY}]\|$, there exists a sufficiently large integer $m \in \mathbb{N}$ such that
$P_{XY} \succeq Q_{XY} \otimes R_r^{\otimes m}
$. This, in turn, implies that [$R_{XY}]$ is also power universal, since $R_r\succeq R_{XY}$ implies the chain of inequalities
\begin{equation}
    P_{XY} \succeq Q_{XY} \otimes R_r^{\otimes m}\succeq Q_{XY} \otimes R_{XY}^{\otimes m}\,.
\end{equation}
Note that we use the subscript $XY$ for both $P_{XY}$ and $Q_{XY}$, since the smaller matrix can always be embedded into the larger one by zero-padding.
Without loss of generality, we may assume that $r \in [1/2, 1)$, as the case $r \in (0, 1/2)$ lies in the same equivalence class due to the ability to permute elements.
The fact that $[R_r]$ is power universal follows from the observation that, given a sufficiently large number of copies of $R_r^{\otimes m}$, the non-zero entries of the resulting vector become arbitrarily small. In particular, the largest entry is $r^m$, which approaches zero as $m \to \infty$.
By the initial sum condition for majorization (see Lemma~\ref{lem: Initial sum condition}), this implies that any column of a suitably large embedded distribution $P_{X,Y=y_i}$ can be transformed into $Q_{X,Y=y_i} \otimes R_r^{\otimes m}$, provided the corresponding weights match; that is, when $P_Y(y_i) = Q_Y(y_i)$.
If the column weights differ, we can apply an additional stochastic map to $P_{XY}$ to adjust the marginals accordingly. Explicitly, let $D$ be a row-stochastic matrix such that
\begin{equation}
    P_{Y}D=Q_Y
\end{equation}
and denote $\widetilde{Q}_{XY}:=P_{XY}D$. It is easy to notice that the columns $\widetilde{Q}_{X,Y=y}$ have the same weight as the columns $Q_{X,Y=y}$ for each $y$.
Moreover,  let $m\in\mb N$ be large enough so that
\begin{equation}
r^m\max_{x \in\mc X}Q_{XY}(x,y)\leq \min_{\substack{x \in\mc X \\ \widetilde{Q}(x, y) \neq 0}} \widetilde{Q}(x, y) \qquad\forall y\in\mc Y.
\end{equation}
It then follows, by Lemma~\ref{lem: Initial sum condition}, that for every $y \in\mc Y$, the column $\widetilde{Q}_{X,Y=y}$ majorizes the column $Q_{X,Y=y} \otimes R^{\otimes m}$.
Let $S^{(i)}$ be a doubly stochastic matrix such that
\begin{equation}
    S^{(i)} \widehat{Q}_{X,Y=y_i} = Q_{X,Y=y_i} \otimes R^{\otimes m} \,,
\end{equation}
where we denoted with $\widehat{Q}_{X,Y}$ a suitably embedding of $\widetilde{Q}_{XY}$.
Furthermore, let $D^{(i)}$ denote the matrix that selects the outcome $y_i$, as defined in Lemma~\ref{lemma:embedding}; that is, the matrix with a 1 at entry $(y_i, y_i)$ and zeros elsewhere.
We now have
\begin{align}
P_{XY}&\succeq P_{XY}D=\widetilde{Q}_{XY}\succeq\sum_{i=1}^kS^{(i)}\widehat{Q}_{XY}D^{(i)} =  Q_{XY}\otimes R^{\otimes m}\,.
\end{align}
Finally, the semiring has polynomial growth, as it evidently contains at least one power-universal element.

\end{proof}

Let us note that, whenever $Q_Y$ is a probability distribution, the joint distribution $(1/2,1/2)\times Q_Y$ corresponds to a power universal of the semiring $S$ according to the above result. Note that this is exactly the kind of joint distribution that we use to normalize our general conditional entropies in the normalization axiom 4 of Definition~\ref{def: conditional entropy}.

\subsection{The homomorphisms}
\label{sec: all monotone homorphisms}
In this section, we derive the general form of conditional entropy satisfying the axioms stated in Definition~\ref{def: conditional entropy}. Within the theory of preordered semirings, monotone homomorphisms and derivations are functions designed to satisfy analogous properties. In particular, for the semiring of conditional majorization, these functions correspond—up to an additive constant and a logarithmic factor—to the extremal conditional entropies, which we later denote by $H_{t,\tau}$. As we prove later, these entropies generate the entire family of conditional entropies via convex combinations (see Theorem~\ref{th:barycenter conditional entropies}). Hence, for our purposes, it is crucial to determine the form of the homomorphisms and derivations, which is achieved by characterizing them through their mathematical properties.



As we reviewed in Section~\ref{sec: background semiring}, we need to consider monotone homomorphisms \(\Phi: S \to \mathbb{K}\) in the fields $
    \mathbb{K} = \mathbb{R}_+,\mathbb{R}_+^{\mathrm{op}},\mathbb{TR}_+,\mathbb{TR}_+^{\mathrm{op}} \,.$
The cases $\mathbb{R}_+,\mathbb{R}_+^{\mathrm{op}}$ correspond to the temperate case, whereas $\mathbb{TR}_+,\mathbb{TR}_+^{\mathrm{op}} $ correspond to the tropical cases.
As we show later, the temperate cases correspond to the conditional entropy $H_{t,\tau}$ (up to a constant and a logarithmic factor), with the probability measure $\tau$ and $t>0$ on $\mathbb{R}_+$ and $t<0$ on $\mathbb{R}_+^{\mathrm{op}}$ (see Definition~\ref{def: extremal entropies}). The tropical cases yield the limits $t\to+\infty,-\infty$. In particular, $\mathbb{K} = \mathbb{TR}_+^{\rm op}$ gives rise to the entropies $H_{-\infty,\tau}$ (see Definition~\ref{def: extremal entropies}), corresponding to the limiting form of $H_{t,\tau}$ as $t\to-\infty$.
Finally, the case \(\mathbb{TR}_+\) gives rise to \(H_{+\infty,0}\), which corresponds to the limiting form of \(H_{t,\tau}\) as \(t \to +\infty\) and to the choice of \(\tau\) as a point measure at \(\alpha = 0\) (see Definition~\ref{def: extremal entropies}). As discussed later, in this last limiting regime, the convexity properties of the function enforce that only point measures supported at \(\alpha = 0\) are admissible. In the following result, we express the statement in terms of a measure \(\mu\), which we later decompose as \(\mu = t\,\tau\), where \(t\) denotes the total mass of \(\mu\) and \(\tau\) is the associated probability measure.


\begin{proposition}
\label{prop: monotone homorphisms}
Let \( S\) be the semiring of conditional majorization. The monotone homomorphisms $\Phi:\mc S\to\mb K\in\{\R_+,\R_+^{\rm op},\T\R_+,\T\R_+^{\rm op}\}$ are of the form
\begin{equation}
\label{eq:FinalForm}
\Phi(P_{XY}) =
\begin{cases} 
\displaystyle \sum_{y \in \mathcal{Y}} P_Y(y) \exp\Bigg(\int_{[0,+\infty]} H_\alpha(P_{X|Y=y_i}) \, \mathrm{d}\mu(\alpha) \Bigg), & \mb K\in\{\mathbb{R}_+,\mathbb{R}_+^{\mathrm{op}}\}, \\
\displaystyle \max \limits_{\substack{y \in \mathcal{Y} \\ P_Y( y) > 0}} \exp\Bigg(\int_{[0,+\infty]} H_\alpha(P_{X|Y=y}) \, \mathrm{d}\mu(\alpha) \Bigg), &\mb K=\mathbb{TR}_+^{\mathrm{op}}, \\
\displaystyle \max \limits_{\substack{y \in \mathcal{Y} \\ P_Y( y) > 0}} \exp\left(k H_0(P_{X|Y=y_i}) \, \right), &\mb K=\mathbb{TR}_+ \,,
\end{cases}
\end{equation}
where $\mu$ is a finite positive measure in case $\mb K=\mathbb{R}_+$, and it is a negative measure in cases $\mb K\in\{\mathbb{R}_+^{\mathrm{op}},\T\R_+^{\rm op}\}$. Moreover, in case $\mb K=\mathbb{TR}_+$,  $k$ is a positive constant. 
\end{proposition}


\begin{proof}
The proof is organized into several key steps that lead to the general form of the monotone homomorphisms. Specifically, we proceed as follows:
\begin{enumerate}
    \item Step 1: Additivity. We use the shifting lemma~\ref{lemma:shifting} to reduce the problem to determining the monotone homomorphism for a fixed column vector \( P_{X,Y=y} \), corresponding to a specific value \( Y = y \).
    \item Step 2: Multiplicativity. We then use the property that monotone homomorphisms must be multiplicative under tensor products to further simplify the problem to normalized conditional distributions \( P_{X|Y=y} \).
    \item Step 3: Inner barycenter. We invoke Proposition~\ref{proposition:barycenter entropies}, which states that the most general form of an entropy can be expressed as a linear combination of Rényi entropies. 
     \item Step 4: Simplifying $\mathbb{TR}_+$. We use the fact that, in this case, the monotone homomorphisms are strongly quasi-concave, which allows us to further simplify their form and to single out point measures supported at $\alpha=0$.
\end{enumerate}

\bigskip
\bigskip

 
\noindent\textbf{Step 1: Additivity}. The first step is to use the additivity of the homomorphism under direct sums to reduce the problem from the matrix $P_{XY}$ to a single column vector $P_{X,Y=y}$ for a fixed value $Y = y$.
Lemma~\ref{lemma:shifting}, together with the fact that the homomorphism is additive in both the temperate and tropical sense under summation (see also Remark~\ref{rem:SomOfColumns}), leads to the following additive form
\begin{equation}\label{eq:additivity}
\Phi(P_{XY}) =
\begin{cases}
\sum_{y \in \Y} \Phi(P_{X,Y=y}), & \mathbb{R}_+,\mathbb{R}_+^{\mathrm{op}}, \\
\max_{y \in \Y} \Phi(P_{X,Y=y}), & \mathbb{TR}_+,\mathbb{TR}_+^{\mathrm{op}} \,.
\end{cases}
\end{equation}
Hence, the problem reduces to finding the homomorphisms for the column vector $P_{X,Y=y}$ with fixed $Y=y$.

\bigskip
\bigskip
 \noindent\textbf{Step 2: Multiplicativity}. The second step applies the multiplicativity of the homomorphism under tensor products, specifically scalar multiplication, to reduce the problem from the unnormalized vector $P_{X,Y=y}$ to the normalized vector probability distribution $P_{X|Y=y} = P_{X,Y=y}/P_Y(y)$, where $P_Y(y) = \sum_x P_{XY}(x,y) = P_Y(y)$ is the total weight of the colum $P_{X,Y=y}$. We then obtain
 \begin{equation}\label{eq:normalize}
     \Phi(P_{X,Y=y}) = \Phi\left(P_Y(y) \otimes \frac{P_{X,Y=y}}{P_Y(y)}\right) =\Phi(P_Y(y)) \Phi\left(\frac{P_{X,Y=y}}{P_Y(y)}\right) = \Phi(P_Y(y)) \Phi\big(P_{X|Y=y}\big) \,.
 \end{equation}
 
Next, we analyze the form of $\Phi(P_Y(y))$, which corresponds to the homomorphism applied to the single matrix entry $P_Y(y) \in \mb R_+$. 
By definition, the homomorphism satisfies $\Phi(0) = 0$ for the zero element and $\Phi(1) = 1$ for the unit element.
Moreover, from the multiplicativity of $\Phi$, it follows that $\Phi(xy)=\Phi(x)\Phi(y)$ for all $x,y\in\mb R_+$. Thus, if the restriction of $\Phi$ to $\mathbb{R}_+$ satisfies a suitable regularity condition—such as being bounded on an interval with non-empty interior—it follows that $\Phi(x) = x^\gamma$ for some $\gamma \in \mathbb{R}$.
 We shall next show that $\Phi(x)\leq 1$ for all $x\in[0,1]$. Let $x\in[0,1]$. We easily see that
\begin{equation}\label{eq:x1}
\left(\begin{array}{cc}
x&0\\
0&1-x
\end{array}\right)\succeq\left(\begin{array}{cc}
1&0\\
0&0
\end{array}\right)\succeq\left(\begin{array}{cc}
x&0\\
0&1-x
\end{array}\right) \,.
\end{equation}
The first implication follows from the fact that we can apply a doubly stochastic map to the second column, conditioned on the second outcome of the register $Y$, and then merge both outcomes into a single $Y$ outcome. The second implication follows by applying a stochastic map on $Y$ that discards the second outcome and creates two new outcomes with probabilities $x$ amd $1-x$. Finally, we permute the outcomes of the $X$ register conditioned on the second outcome of the $Y$ register. Since the homomorphism must be monotone under both actions, equality must hold for both elements, and thus, for the temperate case, the sum property implies that
\begin{equation}
1=\Phi(1)=\Phi\left[\left(\begin{array}{cc}
x&0\\
0&1-x
\end{array}\right)\right]=\Phi(x)+\Phi(1-x)\geq\Phi(x) \,,
\end{equation}
since $\Phi(1-x)\geq 0$. The tropical case is similar since $\max\{\Phi(x),\Phi(1-x)\} \geq \Phi(x)$. Hence, the homomorphisms are bounded and, for what was discussed earlier, $\Phi(x)=x^\gamma$ for all $x>0$ for some $\gamma\in\mb R$. Next, we show that $\gamma = 1$ in the temperate cases $\mathbb{R}_+$ and $\mathbb{R}^{\mathrm{op}}_+$, while $\gamma = 0$ in the tropical cases $\mathbb{TR}_+$ and $\mathbb{TR}^{\mathrm{op}}_+$.
By substituting $x = \frac{1}{2}$, and using the equivalence established in \eqref{eq:x1} along with the (temperate or tropical) additivity, we obtain
\begin{equation}
1 = \Phi(1) =
\begin{cases}
2\,\Phi\!\left(\tfrac{1}{2}\right) = 2^{1-\gamma}, & \quad \mathbb{R}_+,\mathbb{R}_+^{\mathrm{op}}, \\
\Phi\!\left(\tfrac{1}{2}\right) = 2^{-\gamma}, & \quad \mathbb{TR}_+,\mathbb{TR}_+^{\mathrm{op}} \,.
\end{cases}
\end{equation}

This establishes that $\gamma = 1$ in the temperate cases $\mathbb{R}_+$ and $\mathbb{R}^{\mathrm{op}}_+$, and $\gamma = 0$ in the tropical cases $\mathbb{TR}_+$ and $\mathbb{TR}^{\mathrm{op}}_+$. In the above arguments, we implicitly assumed that $P_{Y}(y)>0$. In the situation where $P_Y(y)=0$ we have that $\Phi(P_{X,Y=y}) =0$ since $\Phi(0)=0$. This implies that in the cases $\mathbb{TR}_+$ and $\mathbb{TR}^{\mathrm{op}}_+$, only the $Y$ outcomes such that $P_Y(y)>0$ must be included in the maximization. Substituting the values of $\gamma$ into \eqref{eq:additivity}, we obtain
\begin{equation}\label{eq:general1}
\Phi(P_{XY}) =
\begin{cases}
\sum_{y \in \Y} P_Y(y)\,\Phi\!\left(P_{X|Y=y}\right), & \mathbb{R}_+,\mathbb{R}_+^{\mathrm{op}}, \\
\max \limits_{\substack{y \in \mathcal{Y} \\ P_Y( y) > 0}} \Phi\!\left(P_{X|Y=y}\right), & \mathbb{TR}_+,\mathbb{TR}_+^{\mathrm{op}} \,.
\end{cases}
\end{equation}

 \bigskip
 \bigskip
 
\noindent \textbf{Step 3: Inner barycenter}. The final step is to determine the form of the homomorphism for a normalized column vector $P_{X|Y=y}$. Using multiplicativity, for two column vectors $P_X$ and $Q_{X'}$, we have that
\begin{equation}
    \Phi(P_X \otimes Q_{X'}) = \Phi(P_X) \Phi(Q_{X'}) \,.
\end{equation}
Furthermore, by monotonicity under doubly stochastic maps, we have that for any vector $P_X$ and any doubly stochastic matrix $S$,
\begin{equation}
\begin{cases}
 \Phi(P_X)\leq \Phi(SP_X), & \quad \mathbb{R}_+,\mathbb{TR}_+, \\
\Phi(P_X) \geq \Phi(SP_X), & \quad \mathbb{R}_+^{\mathrm{op}},\mathbb{TR}_+^{\mathrm{op}}\,.
\end{cases}
\end{equation}
Therefore, the function $P_X\mapsto\log{\Phi(P_X)}$ is additive under tensor product and monotone under doubly stochastic maps. Using Proposition~\ref{proposition:barycenter entropies}, we obtain that the function $\log{\Phi(P_X)}$ is a linear combination of Rényi entropies. This leads to the form
\begin{equation}
\label{eq:FinalForm proof}
\Phi(P_{XY}) =
\begin{cases} 
\displaystyle \sum_{y \in \mathcal{Y}} P_Y(y) \exp\Bigg(\int_{[0,+\infty]} H_\alpha(P_{X|Y=y}) \, \mathrm{d}\mu(\alpha) \Bigg), & \mathbb{R}_+,\mathbb{R}_+^{\mathrm{op}}, \\
\displaystyle \max \limits_{\substack{y \in \mathcal{Y} \\ P_Y( y) > 0}} \exp\Bigg(\int_{[0,+\infty]} H_\alpha(P_{X|Y=y}) \, \mathrm{d}\mu(\alpha) \Bigg), & \mathbb{TR}_+,\mathbb{TR}_+^{\mathrm{op}}.
\end{cases}
\end{equation}
where $\mu$ is a positive measure in cases $\mathbb{R}_+$ and $\mathbb{TR}_+$, and it is a negative measure in the opposite cases $\mathbb{R}^{\mathrm{op}}_+$ and $\mathbb{TR}^{\mathrm{op}}_+$. 

\bigskip
\bigskip

\noindent \textbf{Step 4: Simplifying $\mathbb{TR}_+$} In the case of $\mathbb{TR}_+$, the form of the homomorphisms can be further restricted. Specifically, only the case $\alpha = 0$ for the inner barycenter is admissible. To show this, let us first prove that for elements of the semiring that are column vectors, the tropical homomorphisms in $\mathbb{TR}_+$ depend only on the support of the vector. Given $P\neq Q$ with the same support, we have that there exists $s,t \in (0,1)$ and $P',Q'$ such that
\begin{align}
   tQ+(1-t)Q'=P\,, \quad  sP+(1-s)P'=Q 
\end{align}
We only need to consider monotone homomorphisms. In \(\mathbb{TR}_+\), Proposition~\ref{prop:convconc} shows that this is equivalent to the homomorphisms being strongly quasi-concave, with the latter property defined in equation~\eqref{eq: def strongly quasi-conc}.
Then, we use the strong quasi-concavity to get
\begin{align}
    \Phi(P) \geq \max \{\Phi(Q),\Phi(Q')\} \geq \Phi(Q) \geq \max \{\Phi(P),\Phi(P')\} \geq \Phi(P) \,.
\end{align}
Hence, $\Phi(P) = \Phi(Q)$, which implies that $\Phi$ depends only on the support of the distribution. Since the only Rényi entropy that depends solely on the support, and not on the specific probabilities, is $H_0$, we can further restrict the monotone homomorphisms in $\mathbb{TR}_+$ to
\begin{equation}
    \Phi(P_{XY}) = \max \limits_{\substack{y \in \mathcal{Y} \\ P_Y( y) > 0}}
    \exp(kH_0(P_{X|Y=y}))
\end{equation}
for some positive constant $k>0$.
All the results together prove the proposition.
\end{proof}

% \subsection{The degenerate homomorphisms}
% \label{sec: degenerate homorphisms}
% In this section, we discuss the degenerate monotone homomorphisms of the semiring of conditional majorization.

% \begin{proposition}
% \label{prop: degenerate homorphisms}
% Let \( S \) denote the semiring of conditional majorization. Then, 
% \begin{enumerate}
%     \item In \(\mathbb{R}_+\) and \(\mathbb{R}_+^{\mathrm{op}}\), the only degenerate monotone homomorphism is the total weight \(\|\cdot\|\) introduced in Lemma~\ref{lem: degenerate and zig-zag}.
%     \item In \(\mathbb{TR}_+\), the only degenerate homomorphism is the function that takes the value \(1\) for all nonzero \(P_{XY}\).
% \end{enumerate}
% \end{proposition}

% \begin{proof}
% Let us first start with the cases \(\mathbb{R}_+\) and \(\mathbb{R}_+^{\mathrm{op}}\).
% The monotone homomorphisms $\Phi:\cS\to\mb K\in\{\mb R_+,\mb R_+^{\rm op}\}$ have the form
% \begin{equation}
% \Phi(P_{XY})=\sum_{y\in\Y}P_Y(y)\exp{\left(\int_{[0,+\infty]}H_\alpha(P_{X|Y=y})\d\mu(\alpha)\right)}
% \end{equation}
% where $\mu$ is either a positive or negative finite measure. Hence, the only possibly degenerate homomorphism corresponds to the measure being zero. In this case, 
% \begin{equation}
% \Phi(P_{XY})=\sum_{y \in\Y}P_Y(y)=\sum_{x\in\mc X}\sum_{y\in\mc Y}P(x,y)=\|P_{XY}\|\,,
% \end{equation}
% where the total weight of the distribution is represented by the degenerate homomorphism 
% $\|\cdot\|$, which characterizes the zig-zag relation $\sim$—was introduced in Lemma~\ref{lem: degenerate and zig-zag}. 
% Indeed, the homomorphism is degenerate, since whenever \(P \succeq Q\), it follows that \(\|P\| = \|Q\|\). This is because they are connected by a channel, which necessarily preserves the total weight.


% Let us turn to the cases $\mathbb{TR}_+, \mathbb{TR}_+^{\mathrm{op}} $.
% In this case where $\Phi \colon \cS \to \mathbb{K}$, with $\mathbb{K} \in \{ \mathbb{TR}_+, \mathbb{TR}_+^{\mathrm{op}} \}$, the monotone homomorphisms take the form
% \begin{equation}
% \Phi(P_{XY})=\max_{y \in \Y}\exp{\left(\int_{[0,+\infty]}H_\alpha(P_{X|Y=y})\d\mu(\alpha)\right)} \,,
% \end{equation}
% where $\mu$ is measure. The only degenerate homomorphism corresponds to the measure being zero, i.e.,
% \begin{equation}
% \Phi(P_{XY})=\left\{\begin{array}{ll}
% \max_{y \in \Y}\,1=1,&{\rm when}\ P_{XY}\neq0,\\
% 0,&{\rm when}\ P_{XY}=0.
% \end{array}\right.
% \end{equation}
% This homomorphism is trivial and carries no meaningful structural information; hence, we shall disregard it. 

% \end{proof}

% \bigskip

% In the next section, our focus therefore shifts to the study of derivations.

\subsection{The derivations}
\label{sec: derivations}
In the previous sections, we examined the monotone homomorphisms, which (up to a constant and a log factor) form a subset of the extremal conditional entropies—those that, through convex combinations, generate the entire set of conditional entropies. However, to obtain a complete characterization of this set, it is also necessary to consider derivations at $\|\cdot\|$, which we prove to arise as limiting cases of homomorphisms. In particular, we show that these derivations coincide with generalized convex combinations of the conditional entropies \(H_{0,\alpha}\) (see Definition~\ref{def: extremal entropies}), corresponding to the limiting form of \(H_{t,\tau}\) as \(t \to 0\).

Let us point out that, according to Theorem \ref{thm:Fritz2022}, we only need to consider derivations $\Delta$ with $\Delta\big(P_X\otimes(1/2,1/2)\big)=1$ for some (and, as we shall see, consequently for all) probability distributions $P_X$ since $P_X\otimes(1/2,1/2)$ is a power universal according to Proposition \ref{prop: power universals}.

\begin{proposition}
\label{prop: monotone derivations}
Let \( S\) be the semiring of conditional majorization. Then, the extremal monotone derivations at \(\|\cdot\|\), as defined in Lemma~\ref{lem: degenerate and zig-zag}, are given by
\begin{equation}
\label{eq:ExtDerivation}
H_{0,\alpha}(P)
= \sum_{y \in \mathcal{Y}} P_Y(y) \, H_\alpha\!\big(P_{X|Y = y}\big)\,,
\end{equation}
for \(\alpha \in [0,+\infty]\).
\end{proposition}
\begin{proof}
We need to study of monotone derivations at $\|\cdot\|$. These are additive maps \(\Delta \colon S \to \mathbb{R}\) that preserve the order, in the sense that \([P] \rgeq [Q]\)—and hence \(P \preceq Q\) for representatives of the corresponding equivalence classes—implies \(\Delta(P) \geq \Delta(Q)\). Moreover, they satisfy the Leibniz rule with respect to the semiring product,
\begin{align}\label{eq:LeibnizRule}
\Delta(P_{XY}\otimes Q_{X'Y'}) 
= \Delta(P_{XY})\,\|Q_{X'Y'}\|
+ \|P_{XY}\|\,\Delta(Q_{X'Y'}) .
\end{align}

for all $P_{XY},Q_{X'Y'}\in\V$. 

The proof is organized into several key steps that lead to the general form of the monotone derivations. Specifically, we proceed as follows:
\begin{enumerate}
    \item Step 1: Additive constant. We begin by showing that the derivations can be taken to satisfy \(\Delta(x) = 0\) for all \(x \in \mathbb{R}_+\).
    \item Step 2: General form. We then exploit the properties of monotone derivations to determine their general form.
    \item Step 3: Extremal derivations. We identify the extremal derivations, which generate the entire set of derivations through convex combinations.
\end{enumerate}


\medskip

\noindent\textbf{Step 1: Additive constant.} We first show that, without loss of generality, the derivation can be assumed to satisfy $\Delta(x) = 0$ for all $x \in \mathbb{R}_+$. To establish this, we prove that from any given derivation $\Delta$, one can construct another derivation $\Delta'$ which satisfies $\Delta'(x) =0$ and hence it is interchangeable with $\Delta$ as defined in Definition 8.2 of \cite{fritz2023abstractII}. Let $\Delta$ be a derivation.  We notice that, for $x,y\in\mb R_+$,
\begin{equation}
\Delta(xy)=\Delta(x)y+x\Delta(y),
\end{equation}
implying that $\delta:=-\Delta|_{\mb R_+}:\mb R_+\to\mb R$ is a derivation.
Define a new map $\Delta' \colon \mathcal{V} \to \mathbb{R}$ by
\begin{equation}
    \Delta'(P) := \Delta(P) - \Delta(\|P\|).
\end{equation}
By construction, $\Delta'(x) = 0$ for all $x \in \mathbb{R}_+$. It therefore suffices to show that $\Delta$ and $\Delta'$ are interchangeable. This will follow if $\Delta'$ is itself a monotone derivation at $\|\cdot\|$. While this is a general fact, we provide a direct proof here for completeness. We have, for all $x,y\in\mb R_+$,
\begin{equation}
\left(\begin{array}{cc}
x&0\\
0&y
\end{array}\right)\succeq\left(\begin{array}{cc}
x+y&0\\
0&0
\end{array}\right)\succeq\left(\begin{array}{cc}
x&0\\
0&y
\end{array}\right),
\end{equation}
by the same argument that we used in the multiplicativity part of Section~\ref{sec: all monotone homorphisms}. Therefore, we obtain that
\begin{equation}
\Delta(x+y)=\Delta\left[\left(\begin{array}{cc}
x&0\\
0&y
\end{array}\right)\right]=\Delta(x)+\Delta(y)\,.
\end{equation}
In the following, for compactness, we omit the subscripts $X,Y$.
We now have, for all $P,Q\in\V$,
\begin{align}
\Delta'(P\oplus Q)&=\Delta(P\oplus Q)-\Delta(\|P\oplus Q\|)\\
&=\Delta(P)+\Delta(Q)-\Delta(\|P\|+\|Q\|)\\
&=\Delta(P)+\Delta(Q)-\Delta(\|P\|)-\Delta(\|Q\|)\\
&=\Delta'(P)+\Delta'(Q)
\end{align}
and
\begin{align}
\Delta'(P\otimes Q)&=\Delta(P\otimes Q)-\Delta(\|P\otimes Q\|)\\
&=\Delta(P)\|Q\|+\|P\|\Delta(Q)-\Delta(\|P\|\|Q\|)\\
&=\Delta(P)\|Q\|+\|P\|\Delta(Q)-\Delta(\|P\|)\|Q\|-\|P\|\Delta(\|Q\|)\\
&=\Delta'(P)\|Q\|+\|P\|\Delta'(Q),
\end{align}
showing that $\Delta'$ is a derivation at $\|\cdot\|$. To prove monotonicity, we consider $[P]$ and $[Q]$ such that $[P]\rgeq[Q]$. Then, we have
\begin{equation}
\Delta'(P)=\underbrace{\Delta(P)}_{\geq\Delta(Q)}-\Delta(\underbrace{\|P\|}_{=\|Q\|})\geq\Delta(Q)-\Delta(\|Q\|)=\Delta'(Q).
\end{equation}
Thus, $\Delta'$ is also a monotone homomorphism, and we may replace $\Delta$ with $\Delta'$. This means that we are free to assume that $\Delta(x)=0$ for all $x\in\mb R_+$.

\medskip
\noindent \noindent\textbf{Step 2: General form.}
Next, we derive the general form of the monotone derivations. First, we derive the action of $\Delta$ on individual normalized columns. Since $\Delta$ satisfies the axioms of an entropy functional, Proposition~\ref{proposition:barycenter entropies} implies the existence of a finite positive measure $\tau \colon \mathcal{B}([0, +\infty]) \to\R_+$ such that for a probability vector $P_X$
\begin{equation}
\Delta(P_X)=\int_{[0,+\infty]}H_\alpha(P_X)\d\tau(\alpha)\,.
\end{equation}
Combining Lemma~\ref{lemma:shifting} with the additivity and monotonicity of $\Delta$, along with the observations made above, we now obtain the following expression for any joint probability distribution $P_{XY}$
\begin{align}
\Delta(P_{XY})&=\sum_{y\in \Y}\Delta(P_{X,Y=y})=\sum_{y\in \Y}\Delta\big(P_Y(y)P_{X|Y=y}\big)\\
&=\sum_{y\in \Y}\underbrace{\Delta\big(P_Y(y)\big)}_{=0}\|P_{X|Y=y}\|_1+P_Y(y)\Delta(P_{X|Y=y})\\
&=\sum_{y\in \Y}P_Y(y)\Delta(P_{X|Y=y})\\
&=\sum_{y\in \Y}P_Y(y)\int_{[0,+\infty]}H_\alpha(P_{X|Y=y})\d\tau(\alpha).
\end{align}
Hence, we obtain the final form for the derivations
\begin{equation}
\label{eq:derivFinalForm} 
    \Delta(P_{XY}) = \sum_{y\in \Y}P_Y(y)\int_{[0,+\infty]}H_\alpha(P_{X|Y=y})\d\tau(\alpha).
\end{equation}
Note that the additivity of $\Delta$ is immediate given the form in \eqref{eq:derivFinalForm}. It is also straightforward to verify that the form above ensures $\Delta$ satisfies the Leibniz rule. Explicitly, if we denote $R_{XX'YY'}:=P_{XY}\otimes Q_{X'Y'}$, the additivity
of the R\'{e}nyi entropies implies that
\begin{align}
\Delta(P_{XY}\otimes Q_{X'Y'})&=\Delta(R_{XX'YY'})\\
&=\sum_{y\in \Y}\sum_{y'\in \Y'}R_{YY'}(y,y')\int_{[0,+\infty]}H_\alpha\big(R_{XX'|YY'=(y,y')}\big)\d\tau(\alpha)\\
&=\sum_{y\in \Y}\sum_{y'\in \Y'}P_Y(y)Q_{Y'}(y')\int_{[0,+\infty]}H_\alpha\big(P_{X|Y=y}\otimes Q_{X'|Y'=y'}\big)\d\tau(\alpha)\\
&=\sum_{y\in Y}\sum_{y'\in \Y'}P_Y(y)Q_{Y'}(y')\int_{[0,+\infty]}H_\alpha\big(P_{X|Y=y}\big)\d\tau(\alpha)\\
&\qquad +\sum_{y\in \Y}\sum_{y'\in \Y'}P_Y(y)Q_{Y'}(y')\int_{[0,+\infty]}H_\alpha\big(Q_{X'|Y'=y'}\big)\d\tau(\alpha)\\
&=\underbrace{\|Q_{Y'}\|_1}_{=\|Q_{X'Y'}\|}\sum_{y\in \Y}P_Y(y)\int_{[0,+\infty]}H_\alpha\big(P_{X|Y=y}\big)\d\tau(\alpha)\\
&\qquad +\underbrace{\|P_Y\|_1}_{=\|P_{XY}\|}\sum_{y'\in \Y'}Q_{Y'}(y')\int_{[0,+\infty]}H_\alpha\big(Q_{X'|Y'=y'}\big)\d\tau(\alpha)\\
&=\Delta(P_{XY})\|Q_{X'Y'}\|+\|P_{XY}\|\Delta(Q_{X'Y'}).
\end{align}
Finally, note that the normalization condition $\Delta\big(P_X\otimes(1/2,1/2)\big)=1$ implies that the positive measure $\tau$ must be a probability measure. 


\medskip
\noindent \noindent\textbf{Step 3: Extremal derivations.} Finally, we identify the extremal derivations that generate the entire set of derivations via convex combinations. For $\alpha\in[0,+\infty]$, let us consider the derivations (or conditional entropies)
\begin{equation}\label{eq:ExtDerivation proof}
H_{0,\alpha}(P)=\sum_{y\in\mc Y} P_Y(y)H_\alpha(P_{X|Y=y})\,.
\end{equation}
From equation~\eqref{eq:derivFinalForm}, we observe that any derivation \(\Delta\) can be written as
\begin{equation}
    \Delta(P) = \int_{[0,+\infty]} H_{0,\alpha}(P)\, d\tau(\alpha) 
\end{equation}
for some probability measure \(\tau\). Hence, each derivation is a convex combination (or, rather, a barycentre) of the extremal derivations \(H_{0,\alpha}(P)\). 
Consequently, since our goal is to characterize the extremal elements of the set—in order to express the general form of conditional entropies through barycentric (i.e., convex) decompositions—it suffices to consider the derivations
\(H_{0,\alpha}\) with \(\alpha \in [0,+\infty]\).
\end{proof}

\section{Sufficient conditions for the parameters}
\label{sec: sufficient conditions measure}
In this section, we establish sufficient conditions for the parameters under which the homomorphisms and derivations (and, consequently, the associated conditional entropies) are monotone under conditional majorization. In the following section, we show that these conditions are also necessary under additional assumptions. These assumptions include, as special cases, discrete measures with finite support and measures that vanish sufficiently fast as $\alpha\to 1$.


In the previous sections, we derived the general form of monotone homomorphisms and derivations, both of which depend on a choice of positive measure. However, not all such choices are valid, as some lead to violations of monotonicity under conditionally mixing channels. In particular, as we show later, while monotonicity under doubly stochastic maps always holds, monotonicity under channels acting on the conditioning system fails for some choices of the parameters. To characterize the admissible cases, we examine the convexity properties of the homomorphism as a function of the joint probability matrix—a standard and often more tractable approach. Indeed, as discussed in Appendix~\ref{app: convexity and monotonicity}, convexity and monotonicity under conditionally mixing channels are essentially equivalent. Moreover, it suffices to establish the convexity properties required to ensure monotonicity under channels acting on the conditioning system, which form a subclass of conditionally mixing channels as characterized in Lemma~\ref{lem: extension to general channels}. Indeed, monotonicity under doubly stochastic maps on the first system is readily verified for all values of the parameters.

\medskip
Let us first identify the function for which we need to establish convexity or concavity properties. We observe that, in the temperate case, the homomorphisms take the form~\eqref{eq:FinalForm}
\begin{align}\label{eq:TemperateForm}
\Phi(P_{XY})&=\Phi_{t,\tau}(P_{XY})=
\sum_{y\in \Y} P_Y(y)\exp{\left(t\int_{[0,+\infty]}H_\alpha(P_{X|Y=y})\d\tau(\alpha)\right)} \,,
\end{align}
where $\tau$ is a probability measure, the normalized version of the measure $\mu$ of \eqref{eq:FinalForm}, and $t=\mu\big([0,+\infty]\big)\in\R\setminus\{0\}$. As detailed in Appendix~\ref{app: convexity and monotonicity}, we need to analyze the convexity properties of the function $P_{XY} \mapsto \Phi(P_{XY})$.
Therefore, in the following analysis, it suffices to study the convexity properties of the function defined for $\vec{x} \in \mathbb{R}^d_+$ by
\begin{equation}\label{eq:conditionfunct}
    \vec{x} \longmapsto \|\vec{x}\|_1 \exp\left( t\int_{[0,+\infty]} H_\alpha(\hat{x}) \,\d\tau(\alpha) \right)\,,
\end{equation}
where $\|\vec{x}\|_1=\sum_i x_i$ and $\hat{x} = \vec{x}/\|\vec{x}\|_1$. We begin with the discrete case, where the measure is supported on a finite set of points. We then extend these results to a general measure by constructing a suitable discrete approximation and taking the limit in which the approximation converges to the original measure. Indeed, convexity and monotonicity properties are preserved under taking limits.

Finally, at the end of the section, we establish sufficient conditions for both the tropical case and the derivations by taking limits of the conditions obtained earlier for the temperate case. Indeed, these two cases can be viewed as limiting regimes of the temperate case, where the total weight of the measure $t$ diverges to infinity or vanishes to zero, respectively.

\subsection{Sufficient conditions for the discrete case}
The discrete case exploits these two basic facts. The first fact is Minkowski's inequality (see e.g.~\cite[Theorem 9]{bullen2013handbook})
\begin{lemma}
\label{lem: Minkowski}
Let $\alpha \in \mathbb{R}$. Then the function, defined on $\mathbb{R}_{>0}^d$,
\begin{equation}
    \vec{x} \longmapsto \left(\sum_{i=1}^d x_i^\alpha\right)^{\frac{1}{\alpha}}
\end{equation}
\begin{enumerate}[itemsep=1pt]
    \item is convex if $\alpha > 1$;
    \item is concave if $\alpha < 1$.
\end{enumerate}
\end{lemma}
The second fact is a well-known fact about convexity properties of multivariate divergences~\cite[Lemma 8]{mu21_renyi} (see also~\cite[Proposition 13]{farooq2024matrix})
\begin{lemma}
\label{lem: Matrix majorization}
Let $N \in \mathbb{N}$ and let $\vec{\beta} = (\beta_1, \dots, \beta_N) \in \mathbb{R}^N$ be such that $\sum_{\ell=1}^N \beta_\ell = 1$. Then the function, defined on $\mathbb{R}_{>0}^N$,
\begin{equation}
    \vec{y} \longmapsto \prod_{\ell=1}^N y_\ell^{\beta_\ell}
\end{equation}
\begin{enumerate}[itemsep=1pt]
    \item is concave if and only if $\beta_\ell \geq 0$ for all $\ell$;
    \label{it: concavity matrix majorization}
    \item is convex if and only if there exists a unique index $k$ such that $\beta_{k} > 0$ and $\beta_\ell \leq 0$ for all $\ell \neq k$.
     \label{it: convexity matrix majorization}
\end{enumerate}
\end{lemma}
We then combine these two facts to prove concavity and convexity properties in the case where $\tau$ is a discrete measure supported on a finite number of points.
\begin{lemma}\label{lem: discrete measure}
Consider a discrete probability measure $\tau:\mc B\big([0,+\infty]\big)\to[0,1]$ and $t\in\R\setminus\{0\}$. Let us assume that $\tau$ is supported on finitely many points $\alpha_1,\ldots,\alpha_N\in[0,\infty]$ enumerated in increasing order. Let us pick $d\in\mathbb{N}$. The function defined on  $\vec{x}\in\R^d_{\geq 0} \setminus\{0\}$
\begin{equation}
\vec{x}\mapsto \|\vec{x}\|_1\exp{\left(t\int_{[0,+\infty]}H_\alpha(\hat{x})\,d\tau(\alpha)\right)}
\end{equation}
where $\hat{x}:=\vec{x}/\|\vec{x}\|_1$ is
\begin{enumerate}
\item concave if $t>0$, $\alpha_1,\ldots,\alpha_N<1$, $\sum_{\ell=1}^N\frac{\alpha_\ell}{1-\alpha_\ell}\tau(\{\alpha_\ell\})\leq 1/t$ and
\item convex if $t<0$ and either
\begin{enumerate}[label=(\alph*),ref=2(\alph*)]
\item $\alpha_1,\ldots,\alpha_N\leq 1$ or
\label{item: sufficient discrete convex 1}
\item $\alpha_1,\ldots,\alpha_{N-1}<1$, $\alpha_N\geq 1$, and $\sum_{\ell=1}^N\frac{\alpha_\ell}{1-\alpha_\ell}\tau(\{\alpha_\ell\})\geq1/t$.
\label{item: sufficient discrete convex 2}
\end{enumerate}
\end{enumerate}
\end{lemma}

\begin{proof}
    The proof leverages the discreteness of the measure in conjunction with Minkowski's inequality and convexity of multivariate divergences in Lemma~\ref{lem: Minkowski} and~\ref{lem: Matrix majorization}.
    Using that $\hat{x}=\vec{x}/\|\vec{x}\|_1$, we can rewrite 
    \begin{align}
    \|\vec{x}\|_1 \exp\left( t\sum_\ell H_{\alpha_\ell}(\hat{x})\,\tau(\{\alpha_\ell\}) \right)
    = \|\vec{x}\|_1^{1 - t\sum_\ell\frac{\alpha_\ell}{1-\alpha_\ell} \,\tau(\{\alpha_\ell\})} \exp\left( t\sum_\ell H_{\alpha_\ell}(\vec{x}) \,\tau(\{\alpha_\ell\}) \right)
\end{align}
as long as $\alpha_\ell\neq 1$ for any $\ell\in\{1,\ldots,N\}$. Hence, we need to study the properties of the function
    \begin{equation}
    \vec{x} \longmapsto \left(\sum_{i} x_i \right)^{1 - t\sum_\ell \frac{\alpha_\ell}{1-\alpha_\ell} \, \tau(\{\alpha_\ell\})} \prod_{\ell} \left(\sum_{i} x_i^{\alpha_\ell} \right)^{\tau(\{\alpha_\ell\})\frac{1}{1-\alpha_\ell}t}  = \left(\sum_{i} x_i \right)^{\bar{\beta}} \prod_{\ell} \left(\sum_{i} x_i^{\alpha_\ell} \right)^{\frac{\beta_\ell}{\alpha_\ell}}\,,
    \end{equation}
    where we defined the exponents $\beta_\ell = t\frac{\alpha_\ell}{1-\alpha_\ell} \, \tau(\{\alpha_\ell\})$ and $\bar{\beta}= 1 - \sum_\ell \beta_\ell $. We discuss the three cases separately. 

\bigskip

    Let us start with the first case. It follows that $\beta_\ell \geq 0$ for all $\ell$. By Lemma~\ref{lem: Minkowski}, the function $\vec{x} \mapsto \big(\sum_{i=1}^d x_i^{\alpha_\ell}\big)^{1/\alpha_\ell}$ is concave. Hence, for two vectors $\vec{x}$ and $\vec{z}$ and a coefficient $\lambda \in [0,1]$, we have that
\begin{align}
    \left(\sum_{i} (\lambda x_i+(1-\lambda)z_i)^{\alpha_\ell}\right)^\frac{1}{\alpha_\ell} \geq \lambda \left(\sum_{i} x_i^{\alpha_\ell}\right)^\frac{1}{\alpha_\ell} + (1-\lambda) \left(\sum_{i}  z_i^{\alpha_\ell}\right)^\frac{1}{\alpha_\ell} \,.
\end{align}
Moreover, the exponent $\beta_\ell \geq 0$, so the corresponding power is monotone. Therefore, we have that
\begin{align}
    \left(\sum_{i} (\lambda x_i+(1-\lambda)z_i)^{\alpha_\ell}\right)^{\frac{\beta_\ell}{\alpha_\ell}} \geq \left(\lambda \left(\sum_{i} x_i^{\alpha_\ell}\right)^\frac{1}{\alpha_\ell} + (1-\lambda) \left(\sum_{i}  z_i^{\alpha_\ell}\right)^\frac{1}{\alpha_\ell}\right)^{\beta_\ell} \,.
\end{align}
Finally, we use the joint concavity of multivariate divergences from Lemma~\ref{lem: Matrix majorization}, item~\ref{it: concavity matrix majorization}. Since it holds that $\beta_\ell >0$, we only need to ensure that $0 \leq \bar{\beta}= 1- \sum_\ell\beta_\ell $. Under this condition, we have that
\begin{align}
    &\left(\sum_{i} \lambda x_i+(1-\lambda )z_i \right)^{\bar{\beta}} \prod_{\ell} \left(\sum_{i} (\lambda x_i+(1-\lambda ) z_i)^{\alpha_\ell} \right)^{\frac{\beta_\ell}{\alpha_\ell}} \\
    &\qquad \geq \lambda \left(\sum_{i} x_i \right)^{\bar{\beta}} \prod_{\ell} \left(\sum_{i} x_i^{\alpha_\ell} \right)^{\frac{\beta_\ell}{\alpha_\ell}} + (1-\lambda) \left(\sum_{i} y_i \right)^{\bar{\beta}} \prod_{\ell} \left(\sum_{i} y_i^{\alpha_\ell} \right)^{\frac{\beta_\ell}{\alpha_\ell}} \,.
\end{align}
The constraint $ 0 \leq \bar{\beta} =1-\sum_\ell\beta_\ell$ yields the condition on the measure appearing on the first item of the lemma. %The other cases follow from a similar argument, and we provide the details below.

\bigskip

 Let us turn to the second case in the statement,~\ref{item: sufficient discrete convex 1}. In this case, the measure is negative. We begin by assuming $\alpha_N<1$. The extension to include the possibility $\alpha_N=1$ can then be obtained via a limit argument (see the discussion below). It follows that $\beta_\ell < 0$ for all $\ell$. By Lemma~\ref{lem: Minkowski}, the function $\vec{x} \mapsto \big(\sum_{i=1}^d x_i^{\alpha_\ell}\big)^{1/\alpha_\ell}$ is concave. Hence, for two vectors $\vec{x}$ and $\vec{z}$ and a coefficient $\lambda \in [0,1]$, we have that
\begin{align}
    \left(\sum_{i} (\lambda x_i+(1-\lambda)z_i)^{\alpha_\ell}\right)^\frac{1}{\alpha_\ell} \geq \lambda \left(\sum_{i} x_i^{\alpha_\ell}\right)^\frac{1}{\alpha_\ell} + (1-\lambda) \left(\sum_{i}  z_i^{\alpha_\ell}\right)^\frac{1}{\alpha_\ell} \,.
\end{align}
Moreover, the exponent $\beta_\ell < 0$, so the corresponding power is antimonotone. Therefore, we have that
\begin{align}
    \left(\sum_{i} (\lambda x_i+(1-\lambda)z_i)^{\alpha_\ell}\right)^{\frac{\beta_\ell}{\alpha_\ell}} \leq \left(\lambda \left(\sum_{i} x_i^{\alpha_\ell}\right)^\frac{1}{\alpha_\ell} + (1-\lambda) \left(\sum_{i}  z_i^{\alpha_\ell}\right)^\frac{1}{\alpha_\ell}\right)^{\beta_\ell} \,.
\end{align}
Finally, we use the joint convexity result from Lemma~\ref{lem: Matrix majorization}, item~\ref{it: convexity matrix majorization}. Since it holds that $\beta_\ell <0$, we only need to ensure that $0 \leq \bar{\beta}= 1- \sum_\ell\beta_\ell $. However, this condition already follows from the fact that $\beta_\ell<0$ and hence we do not impose any additional constraint on the measure.  Under this condition, we have that
\begin{align}
    &\left(\sum_{i} \lambda x_i+(1-\lambda )z_i \right)^{\bar{\beta}} \prod_{\ell} \left(\sum_{i} (\lambda x_i+(1-\lambda ) z_i)^{\alpha_\ell} \right)^{\frac{\beta_\ell}{\alpha_\ell}} \\
    &\qquad \leq \lambda \left(\sum_{i} x_i \right)^{\bar{\beta}} \prod_{\ell} \left(\sum_{i} x_i^{\alpha_\ell} \right)^{\frac{\beta_\ell}{\alpha_\ell}} + (1-\lambda) \left(\sum_{i} y_i \right)^{\bar{\beta}} \prod_{\ell} \left(\sum_{i} y_i^{\alpha_\ell} \right)^{\frac{\beta_\ell}{\alpha_\ell}} \,.
\end{align}
Hence, we proved the second item of the lemma.

\bigskip

 Let us turn to case~\ref{item: sufficient discrete convex 2}. 
 Note that the condition $\sum_\ell\mu(\{\alpha_\ell\}) \frac{\alpha_\ell}{1-\alpha_\ell} \leq \frac{1}{t}$ ensures $\tau(\{\alpha_N\}) \neq 0$.
By Lemma~\ref{lem: Minkowski}, the function $\vec{x} \mapsto \big(\sum_{i=1}^d x_i^{\alpha_\ell}\big)^{1/\alpha_\ell}$ is concave for $\ell\leq N-1$. Hence, for two vectors $\vec{x}$ and $\vec{z}$ and a coefficient $\lambda \in [0,1]$, we have that, for $\ell\leq N-1$,
\begin{align}
    \left(\sum_{i} (\lambda x_i+(1-\lambda)z_i)^{\alpha_\ell}\right)^\frac{1}{\alpha_\ell} \geq \lambda \left(\sum_{i} x_i^{\alpha_\ell}\right)^\frac{1}{\alpha_\ell} + (1-\lambda) \left(\sum_{i}  z_i^{\alpha_\ell}\right)^\frac{1}{\alpha_\ell} \,.
\end{align}
Additionally, for $\alpha_N > 1$, the function is convex and hence
\begin{align}
    \left(\sum_{i} (\lambda x_i+(1-\lambda)z_i)^{\alpha_N}\right)^{\frac{1}{\alpha_N}} \leq \lambda \left(\sum_{i} x_i^{\alpha_k}\right)^{\frac{1}{\alpha_N}} + (1-\lambda) \left(\sum_{i}  z_i^{\alpha_k}\right)^{\frac{1}{\alpha_N}} \,.
\end{align}
Moreover, the exponent $\beta_\ell < 0$ for $\ell\leq N-1$, so the corresponding power is antimonotone. Therefore, we have that
\begin{align}
    \left(\sum_{i} (\lambda x_i+(1-\lambda)z_i)^{\alpha_\ell}\right)^{\frac{\beta_\ell}{\alpha_\ell}} \leq \left(\lambda \left(\sum_{i} x_i^{\alpha_\ell}\right)^\frac{1}{\alpha_\ell} + (1-\lambda) \left(\sum_{i}  z_i^{\alpha_\ell}\right)^\frac{1}{\alpha_\ell}\right)^{\beta_\ell} \,.
\end{align}
Additionally, for $\alpha_N > 1$, the exponent $\beta_N = \tau(\{\alpha_N\}) \frac{\alpha_N}{1-\alpha_N}t > 0$ and hence the corresponding power is monotone. Therefore, we have that
\begin{align}
   \qquad \;\, \left(\sum_{i} (\lambda x_i+(1-\lambda)z_i)^{\alpha_N}\right)^\frac{\beta_N}{\alpha_N} \leq \left(\lambda \left(\sum_{i} x_i^{\alpha_N}\right)^\frac{1}{\alpha_N} + (1-\lambda) \left(\sum_{i}  z_i^{\alpha_N}\right)^\frac{1}{\alpha_N}\right)^{\beta_N} \,.
\end{align}
Finally, we use the joint convexity result from Lemma~\ref{lem: Matrix majorization}, item~\ref{it: convexity matrix majorization}. Since it holds that $\beta_\ell <0$ for $\ell\leq N-1$ and $\beta_N >0$, we only need to ensure that $ \bar{\beta}= 1- \sum_\ell\beta_\ell \leq 0$.  Under this condition, we have that
\begin{align}
    &\left(\sum_{i} \lambda x_i+(1-\lambda )z_i \right)^{\bar{\beta}} \prod_{\ell} \left(\sum_{i} (\lambda x_i+(1-\lambda ) z_i)^{\alpha_\ell} \right)^{\frac{\beta_\ell}{\alpha_\ell}} \\
    &\qquad \leq \lambda \left(\sum_{i} x_i \right)^{\bar{\beta}} \prod_{\ell} \left(\sum_{i} x_i^{\alpha_\ell} \right)^{\frac{\beta_\ell}{\alpha_\ell}} + (1-\lambda) \left(\sum_{i} y_i \right)^{\bar{\beta}} \prod_{\ell} \left(\sum_{i} y_i^{\alpha_\ell} \right)^{\frac{\beta_\ell}{\alpha_\ell}} \,.
\end{align}
The constraint $ \bar{\beta} =1-\sum_\ell\beta_\ell \leq 0$ yields the condition on the measure appearing on the third item of the lemma. 

\bigskip

Finally, note that the endpoints of the interval can be included—provided they are not excluded by the condition on the measure—since the pointwise limit of convex functions is itself convex. For instance, in the second case where $t < 0$ and $\alpha_\ell\leq 1$, the endpoint $\alpha_N = 1$ can be included by observing that
\begin{align}
   &\|\vec{x}\|_1^{1 - t\sum_\ell\frac{\alpha_\ell}{1-\alpha_\ell} \,\tau(\{\alpha_\ell\})} \exp\left( t\sum_\ell H_{\alpha_\ell}(\vec{x}) \,\tau(\{\alpha_\ell\}) + tH(\hat{x}) \,\mu(\{1\})\right)\\
   & \qquad = \lim_{\alpha' \rightarrow 1^-} \|\vec{x}\|_1^{1 - t\sum_{\ell=1}^{N-1}\frac{\alpha_\ell}{1-\alpha_\ell} \,\tau(\{\alpha_\ell\})-\frac{t\alpha'}{1-\alpha'} \,\tau(\{1\})} \exp\left( t\sum_\ell H_{\alpha_\ell}(\vec{x}) \,\tau(\{\alpha_\ell\}) 
    +tH_{\alpha'}(\vec{x}) \,\tau(\{1\})\right) \,,
\end{align}
and applying the convexity result for each $\alpha' < 1$. This indeed implies convexity in the limit as $\alpha' \to 1^-$.

\end{proof}

\subsection{Sufficient conditions for the temperate homomorphisms - general measure}
In this section, we extend the convexity result to the case where $\tau$ is a general probability measure. The idea is to approximate $\tau$ by a sequence of finitely supported probability measures $(\tau_n)_{n \in \mathbb{N}}$, obtained by discretizing the intervals with increasingly fine partitions. Convexity holds for each $\tau_n$ by the result from the previous section. Since $(\tau_n)$ converges to $\tau$ (in a suitable sense), and convexity is preserved under pointwise limits, the result carries over to the continuous case. It turns out that the condition
\begin{equation}\label{eq:IntegralCond}
\int_{[0,+\infty]}\frac{\alpha}{1-\alpha}\,\d\tau(\alpha)\leq\frac{1}{t}
\end{equation}
is crucial.  Here, we adopt the convention that the integral equals $+\infty$ if the measure $\tau$ has support at \(\alpha = 1\).


\begin{proposition}\label{prop: sufficient conditions}
Let $\tau:\mc B\big([0,+\infty]\big)\to\mb R$ be a probability measure, $t\in\R\setminus\{0\}$. The function defined on $\vec{x}\in\R^d_{\geq 0}$ 
\begin{equation}
\vec{x} \longmapsto \|\vec{x}\|_1 \exp\left( t\int_{[0,+\infty]} H_\alpha(\hat{x}) \,\d\tau(\alpha) \right)
\end{equation}
where $\hat{x}=\vec{x}/\|\vec{x}\|_1$ is
\begin{enumerate}
\item concave if $t>0$, $\;\textup{supp}(\tau) \subseteq [0,1]$, and \eqref{eq:IntegralCond} holds and 
\label{it: first sufficient}
\item convex if $t<0$ and either
\begin{enumerate}[label=(\alph*),ref=2(\alph*)]
 \item $\textup{supp}(\tau)\subseteq[0,1]$ or
\label{it: second sufficient (a)}
 \item there is $\alpha_*\in(1,+\infty]$ such that $\textup{supp}(\tau)\subseteq[0,1]\cup\{\alpha_*\}$ and \eqref{eq:IntegralCond} holds.
\label{it: second sufficient (b)}
\end{enumerate}
\end{enumerate}
\end{proposition}
Note that in the first case of  $t<0$, the measure $\tau$ may have support at $\alpha=1$.  In contrast, in the second case of \(t < 0\) and for \(t > 0\) the measure cannot have support at \(\alpha = 1\). Indeed, our convention implies that the integral diverges, and consequently, the bound in \eqref{eq:IntegralCond} would be violated.


\begin{proof}
We begin by considering the first case. The idea is to construct a discrete approximation that enables the application of the result from the previous section for each finite level of the approximation. We then take the limit as the approximation becomes increasingly fine, observing that convexity is preserved in the limit. 

To show that the approximation converges to the desired function, we apply the dominated or monotone convergence theorem, although, as the functions that we study are monotone, this convergence is rather a simple fact. For this, we rely on the monotonicity of both the Rényi entropies $\alpha \mapsto H_\alpha(\hat{x})$ and the function $\alpha \mapsto \frac{\alpha}{1 - \alpha}$. Our goal is to apply the monotone convergence theorem to approximate the integral using a sequence of discrete functions $f^{(N)}$
\begin{equation}
\int_{[0,1]} f(\alpha)\d\tau(\alpha) = \lim_{N \to \infty} \int_{[0,1]} f^{(N)}(\alpha) \d\tau(\alpha).
\end{equation}
To invoke this theorem, we must construct a sequence of functions $(f^{(N)})_{N \in \mathbb{N}}$ that satisfy the monotonicity assumption: each $f^{(N)}$ takes only finitely many values (i.e.,\ $f^{(N)}$ is a simple function), there is a $\tau$-integrable function $g$ such that $f^{(N)}(\alpha) \leq g(\alpha)$ for all $\alpha \in [0,1]$ and $N\in\mb N$, and converges pointwise, that is, $f^{(N)}(\alpha) \to f(\alpha) $ as $ N \to \infty$. 

We want to construct monotone (or at least dominated) pointwise approximations for both the integral $\int_{[0,1]} H_\alpha(\hat{x}) \,\d\tau(\alpha)$ and the constraint $\int_{[0,1]} \frac{\alpha}{1 - \alpha} \,\d\tau(\alpha)\leq 1/t$. For the latter one, because of the inequality involved, we want to use the monotone convergence theorem and approach to the integral from below. As we show below, the monotonicity of the approximations of the latter constraint is needed to ensure that the condition on the parameters for a discrete measure in Lemma~\ref{lem: discrete measure} is always satisfied, so that we can apply the result to each discrete approximation and eventually take the limit. 
Note that the assumption on the measure implies $\tau(\{1\}) = 0$, which is why we include $\alpha = 1$ in the integration interval.

The discrete approximation $f^{(N)}$ is done using a staircase construction.
For each $N \in \mathbb{N}$, let $\{b_n^{(N)}\}_{n=1}^N \subset \mathbb{R}_+$ be a strictly increasing sequence with $b_1^{(N)} = 0$.
The staircase approximation is then given by
\begin{equation}
f^{(N)}(\alpha) := \sum_{n=1}^{N-1} b_{n+1}^{(N)}  \chi_{f^{-1}\big ([b_n^{(N)}, b_{n+1}^{(N)})\big)}(\alpha),
\end{equation}
Note that we are converging to $f$ from the above; for the integral involving $\alpha\mapsto\alpha/(1-\alpha)$ we will approach from below. Here $\chi_A(\alpha)$ is the characteristic function of a measurable set $A \subseteq [0,1]$, defined by
\begin{equation}
    \chi_A(\alpha) = 
  \begin{cases}
  1 & \text{if } \alpha \in A, \\
  0 & \text{otherwise,}
  \end{cases}
\end{equation}
 and $f^{-1}[b_n, b_{n+1})$ is the preimage of the interval $[b_n, b_{n+1})$, that is,
 \begin{equation}
     f^{-1}\big([b_n, b_{n+1})\big) := \{\alpha \in [0,1] \mid f(\alpha) \in [b_n, b_{n+1}) \}.
 \end{equation}
Hence, by the dominated convergence theorem, we obtain the approximation
\begin{equation}
\int_{[0,1]} f(\alpha) \d\tau(\alpha) = \lim_{N \to \infty} \sum_{n=1}^{N-1} b_n^{(N)}  \tau\left(f^{-1}\left([b_{n+1}^{(N)}, b_{n+1}^{(N)})\right)\right).
\end{equation}
In our case, we set $f = \mathrm{ev}_{\hat{x}}$, defined by $\mathrm{ev}_{\hat{x}}(\alpha) := H_\alpha(\hat{x})$. In this case, $f\leq H_0(\hat{x})$ and, as $\tau$ is a probability measure, the constant upper bound is $\tau$-integrable, and we may use the dominated convergence theorem.
Hence, the staircase approximation yields a rigorous way to evaluate the integral $\int_{[0,1]} H_\alpha(\hat{x}) \,\d\mu(\alpha)$ using only simple functions.% and Lebesgue measure theory.

We now use the result of the previous section for finite approximation to establish the result for the general case.
Possibly excluding a few initial indices $n = 1, \ldots, m$, we may assume that for each $n = 1, \ldots, N - 1$, there exist parameters $\alpha^{(N)}_n \in [0,1)$ such that $b^{(N)}_n = H_{\alpha^{(N)}_m}(\hat{x})$; note that, since $\alpha\mapsto H_\alpha(\hat{x})$ is non-increasing, the indexing of $b^{(N)}_n$ and $\alpha^{(N)}_m$ runs oppositely. In the following, we disregard these initial exceptions and relabel the sequence accordingly. We define the intervals
\begin{align}
    &I^{(N)}_{n}={\rm ev}_{\hat{p}}^{-1}\big(\big[b^{(N)}_n,b^{(N)}_{n+1}\big)\big)=\big[\alpha^{(N)}_n,\alpha^{(N)}_{n+1}\big) \quad \text{for} \;\; n=1,\ldots,N-2 \\
    & I^{(N)}_{N-1} = {\rm ev}_{\hat{p}}^{-1}\big([b^{(N)}_{N-1},+\infty\big)\big)=\big[\alpha^{(N)}_{N-1},1\big)
\end{align}
With these definitions, the above convergence arguments imply that $\Phi_N(\vec{x})\to\Phi(\vec{x})$ as $N\to\infty$ where
\begin{equation}
\Phi_N(\vec{x}):=\|\vec{x}\|_1\exp{\left(t\sum_{n=1}^{N-1}b_n^{(N)}\tau\big(I^{(N)}_n\big)\right)}.
\end{equation}
Denoting $\tau^{(N)}_n:=\tau\big(I^{(N)}_n\big)$, and $\beta^{(N)}_n:=\tau^{(N)}_nt\alpha^{(N)}_n/(1-\alpha^{(N)}_n)$, we have
\begin{align}
\Phi_N(\vec{x})&=\|\vec{x}\|_1\exp{\left(t\sum_{n=1}^{N-1}\tau^{(N)}_n b^{(N)}_n\right)}\\
&=\|\vec{x}\|_1\exp{\left(t\sum_{n=1}^{N-1}\tau^{(N)}_n H_{\alpha^{(N)}_n}(\hat{x})\right)}\\
&=\|\vec{x}\|_1\exp{\left(t\sum_{n=1}^{N-1}\frac{\tau^{(N)}_n}{1-\alpha^{(N)}_n}\log{\sum_{i=1}^d\hat{x}_i^{\alpha^{(N)}_n}}\right)}\\
&=\|\vec{x}\|_1\prod_{n=1}^{N-1}\left(\sum_{i=1}^d\hat{x}_i^{\alpha^{(N)}_n}\right)^{\beta^{(N)}_n/\alpha^{(N)}_n}\\
&=\|\vec{x}\|_1^{\beta^{(N)}_0}\prod_{n=1}^{N-1}\left(\sum_{i=1}^d x_i^{\alpha^{(N)}_n}\right)^{\beta^{(N)}_n/\alpha^{(N)}_n}\label{eq:PhiN}
\end{align}
where $\beta_0^{(N)}:=1-\beta^{(N)}_1-\cdots-\beta^{(N)}_N$ and in the last line, we used the relation between the normalized and unnormalized vectors, namely $\hat{x}_i = x_i / \|\vec{x}\|_1$.

To prove concavity, we need to show that for two vectors $\vec{w}$ and $\vec{z}$ such that
$\vec{x}=t\vec{w}+(1-t)\vec{z}$, we have
\begin{equation}\label{eq:Phiconc}
\Phi(\vec{x})\geq t\Phi(\vec{w})+(1-t)\Phi(\vec{z}).
\end{equation}
As done for $\Phi(\vec{x})$ above, we define similar approximations $\Phi_N(\vec{w})$ and $\Phi_N(\vec{z})$ for $\Phi(\vec{w})$ and $\Phi_N(\vec{z})$. Using the form~\eqref{eq:PhiN} and Lemma~\ref{lem: discrete measure} for discrete measures, we obtain that
\begin{equation}
\Phi_N(\vec{x})\geq t\Phi_N(\vec{w})+(1-t)\Phi_N(\vec{z})
\end{equation}
for 
\begin{equation}
\sum_{n=1}^{N-1}\frac{\alpha^{(N)}_n}{1-\alpha^{(N)}_n}\tau_{n}^{(N)} \leq \frac{1}{t} \,.
\end{equation}
We then take the limit on both sides of the above inequality to establish the desired statement. Moreover, the condition on the measure becomes
\begin{equation}
\int_{[0,1]}\frac{\alpha}{1-\alpha}\d\tau(\alpha) \leq \frac{1}{t} \,.
\end{equation}
Indeed, we may once again apply the monotone convergence theorem by observing that the function $f : [0,1) \to \mathbb{R}_+$, given by $f(\alpha) = \frac{\alpha}{1 - \alpha}$, is non-decreasing. Consider the approximation $f^{(N)} : [0,1) \to \mathbb{R}_+$ defined as
\begin{equation}
    f^{(N)} = \sum_{n=1}^{N-1} \frac{\alpha_n^{(N)}}{1 - \alpha_n^{(N)}} \chi_{I_n^{(N)}} 
\end{equation}
which satisfies $f^{(N)} \leq f$ for all $N \in \mathbb{N}$. Then,
\begin{equation}
    \sum_{n=1}^{N-1} \frac{\tau_n^{(N)} \alpha_n^{(N)}}{1 - \alpha_n^{(N)}} = \int_{[0,1)} f^{(N)} \,\d\tau \leq \int_{[0,1]} \frac{\alpha}{1 - \alpha} \,\d\tau(\alpha) \leq 1/t\,.
\end{equation}
The second case follows from a similar argument. 

The final case is likewise analogous, where we isolate the point mass at $\alpha_* > 1$ and apply the same approximation to the part of the measure supported on $[0,1)$. It is worthwhile to note that case~\ref{it: second sufficient (a)} can be seen as a limiting argument from~\ref{it: second sufficient (b)} when $\alpha_*\to 1^{+}$: We let $\alpha_*>1$ with $\tau(\{\alpha_k\})>0$ converge to 1. At some point, when $\alpha_*$ is close enough to 1, the integral condition in~\ref{it: second sufficient (b)} becomes trivial. We also get the case $\alpha_*=\infty$ in~\ref{it: second sufficient (b)} through a limit argument.
\end{proof}

\subsection{Monotonicity under doubly stochastic maps}
\label{sec: monotonicty doubly stochastic}
As mentioned at the beginning of this section, monotonicity under doubly stochastic maps acting on the first register \(X\) always holds, and its proof is straightforward.
Indeed, since the exponential function is monotone and the map
\begin{equation}
    P_X \mapsto \int_{[0,+\infty]} H_\alpha(P_X) \,\d\mu(\alpha)
\end{equation}
is monotone under doubly stochastic maps (as each $H_\alpha(P_X)$ individually is), it follows directly from the explicit form of $\Phi$ in~\eqref{eq:FinalForm} that $\Phi$ is monotone as well. 

Hence, under the above condition on the measure, Lemma~\ref{lem: extension to general channels} ensures that the monotone homomorphisms are monotone under conditionally mixing channels as required for a conditional entropy in Definition~\ref{def: conditional entropy}. 



\subsection{Sufficient conditions for tropical homorphisms}
In this section, we establish sufficient conditions on the measure for the tropical homomorphisms. The key observation is that monotonicity under conditionally mixing channels is preserved under limits; that is, the limit of a monotone quantity remains monotone. This is because taking limits on both sides of an inequality does not alter its direction. Consequently, the convexity and concavity conditions established above, which are equivalent to monotonicity for temperate homomorphisms (see Appendix~\ref{app: convexity and monotonicity}), provide the basis for extending monotonicity properties to tropical homomorphisms.

In the following, we establish a monotonicity result for the conditional entropies that can be constructed from the tropical homomorphisms. We start with the case $\mb{TR}_+^{\rm op}$.
\begin{proposition}
\label{prop: sufficient tropical}
Let $\tau:\mc B\big([0,+\infty]\big)\to\mb R$ be a probability measure. Then, the function
\begin{align}\label{eq:TropicalForm}
     P_{XY} \longmapsto &\min \limits_{\substack{y \in \mathcal{Y} \\ P_Y( y) > 0}}  \int_{[0,+\infty]} H_\alpha(P_{X|Y=y}) \,\d\tau(\alpha)
\end{align}
is monotone under conditionally mixing channels acting on $P_{XY}$ if one of the following holds:
\begin{enumerate}[label=(\alph*),ref=2(\alph*),itemsep=0pt,topsep=5pt]
    \item  $\;\textup{supp}(\tau) \subseteq [0,1]$.
    \label{it: 1 for tropical}
    \item $\;\operatorname{supp}(\tau) \subseteq [0,1] \cup \{\alpha_*\}$, $\;\alpha_* \in (1, +\infty]$, $\;\displaystyle\int_{[0,+\infty]} \frac{\alpha}{1 - \alpha} \, \mathrm{d}\tau(\alpha) \leq 0$.
    \label{it: 2 for tropical}
\end{enumerate}
\end{proposition}
 Note that the result below, according to Proposition~\ref{prop:convconc}, implies the quasi-convexity of the tropical homomorphisms in $\mb{TR}_+^{\rm op}$.
\begin{proof}
    The conditional entropies constructed from the tropical homomorphisms can be obtained by letting $t\to-\infty$. Thus, to obtain sufficient conditions, it is enough to take these limits and verify which of the conditions we derived in the previous sections on the measure continue to hold. 
    
    Explicitly, we have the limit
    \begin{align}
    \label{eq: limit tropical}
&\lim_{t \to -\infty} \frac{1}{t} 
\log\left( \sum_{y \in \mathcal{Y}} P_Y(y) 
\exp\left( t \int_{[0,+\infty]} H_\alpha\big(P_{X|Y = y}\big) \,\mathrm{d}\tau(\alpha) \right) \right) \\
&\hspace{15em}= \min \limits_{\substack{y \in \mathcal{Y} \\ P_Y( y) > 0}} \int_{[0,+\infty]} H_\alpha\big(P_{X|Y = y}\big) \,\mathrm{d}\tau(\alpha) \,.
\end{align}
Let us start with the first case, i.e., we want to show that monotonicity holds whenever the probability measure $\mu$ satisfies $\operatorname{supp}(\mu) \subseteq [0,1]$. This we get directly when we combine the limit in \eqref{eq: limit tropical} to case~\ref{it: second sufficient (a)} of Proposition \ref{prop: sufficient conditions}.

Let us now turn to the second case.
In this case, we assume that the conditions of item~\ref{it: 2 for tropical} of the claim hold for a probability measure $\mu$.  
It is straightforward to see that for sufficiently large $|t|$ as $t\rightarrow-\infty$, the following condition~\ref{it: second sufficient (b)} on convexity in Proposition~\ref{prop: sufficient conditions} holds:
\begin{equation}
\label{eq: condition mu for M}
t<0,\;\operatorname{supp}(\tau) \subseteq [0,1] \cup \{\alpha_*\},\;
\alpha_* \in (1, +\infty],\;
\int_{[0,+\infty]} \frac{\alpha}{1 - \alpha} \,\mathrm{d}\tau(\alpha) \leq \frac{1}{t} \,.
\end{equation}
Hence, we obtain the desired result.
\end{proof}
Next, we turn to $\mb{TR}_+$.
\begin{proposition}
\label{prop: sufficient tropical no}
The function $P_{XY} \longmapsto \max_{y \in \Y}   H_0(P_{X|Y=y})$
is monotone under conditionally mixing channels acting on $P_{XY}$. 
\end{proposition}
\begin{proof}
Also in this case, the proof follows by a limit argument. We have the limit
    \begin{align}
    \label{eq: limit tropical no}
&\lim_{t \to +\infty} \frac{1}{t} 
\log\left( \sum_{y \in \mathcal{Y}} P_Y(y) 
\exp\left( t  H_0\big(P_{X|Y = y}\big) \, \right) \right) = \max \limits_{\substack{y \in \mathcal{Y} \\ P_Y( y) > 0}}  H_0\big(P_{X|Y = y}\big)  \,.
\end{align}
The result then follows by combining the limit in \eqref{eq: limit tropical no} with case~\eqref{it: first sufficient} for $\alpha = 0$ in Proposition~\ref{prop: sufficient conditions}, together with the fact that the limit of a monotone quantity remains monotone.
\end{proof}

\begin{remark}
    We note that unless $\alpha=0$, the condition~\eqref{it: first sufficient} in Proposition~\ref{prop: sufficient conditions} 
cannot be used to obtain an alternative set of sufficient conditions for $\mathbb{T}\mathbb{R}_+$, since the limit $t\to+\infty$ eventually violates the integral constraint. Since such limits give rise to monotone homomorphisms in $\mathbb{T}\mathbb{R}_+$, this observation is consistent with the fact that only the monotone homomorphism with $\alpha = 0$ is admissible in $\mathbb{T}\mathbb{R}_+$ as established in Section~\ref{sec: all monotone homorphisms}.
\end{remark}






\subsection{Sufficient conditions for derivations}
We now turn to the derivations. 
\begin{proposition}
 \label{prop: sufficient derivations}
 The function
 \begin{align}
      P_{XY} \longmapsto &\sum_{y \in \Y} P_Y(y)  H_\alpha(P_{X|Y=y}) \,,
 \end{align}
 is monotone under conditionally mixing channels acting on $P_{XY}$ if $\alpha \in [0,1]$.
 \end{proposition}
 According to Lemma~\ref{lem: conv derivations}, which shows that convexity and monotonicity are equivalent, the above monotonicity result also establishes the convexity of the function.
 \begin{proof}
 The derivations can be obtained by choosing a single point measure and letting the total weight of the measure tend to $0$. Thus, to obtain sufficient conditions, it is enough to take these limits and verify which of the conditions we derived in the previous sections on the measure continue to hold. 
    
Explicitly, we have the limit
    \begin{align}
     \label{eq: limit derivations}
 &\lim_{t \to 0} \frac{1}{t} 
 \log\left( \sum_{y \in \mathcal{Y}} P_Y(y) 
 \exp\left( t H_\alpha\big(P_{X|Y = y}\big)  \right) \right) = \sum_{y \in \mathcal{Y}} P_Y(y)  H_\alpha\big(P_{X|Y = y}\big)  \,.
 \end{align}
 We want to show that monotonicity holds whenever $\alpha \in [0,1]$.
 The monotonicity in this case can be recovered from the conditions on the temperate case in Proposition~\ref{prop: sufficient conditions}, item~\ref{it: second sufficient (a)} by considering the limit $t\to 0$. Explicitly, we need to choose a single point measure concentrated at $\alpha=1$.
We next apply the convexity conditions from Proposition~\ref{prop: sufficient conditions}, item~\ref{it: second sufficient (a)}, which are equivalent to monotonicity, to the measure $\tau$.
Combining these conditions with the limit in~\eqref{eq: limit derivations}, and using that monotonicity is preserved under limits, we obtain monotonicity whenever $\alpha \in [0,1]$.
\end{proof}


The results of this section thus single out the admissible homomorphisms and derivations which, up to logarithmic and constant factors, coincide with the admissible conditional entropies. These extremal conditional entropies will later appear in Theorem~\ref{th:barycenter conditional entropies}, which states the general form of conditional entropies.

\section{Necessary conditions for the parameters}
\label{sec: necessary conditions measures}
In this section, we show that the sufficient conditions on the parameters derived in the previous section that ensure monotonicity under conditionally mixing channels are also necessary, provided some additional assumptions hold. In particular, we assume that the measure vanishes as $\alpha \to 1$ faster than $1/(1-\alpha)$; this class of measures includes, as a special case, discrete measures with finite support. Although these regularity assumptions are required to make our analysis rigorous, we expect---as discussed in the later sections---that the sufficient conditions are in fact necessary in full generality.


Our approach consists in identifying an explicit perturbation direction and computing the second derivative of the function along this direction. By analyzing the sign of the second derivative, we show that outside the range specified in Proposition~\ref{prop:convconc}, the function fails to exhibit the expected convexity or concavity (which is equivalent to monotonicity under conditionally mixing channels). 
Finally, we also note that the problem is inherently challenging: obtaining a sharp characterization of the admissible range requires constructing counterexamples in increasingly high dimensions. Indeed, our numerical investigations indicate that no such violations occur in low-dimensional settings. While the proofs are rather intricate, the underlying idea is simple. We therefore provide a sketch of the proof strategy used to establish the necessary conditions.

\bigskip 
\textbf{Proof idea}. The main idea is to leverage the counterexamples for the convexity properties of the function
\begin{equation}
\label{eq: multivariate divergences function}
    \vec{v} \longmapsto v_1^{\beta_0} v_2^{\beta_1} \cdots v_d^{\beta_d} 
\end{equation}
with $\beta_0=1-\sum_{i=1}^d\beta_i$. Such counterexamples can be found in~\cite[Lemma 8]{mu21_renyi} (see also~\cite[Proposition 13]{farooq2024matrix}). In the previous section, we argued that it is sufficient to study the convexity properties of the function
\begin{equation}
    \vec{x} \longmapsto \|\vec{x}\|_1 \exp\left( t\int_{[0,+\infty]} H_\alpha(\hat{x}) \,\d\tau(\alpha) \right)\,.
\end{equation}
This function is connected to the one defined in equation~\ref{eq: multivariate divergences function} whenever the measure $\tau$ is discrete, under the identification
\begin{align}
    &\beta_i \rightarrow \tau(\{\alpha_i\})  t\frac{\alpha_i}{1-\alpha_i} \,, \qquad  v_i \rightarrow\left(\sum_j x_j^{\alpha_i}\right)^\frac{1}{\alpha_i} \,.
\end{align}
The key difference is that, while the variables $v_i$ can be chosen independently in the first case, the terms $\bigl(\sum_j x_j^{\alpha_i}\bigr)^{1/\alpha_i}$ are coupled: varying $\vec{x}$ simultaneously affects all of them. Nevertheless, one can control the variations of $\vec{x}$ so that they influence only a subset of these terms. This is achieved by constructing instances in very large dimensions. Concretely, consider appending a $d$-dimensional tail with weight $1/d$ to the vector $\vec{x}=(x_1,\ldots,x_k)$, i.e., consider the vector
\begin{equation}
   \vec{y} = (x_1,\ldots,x_k, 
   \underbrace{1/d, \ldots, 1/d}_{d\ \text{times}} ) \,.
\end{equation}
It is easy to verify that the additional tail affects only the terms with $\alpha_i<1$ (which diverge), while leaving those with $\alpha_i>1$ little affected. Indeed, we have that
\begin{equation}
    \left(\sum_j y_j^{\alpha_i}\right)^{\tfrac{1}{\alpha_i}} 
= \left(\sum_{j=1}^k x_j^{\alpha_i} + d^{\,1-\alpha_i}\right)^{\tfrac{1}{\alpha_i}}
\;\;\xrightarrow[d \to \infty]{}\;\;
\begin{cases}
+\infty, & \alpha_i < 1, \\[6pt]
\sum_{j=1}^k x_j^{\alpha_i}, & \alpha_i > 1.
\end{cases}
\end{equation}

Below, we establish a set of necessary conditions under the assumption that the measure satisfies the following property.
\begin{assumption}\label{assumption:finite integrals}
The positive/negative measure $\nu:\mc B\big([0,+\infty]\big)\to\mb R$ (typically $\nu=t\tau$) satisfies the following:
\begin{align}
\left|\int_{[0,1)}\nu(\alpha)\frac{\alpha}{1-\alpha}\right| <+\infty\,,\quad \left|\int_{(1,+\infty]}\nu(\alpha)\frac{\alpha}{1-\alpha}\right|<+\infty\,.
\end{align}
\end{assumption}
As mentioned earlier, this assumption requires the measure to vanish faster than $1/(1-\alpha)$ as $\alpha \to 1$, since otherwise the relevant integrals would diverge. In particular, this class of measures includes, as a special case, discrete measures with finite support, for which the corresponding integrals are finite.



\subsection{Necessary conditions; temperate case}
The temperate case corresponds to the bulk entropies $H_{t,\tau}$.
As discussed at the beginning of Section~\ref{sec: sufficient conditions measure}, it is sufficient to study the convexity properties of the function
\begin{equation}
    \vec{x} \longmapsto \|\vec{x}\|_1 \exp\left( t\int_{[0,+\infty]} H_\alpha(\hat{x}) \,\d\tau(\alpha) \right)
\end{equation}
with a probability measure $\tau$ and $t\in\R\setminus\{0\}$. For notational convenience, we will often combine the weight $t$ with the probability measure $\tau$ and treat their product as a positive/negative measure (i.e., a generally unnormalized finite measure that may take either positive or negative values depending on the sign of $t$). We denote this measure by $\mu$, defined as
\begin{align}
    \mu := t\,\tau \, .
\end{align}
We also define the coefficient
\begin{equation}
\label{eq: definition beta0}
     \beta_0 = 1 - \int_{[0,+\infty]\setminus\{1\}} \frac{\alpha}{1 - \alpha} \, \mathrm{d}\mu(\alpha) \,.
\end{equation}



Since the proof is rather involved, we divide the argument into several propositions.
We first start with the case where the measure $\mu$ is positive, or equivalently, the $\mathbb{R}_+$ case.
\begin{enumerate}
    \item In Proposition~\ref{prop: first prop measure}, we show that the measure cannot have support on values $\alpha > 1$, and the total weight on points with $\alpha < 1$ must be bounded.
\end{enumerate}
We then consider the case where the measure is negative, or equivalently, the $\mathbb{R}^{\rm opp}_+$ case.
\begin{enumerate}
    \item In Proposition~\ref{prop: prop no alpha=1}, we show that if the measure is supported at $\alpha=1$, then no support for $\alpha>1$ is allowed. This solves the case when the measure has support at $\alpha=1$.
    \item In Proposition~\ref{prop: second prop measure}, we show that if the measure vanishes at $\alpha=1$ and $\beta_0>0$, then it cannot have support on points with $\alpha>1$.
    \item In Proposition~\ref{prop: third prop measure}, we show that if the measure is zero at $\alpha=1$, and $\beta_0\leq 0$, then the measure cannot have support on more than one point $\alpha>1$.
\end{enumerate}
Together, these results show that the sufficient conditions derived in the previous section are tight under Assumption~\ref{assumption:finite integrals}. A more detailed discussion is given in Section~\ref{sec: towards complete characterization}.


\bigskip






We begin with the case where $t>0$, i.e.,\ $\mu>0$. In this case, we demonstrate that the measure cannot be supported on values $\alpha \geq 1$, and that the total weight on points with $\alpha \leq 1$ must be bounded.

\begin{proposition}
\label{prop: first prop measure}
Let $\mu:\mc B\big([0,+\infty]\big)\to\mb R$ be a positive measure that satisfies Assumption~\ref{assumption:finite integrals}. Then, the function
\begin{equation}
    \vec{x} \longmapsto \|\vec{x}\|_1 \exp\left( \int_{[0,+\infty]} H_\alpha(\hat{x}) \,\d\mu(\alpha) \right) \,,
\end{equation}
where $\hat{x} = \vec{x}/\|\vec{x}\|_1$,
    is concave only if $\;\textup{supp}(\mu) \subseteq [0,1]$, and $\;\displaystyle\int_{[0,1]} \frac{\alpha}{1 - \alpha} \, \mathrm{d}\mu(\alpha) \leq 1$.
\end{proposition}
\begin{proof}
The idea of the proof is to construct a point and a perturbation direction along which the second derivative of the function becomes strictly positive whenever the condition on the measure from the above proposition is violated. This contradicts concavity, since for a concave function, the second derivative must be non-positive in all directions. To achieve this, we explicitly construct a counterexample in increasingly high dimensions and show that such violations can indeed be found in the high-dimensional regime. Specifically, we first show that the measure cannot have support for $\alpha \geq 1$. Then, we establish that the measure’s weight on $\alpha < 1$ must be bounded. Both results are proved using the same perturbation. The assumption ensures that the resulting integrals in the analysis are all finite.

First, we introduce the vector $\vec{x} = \hat{x}\|\vec{x}\|_1$ and split the integral over the range $[0, +\infty]$ into several parts
\begin{align}
    &f(\vec{x})= \left( \sum_i x_i \right)^{\beta_0} \exp\Bigg( 
    \int_{[0,1)} \d\mu(\alpha)  H_\alpha(\vec{x}) \, 
    + \int_{(1, +\infty)} \d\mu(\alpha)  H_\alpha(\vec{x}) \, \notag\\
    &\quad\; \qquad \qquad \qquad\qquad\qquad\hspace{5em}
    + \mu(1) \sum_i H_1(\hat{x}) + \mu(+\infty) H_\infty(\hat{x})
    \Bigg) \,,
\end{align}
where $\beta_0 = 1- \int_{[0,1)}\d \mu(\alpha) \frac{\alpha}{1-\alpha} - \int_{(1,+\infty)}\d \mu(\alpha) \frac{\alpha}{1-\alpha}$.  Note that these integrals are well-defined because of Assumption~\ref{assumption:finite integrals}. Here, 
\begin{align}
    H_\alpha(\vec{x}) = \frac{1}{1-\alpha}\log{\sum_ix_i^\alpha} \,, \quad  H_1(\hat{x}) =-\sum_i\frac{x_i}{\sum_j x_j} \log \left( \frac{x_i}{\sum_j x_j}  \right) \,\quad  H_\infty(\hat{x}) =-\log \max_i \frac{x_i}{\sum_j x_j} \,.
\end{align}
To lighten our notations, we denoted $\mu(\alpha):=\mu(\{\alpha\})$ for $\alpha\in[0,+\infty]$.
For constant $p \in (0,1)$, we define the point $\vec{x}$ and the direction $\vec{v}$
\begin{align}
\vec{x}=\bigg[\underbrace{\frac{p}{d},\ldots,\frac{p}{d}}_{d\,\textrm{times}},\underbrace{\frac{1-p}{d^2},\ldots,\frac{1-p}{d^2}}_{d^2\,\textrm{times}}\bigg] \,,
\quad 
\vec{v}=c\bigg[\underbrace{-\frac{1}{d},\ldots,-\frac{1}{d}}_{d\,\textrm{times}},\underbrace{\frac{1}{d^2},\ldots,\frac{1}{d^2}}_{d^2\,\textrm{times}}\bigg]
\end{align}

Here, the constant $c>0$ ensures that $p - c >0$ as we send $p \rightarrow 0$. Indeed, we can rewrite \(\vec{x} + \lambda \vec{v} = (1 - \lambda) \vec{x} + \lambda \vec{v}'\), where \(\vec{v}' = \vec{x} + \vec{v}\). To ensure that this is a valid convex combination in the semiring, we require that \(\vec{v}' \geq 0\), i.e., \(\vec{v}'\) must be an admissible element of the semiring. Nevertheless, the factor 
$c$ only contributes an overall constant and does not affect the second derivative. Moreover, since we consider the limit of an infinitesimal perturbation, this constant is not needed and will therefore be omitted in the analysis.

By substituting the chosen point and direction, we have that
\begin{align}
    f(\vec{x}+\lambda\vec{v}) = \exp\Bigg(&\int_{[0,1)}\d \mu(\alpha)\frac{\alpha}{1-\alpha}\log{\left(d^{1-\alpha}(p-\lambda)^\alpha+d^{2(1-\alpha)}(1-p+\lambda)^\alpha\right)^\frac{1}{\alpha}} \notag \\
&+\int_{(1,+\infty)}\frac{\alpha}{1-\alpha}\log{\left(d^{1-\alpha}(p-\lambda)^\alpha+d^{2(1-\alpha)}(1-p+\lambda)^\alpha\right)^\frac{1}{\alpha}} \notag \\
    &\qquad\quad\qquad \qquad \;\; \qquad \qquad +\mu(1) h_3(\lambda) - \mu( +\infty) \log g_3(\lambda) \Bigg) \,.
\end{align}
Here, we defined the functions
\begin{align}
    &h_3(\lambda) = \left(-(p-\lambda) \log{\frac{p-\lambda}{d}}-(1-p+\lambda)\log{\frac{1-p+\lambda}{d^2}} \right) \,, \\
    & g_3(\lambda) = \max \left\{\frac{p-\lambda}{d},\frac{1-p+\lambda}{d^2}\right\} = \frac{p-\lambda}{d} \,,
\end{align}
where the second holds for large enough $d$.
The derivatives of the latter functions are
\begin{align}
    &h_3'(\lambda) =  \log{\frac{(p-\lambda)d}{1-p+\lambda}} \,, \quad h_3''(\lambda) = -\left(\frac{1}{p-\lambda}+\frac{1}{1-p+\lambda}\right)  \\
    &g_3'(\lambda) = -\frac{1}{d}\,, \quad g_3''(\lambda) =0
\end{align}
Moreover, we have the asymptotic behavior
\begin{align}
    h_3(0) \rightarrow (2-p) \log{d}\,, \; h_3'(0) \rightarrow  \;\log{d} \,, \; h_3''(0) \rightarrow -\left(\frac{1}{p}+\frac{1}{1-p}\right)  \,.
\end{align}
Furthermore, if we define the functions
\begin{align}
&g_1(\lambda)=\left(\varepsilon^{1-\alpha}(p-\lambda)^\alpha+(1-p+\lambda)^\alpha\right)^\frac{1}{\alpha}\\
    & g_2(\lambda)=\left((p-\lambda)^\alpha+\varepsilon^{\alpha-1}(1-p+\lambda)^\alpha\right)^\frac{1}{\alpha}
\end{align}
and we set $\varepsilon =1/d$. Note that, up to constant factors that do not affect the sign of the derivative, the function becomes
\begin{align}
    f(\vec{x}+\lambda\vec{v}) = \exp\Bigg(&\int_{[0,1)}\d \mu(\alpha)\frac{\alpha}{1-\alpha}\log{g_1(\lambda)} +\int_{(1,+\infty)}\frac{\alpha}{1-\alpha}\log{g_2(\lambda)} \notag \\
    &\quad \qquad \qquad \qquad \qquad +\mu( 1) h_3(\lambda) - \mu( +\infty) \log g_3(\lambda) \Bigg) \,.
\end{align}
Next, we compute the second derivative with respect to $\lambda$ and show that if the condition on the measure in the above proposition is not satisfied, the second derivative is positive. This implies the existence of a convex direction, contradicting concavity. We use the formulas for the second derivatives of the exponential and the logarithm of a function
\begin{align}
    &\frac{\d^2}{\d\lambda^2} \exp(h(\lambda)) = \exp(h(\lambda)) \left((h'(\lambda))^2+h''(\lambda)\right) \,, \\
&\frac{\d^2}{\d\lambda^2} \log(g(\lambda)) = \frac{g''(\lambda)}{g(\lambda)} - \left( \frac{g'(\lambda)}{g(\lambda)} \right)^2
\,.
\end{align}

A straightforward computation shows that the second derivative is
\begin{align}
    &\frac{\d^2}{\d \lambda^2}f(\vec{x}+\lambda\vec{v}) \\
    & = \exp\Bigg(\int_{[0,1)}\d \mu(\alpha)\frac{\alpha}{1-\alpha}\log{g_1(\lambda)}+\int_{(1,+\infty)}\d \mu(\alpha)\frac{\alpha}{1-\alpha}\log{g_2(\lambda)} \notag\\
    &\hspace{5em}-\mu( +\infty) \log g_3(\lambda) + \mu(1)h_3(\lambda) \Bigg) \notag \\
    & \quad  \times \Bigg[\left(\int_{[0,1)}\d \mu(\alpha)\frac{\alpha}{1-\alpha}\frac{g'_1(\lambda)}{g_1(\lambda)}+\int_{(1,+\infty)}\d \mu(\alpha)\frac{\alpha}{1-\alpha}\frac{g'_2(\lambda)}{g_2(\lambda)} -\mu( +\infty)\frac{g_3'(\lambda)}{g_3(\lambda)} + \mu(1)h_3'(\lambda)\right)^2 \notag\\
    &\;\;\qquad +\int_{[0,1)}\d \mu(\alpha)\frac{\alpha}{1-\alpha}\left(-\frac{g'_1(\lambda)^2}{g_1(\lambda)^2}+\frac{g''_1 (\lambda)}{g_1(\lambda)}\right)+ \int_{( 1,+\infty)}\d \mu(\alpha)\frac{\alpha}{1-\alpha}\left(-\frac{g'_2(\lambda)^2}{g_2(\lambda)^2}+\frac{g''_2 (\lambda)}{g_2(\lambda)}\right)\notag \\
    &  \hspace{19.2 em}- \mu(+\infty)\left(-\frac{g'_3(\lambda)^2}{g_3(\lambda)^2}+\frac{g''_3 (\lambda)}{g_3(\lambda)}\right) + \mu(1)h_3''(\lambda) \Bigg]\,.
\end{align}
Next, we evaluate the second derivative in $\lambda=0$ and consider the limit $\varepsilon \rightarrow 0$. We define the constants 
\begin{equation}
\beta_2=\int_{[0.1)} \d \mu(\alpha) \frac{\alpha}{1-\alpha} \,, \quad \beta_1=\int_{(1,+\infty)} \d\mu(\alpha) \frac{\alpha}{1-\alpha}  - \mu(+\infty)
\end{equation}
We obtain up to positive constants the asymptotic (small $\varepsilon$) behavior
\begin{align}
    &\left.\frac{\d^2 }{\d\lambda^2}f(\vec{x}+\lambda\vec{v})\right\vert_{\lambda=0}  \rightarrow \left[\left(\mu(1)\log{d}-\frac{\beta_1}{p}+\frac{\beta_2}{1-p}\right)^2-\frac{\beta_1}{p^2}-\frac{\beta_2}{(1-p)^2}- \mu(1)\left(\frac{1}{p}+\frac{1}{1-p}\right)\right] \,.
\end{align}
Note that we also invoked uniform convergence on the respective intervals, together with Assumption~\ref{assumption:finite integrals}, which is readily verified from the simple form of the functions inside the integrals.
We distinguish several cases. First, consider the case where $\mu(1) > 0$. In this setting, taking the limit $d \to +\infty$ yields an infinite and hence positive value, which provides the desired counterexample. Next, we consider the case $\mu(1) = 0$. To show that no points $\alpha > 1$ are allowed in this case, we expand the expression above to obtain
\begin{equation}
    \frac{\beta_1(\beta_1-1)}{p}+\frac{\beta_2(\beta_2-1)}{1-p}-2\frac{\beta_1\beta_2}{p(1-p)} \,.
\end{equation}
We then take the limit $p \rightarrow 0$  and obtain
\begin{equation}
    \left.\frac{\d^2 }{\d\lambda^2}f(\vec{x}+\lambda\vec{v})\right\vert_{\lambda=0} \rightarrow \frac{\beta_1(\beta_1-1)}{p} \,.
\end{equation}
Since $\beta_1 \leq 0$, the presence of any points $\alpha > 1$ with nonzero measure yields a direction along which the second derivative is positive, thereby providing the convexity needed for the counterexample. Finally, in the case $\beta_1 = 0$, we obtain
\begin{equation}
    \left.\frac{\d^2 }{\d\lambda^2}f(\vec{x}+\lambda\vec{v})\right\vert_{\lambda=0} \rightarrow \frac{\beta_2(\beta_2-1)}{(1-p)^2} 
\end{equation}
that shows that $\beta_2 \leq 1$, which is precisely the constraint on the measure in the proposition.
\end{proof}



%
Next, we consider the case where the measure is negative. 
We first prove that if $\mu(1) \neq 0$, then no support on the range $\alpha>1$ is allowed.



\begin{proposition}
\label{prop: prop no alpha=1}
Let $\mu:\mc B\big([0,+\infty]\big)\to\mb R$ be a negative measure such that $\mu(1) \neq 0$ and satisfies Assumption~\ref{assumption:finite integrals}.
Then, the function
\begin{equation}
    \vec{x} \longmapsto \|\vec{x}\|_1 \exp\left( \int_{[0,+\infty]} H_\alpha(\hat{x}) \,\d\mu(\alpha) \right)
\end{equation}
    is convex only if $\;\textup{supp}(\mu) \subseteq [0,1]$.
\end{proposition}

\begin{proof}
    \bigskip
We need to show that if $\mu(1) \neq 0$, then no support on the range $\alpha>1$ is allowed. Let us suppose that there is some support on $\alpha=1$ and $\alpha>1$ define
\begin{equation}
    \beta_0 = 1- \int_{[0,+\infty]\setminus\{1\}} \frac{\alpha}{1-\alpha}\d\mu(\alpha) \,,  \quad \beta_1 = \int_{(1,+\infty]} \frac{\alpha}{1-\alpha}\d\mu(\alpha) > 0 \,.
\end{equation}
Here, at the limit point $\alpha \to +\infty$, the function $\alpha/(1-\alpha)$ is taken to be equal to $1$, which is consistent with the definition of the Rényi entropy for $\alpha=+\infty$.
Note that the latter integrals are well-defined because of Assumption~\ref{assumption:finite integrals}.
We define the point $\vec{x}$ and the direction $\vec{v}$
\begin{align}
\vec{x}=\bigg[\underbrace{\frac{\beta_1}{d},\ldots,\frac{\beta_1}{d}}_{d\,\textrm{times}},\underbrace{\frac{-\mu(1)\log{d}}{d^2},\ldots,\frac{-\mu(1)\log{d}}{d^2}}_{d^2\,\textrm{times}}\bigg] \,,
\quad 
\vec{v}=\bigg[\underbrace{\frac{1}{d},\ldots,\frac{1}{d}}_{d\,\textrm{times}},\underbrace{0,\ldots,0}_{d^2\,\textrm{times}}\bigg] \,.
\end{align}
We recall that $\mu(1)<0$. By substituting the chosen point and direction, we have that
\begin{align}
&f(\vec{x}+\lambda\vec{v}) =\notag\\
&\qquad \exp\Bigg(\beta_0\log{(\beta_1+\lambda-\mu(1)\log{d})}\notag\\
&\qquad\qquad\; +\int_{[0,1)}\d \mu(\alpha)\frac{\alpha}{1-\alpha}\log{\left(d^{1-\alpha}(\beta_1+\lambda)^\alpha+d^{2(1-\alpha)}(-\mu(1)\log{d})^\alpha\right)^\frac{1}{\alpha}} \notag \\
&\quad\qquad\quad\;+\int_{(1,+\infty]}\frac{\alpha}{1-\alpha}\log{\left(d^{1-\alpha}(\beta_1+\lambda)^\alpha+d^{2(1-\alpha)}(-\mu(1)\log{d})^\alpha\right)^\frac{1}{\alpha}} +\mu(1) h_3(\lambda) \Bigg) \,.
\end{align}
Here, we defined the functions
\begin{align}
    h_3(\lambda) =& \left(-(\beta_1+\lambda) \log{\frac{\beta_1+\lambda}{d}}-(-\mu(1)\log{d})\log{\frac{-\mu(1)\log{d}}{d^2}} \right) \frac{1}{\beta_1+(-\mu(1)\log{d})} + \notag\\
    &\quad +\log{(\beta_1+(-\mu(1)\log{d}))}\,.
\end{align}
The derivatives of the latter function are
\begin{align}
    \hspace{-0.65em}h_3'(\lambda) = \frac{1}{\beta_1+(-\mu(1)\log{d})}\Bigg(\log{\Big(\frac{d}{\beta_1+\lambda}\Big)} -1\Bigg)\,, \,h_3''(\lambda) = \frac{c^2}{\beta_1+(-\mu(1)\log{d})}\Bigg(-\frac{1}{\beta_1+\lambda}\Bigg) 
\end{align}
Moreover, we have the asymptotic behavior
\begin{align}
     h_3'(0) \rightarrow  \frac{1}{\beta_1+(-\mu(1)\log{d})}\log{d} \,, \; h_3''(0) \rightarrow -\frac{1}{\beta_1} \frac{1}{\beta_1+(-\mu(1)\log{d})}\,.
\end{align}
Furthermore, let us set $\varepsilon=1/d$ and define the functions
\begin{align}
&g_0(\lambda)=(\log{(1/\varepsilon)})^{-1}(\beta_1+c\lambda)+(-\mu(1))\,,\\
&g_1(\lambda)=\left(\varepsilon^{1-\alpha}(\beta_1+c\lambda)^\alpha+(-\mu(1)\log{(1/\varepsilon)})^\alpha\right)^\frac{1}{\alpha}\,,\\
    & g_2(\lambda)=\left((\beta_1+c\lambda)^\alpha+\varepsilon^{\alpha-1}(-\mu(1)\log{(1/\varepsilon)})^\alpha\right)^\frac{1}{\alpha}\,.
\end{align}
Then, the function becomes up to positive constant
\begin{align}
    &f(\vec{x}+\lambda\vec{v}) \\
    &\quad =\exp\Bigg(\beta_0 \log{g_0(\lambda)}+\int_{[0,1)}\d \mu(\alpha)\frac{\alpha}{1-\alpha}\log{g_1(\lambda)} +\int_{(1,+\infty)}\frac{\alpha}{1-\alpha}\log{g_2(\lambda)} +\mu( 1) h_3(\lambda) \Bigg) \,.
\end{align}
Next, we compute the second derivative with respect to $\lambda$. 
We obtain up to the constants the asymptotic (large $d$) behavior
\begin{align}
    &\left.\frac{\d^2 }{\d\lambda^2}f(\vec{x}+\lambda\vec{v})\right\vert_{\lambda=0}  \rightarrow \left[\left( \frac{\mu(1)}{\beta_1+(-\mu(1)\log{d})}\log{d}+1\right)^2-\frac{1}{\beta_1}- \frac{\mu(1)}{\beta_1+(-\mu(1)\log{d})} \frac{1}{\beta_1} \right] \,.
\end{align}
In the limit of large $d$, the expression approaches $-1/\beta_1$, which is negative, indicating a concave direction and thus contradicting convexity.
\end{proof}

Next, we show that if $\beta_0$, defined in~\eqref{eq: definition beta0}, is positive and the measure is zero at $\alpha=1$, then the measure cannot have support on points with $\alpha > 1$.

\begin{proposition}
\label{prop: second prop measure}
Let $\mu:\mc B\big([0,+\infty]\big)\to\mb R$ be a negative measure such that $\mu(1)=0$, satisfies Assumption~\ref{assumption:finite integrals} and 
\begin{equation}
   \int_{[0,+\infty]\setminus\{1\}} \frac{\alpha}{1-\alpha} \,\d\mu(\alpha) < 1 \,.
\end{equation}
Then, the function
\begin{equation}
    \vec{x} \longmapsto \|\vec{x}\|_1 \exp\left( \int_{[0,+\infty]} H_\alpha(\hat{x}) \,\d\mu(\alpha) \right)
\end{equation}
    is convex only if $\;\textup{supp}(\mu) \subseteq [0,1]$.
\end{proposition}
\begin{proof}
Similarly to the previous case, we identify a direction along which the second derivative is negative whenever the condition on the measure is violated—that is, if there exists a point in the support of the measure with $\alpha > 1$. This corresponds to a concave direction, contradicting the requirement of convexity. 

We denote the weight on the $\alpha\in(1,+\infty)$ range and the power of the sum of the components of the vector as
\begin{equation}
    \beta_1=\int_{(1,+\infty)} \d \mu(\alpha) \frac{\alpha}{1-\alpha} \,, \quad \beta_0=1-\int_{[0,+\infty] \setminus\{1\}} \d \mu(\alpha) \frac{\alpha}{1-\alpha} \,.
\end{equation}
Note that these integrals are well-defined because of Assumption~\ref{assumption:finite integrals}.

Note that because the measure is negative and $\alpha > 1$, it follows that $\beta_1 > 0$. Moreover, by assumption $\beta_0 > 0$.
Next, we define the constant $\varsigma= (\beta_1+2)/\beta_0$ as well as the point and direction
\begin{align}
\vec{x}=\bigg[\beta_1,\underbrace{\frac{1}{d},\ldots,\frac{1}{d}}_{d\,\textrm{times}},\underbrace{\frac{1}{d^2},\ldots,\frac{1}{d^2}}_{d^2\,\textrm{times}}\bigg] \,,
\quad 
\vec{v}=\bigg[-1,\underbrace{\frac{\varsigma}{d},\ldots,\frac{\varsigma}{d}}_{d\,\textrm{times}},\underbrace{0,\ldots,0}_{d^2\,\textrm{times}}\bigg] \,.
\end{align}
Note that the direction is not normalized. Indeed, we have
\begin{equation}
    \sum_i x_i+ \lambda v_i = \beta_1+2-\lambda(\varsigma-1) = \frac{\beta_1+2}{\beta_0}(\beta_0+\lambda) \,.
\end{equation}
We have up to the positive and finite constants
that the function becomes
\begin{align}
    &f(\vec{x} + \lambda v_i)\\
    &\quad= (\beta_0+ \lambda)^{\beta_0}\exp{\left(\int_{[0,1)}\d \mu(\alpha)\frac{\alpha}{1-\alpha}\log{g_1(\lambda)}+\int_{(1,+\infty]}\d \mu(\alpha)\frac{\alpha}{1-\alpha}\log{g_2(\lambda)}\right)}\\
    &\quad=\exp{\left(\beta_0 \log{g_0(\lambda)}+\int_{[0,1)}\d \mu(\alpha)\frac{\alpha}{1-\alpha}\log{g_1(\lambda)}+\int_{(1,+\infty]}\d \mu(\alpha)\frac{\alpha}{1-\alpha}\log{g_2(\lambda)}\right)} \,,
\end{align}
where we defined $\varepsilon=1/d$ and the functions
\begin{align}
    &g_0(\lambda) = \beta_0+c\lambda \,,\\
    & g_1(\lambda) = (\varepsilon^{2(1-\alpha)}(\beta_1-c\lambda)^\alpha+\varepsilon^{1-\alpha}(1-c\varsigma\lambda)^\alpha+1)^\frac{1}{\alpha} \,, \\
    & g_2(\lambda) = ((\beta_1-c\lambda)^\alpha+\varepsilon^{\alpha-1}(1-c\varsigma\lambda)^\alpha+\varepsilon^{2(\alpha-1)})^\frac{1}{\alpha} \,.
\end{align}

The second derivative yields
\begin{align}
    &\frac{\d^2}{\d \lambda^2}f(\vec{x}+\lambda\vec{v})\notag\\
    &\; = \exp{\left(\beta_0 \log{g_0(\lambda)}+\int_{[0,1)}\d \mu(\alpha)\frac{\alpha}{1-\alpha}\log{g_1(\lambda)}+\int_{(1,+\infty]}\d \mu(\alpha)\frac{\alpha}{1-\alpha}\log{g_2(\lambda)}\right)}\notag\\
    &\; \; \quad \quad\qquad  \times \Bigg[\left(\beta_0\frac{g_0(\lambda)'}{g_0(\lambda)}+\int_{[0,1)}\d \mu(\alpha)\frac{\alpha}{1-\alpha}\frac{g'_1(\lambda)}{g_1(\lambda)}+\int_{\alpha >1}\d \mu(\alpha)\frac{\alpha}{1-\alpha}\frac{g'_2(\lambda)}{g_2(\lambda)}\right)^2\notag\\
    &\hspace{7em} +\beta_0\left(-\frac{g'_0(\lambda)^2}{g_0(\lambda)^2}+\frac{g''_0 (\lambda)}{g_0(\lambda)}\right) +\int_{[0,1)}\d \mu(\alpha)\frac{\alpha}{1-\alpha}\left(-\frac{g'_1(\lambda)^2}{g_1(\lambda)^2}+\frac{g''_1 (\lambda)}{g_1(\lambda)}\right)\notag\\
    & \hspace{17em} + \int_{(1,+\infty]}\d \mu(\alpha)\frac{\alpha}{1-\alpha}\left(-\frac{g'_2(\lambda)^2}{g_2(\lambda)^2}+\frac{g''_2 (\lambda)}{g_2(\lambda)}\right)\Bigg] \,.
\end{align}
Taking $\lambda = 0$ and letting $\varepsilon \to 0$, we obtain the following asymptotic behavior as $\varepsilon \to 0$
\begin{equation}
    \left.\frac{\d^2 }{\d\lambda^2}f(\vec{x}+\lambda\vec{v})\right\vert_{\lambda=0} \rightarrow \beta_0^{\beta_0} \beta_1^{\beta_1} \left[\left(1-1\right)^2-\frac{1}{\beta_0}-\frac{1}{\beta_1}\right] \,.
\end{equation}
Expanding and dividing by the positive constant $\beta_0^{\beta_0} \beta_1^{\beta_1}$ we obtain
\begin{equation}
    \left.\frac{\d^2 }{\d\lambda^2}f(\vec{x}+\lambda\vec{v})\right\vert_{\lambda=0} \rightarrow \frac{-(\beta_0+\beta_1)}{\beta_0\beta_1} \,,
\end{equation}
which is negative since $\beta_0, \beta_1 > 0$. This contradicts convexity, and thus we must have $\beta_1 = 0$, meaning that no points with $\alpha > 1$ are allowed.
\end{proof}

\begin{remark}
In the proof of the proposition above, we considered a point $\vec{x}$ that is not normalized, meaning $\sum_i x_i \neq 1$. This is justified because the function is defined on the entire domain $\vec{x} \in \mathbb{R}_+^{d'}$ for some dimension $d'$. Indeed, the semiring includes vectors that are not normalized, and within this broader space, counterexamples can still arise. By semiring theory, the existence of such counterexamples in this extended setting is sufficient to exclude the corresponding measure values in the barycentric decomposition of conditional entropies stated in Theorem~\ref{th:barycenter conditional entropies}, even though the theorem itself concerns normalized probability distributions.
This conclusion relies on the fact that, without loss of generality, we may assume $\|\vec{x}\|_1 = 1$, since any vector $\vec{x}$ can be rescaled by a positive constant, and such scaling does not affect convexity properties.
\end{remark}

\begin{remark}
In the proof of the above proposition, we used a nontraceless direction, i.e. a direction $\vec{v}$ such that $\sum_i v_i \neq 0$.
 This is admissible because the function is defined on the whole domain $\vec{x} \in \mathbb{R}_+^{d'}$ for some dimension $d'$.  In particular, the semiring is defined over the set of unnormalized joint matrices, within which counterexamples can be found. By semiring theory, the existence of such counterexamples in this broader space is sufficient to rule out the corresponding measure values in the barycentric decomposition of conditional entropies in Theorem~\ref{th:barycenter conditional entropies}, even though the theorem itself is stated for normalized probability distributions.

This fact follows from the additivity property of homomorphisms, as we now explain. First, without loss of generality, we may assume that $\|\vec{x}\|_1 < 1$; indeed, $\vec{x}$ can always be rescaled by a positive constant, and such a multiplicative factor does not affect convexity properties.
Next, consider a perturbation of $\vec{x}$ in a direction $\vec{v}$ with $\sum_i v_i \neq 0$, i.e., a direction that is not traceless. To handle this, we embed the problem into a larger space by introducing an additional outcome to the conditioning register $Y$. Instead of a single probability vector $\vec{x}$, we now consider a matrix with two columns, $\vec{x}$ and $\vec{x}'$, corresponding to a joint distribution over $X \times Y$.
In this extended space, by the additivity of the homomorphism $f$, the total function becomes $f(\vec{x}) + f(\vec{x}')$. Now, we perturb $\vec{x}$ as $\vec{x} + \lambda \vec{v}$, and simultaneously perturb $\vec{x}'$ as $\vec{x}' + \lambda \hat{x}' \sum_i v_i$. This ensures that the total weight of the joint system remains constant. Since $f$ depends linearly on the normalization of $\vec{x}'$, this adjustment does not affect any convexity analysis. Therefore, this construction allows us to effectively study perturbations of $\vec{x}$ in directions $\vec{v}$ with $\sum_i v_i \neq 0$, by embedding them into a trace-preserving framework.

\end{remark}
Finally, in the case where the measure is negative, we show that if $\beta_0$, defined in~\eqref{eq: definition beta0},
is negative and the measure is zero at $\alpha=1$, then the measure cannot have support on more than one point $\alpha > 1$.
\begin{proposition}
\label{prop: third prop measure}
Let $\mu:\mc B\big([0,+\infty]\big)\to\mb R$ be a negative measure such that $\mu(1)=0$, satisfies Assumption~\ref{assumption:finite integrals} and
\begin{equation}
\label{eq: cond on beta0}
   \int_{[0,+\infty]\setminus\{1\}} \frac{\alpha}{1-\alpha} \,\d\mu(\alpha) \geq 1 \,.
\end{equation}
Then, the function
\begin{equation}
    \vec{x} \longmapsto \|\vec{x}\|_1 \exp\left( \int_{[0,+\infty]} H_\alpha(\hat{x}) \,\d\mu(\alpha) \right)
\end{equation}
    is convex only if $\;\operatorname{supp}(\mu) \subseteq [0,1] \cup \{\alpha_*\}$, $\alpha_* \in [1, +\infty]$.
\end{proposition}

\begin{proof}
We need to show that the measure can have support on at most one point $\alpha > 1$. 
In particular, we then construct a point and a direction along which the function is concave by computing the second derivative and showing that it is negative. This provides a counterexample to convexity.

Let us assume that the measure has support on two different points $\alpha_2 > \alpha_1 > 1$. We define the points
\begin{equation}
    x_2 = 1 + \frac{\alpha_1 - 1}{2}, \quad x_3 = \alpha_1 + \frac{\alpha_2 - \alpha_1}{2}, \quad x_4 = \alpha_2 - \delta.
\end{equation}
for some fixed $\delta>0$.
Then, $(x_2,x_3)$ is a neighborhood  of $\alpha_1$ and a value $x_4<\alpha_2$. By our assumption on the support of the measure and recalling that the measure $\mu$ is negative, it follows that
\begin{equation}
    \beta_1 = \int_{(x_2,x_3)} \frac{\alpha}{1-\alpha}\d\mu(\alpha) > 0 \,, \text{and} \quad \beta_2=  \int_{\alpha > x_4} \d\mu(\alpha) \frac{\alpha}{1-\alpha} > 0 \,.
\end{equation} 
We depict these points in Fig~\ref{fig:patch}. 
\begin{figure}[h]
\hspace*{0.15\textwidth}
\includegraphics[width=.8\textwidth]{patches4.pdf}
    \caption{The figure shows the lines defined in equation~\eqref{eq: lines necessary}, and the dominant contributions in equation~\eqref{eq: contributions lines} in the different ranges.}
    \label{fig:patch}
\end{figure}
Note that these integrals are always finite since the pole at $\alpha=1$ of the function $\alpha/(1-\alpha)$ is excluded and the total weight of the measure is finite.
The idea is to define a point and a direction so that, for large dimension values, the relevant function behaves as $\beta_1 + \lambda$ on the interval $\alpha \in (x_2, x_3)$ and as $\beta_2 - \lambda$ on the interval $\alpha > x_4$. We illustrate this situation in Fig.~\ref{fig:patch}. We then define the lines
\begin{equation}
\label{eq: lines necessary}
    y_1(\alpha) = \alpha, \quad y_2(\alpha) = (\alpha - 1) s_2, \quad y_3(\alpha) = (\alpha - \alpha_1) s_3, \quad y_4(\alpha) = (\alpha - \alpha_{12}) s_4,
\end{equation}
where we define $\alpha_{12}=(\alpha_1+\alpha_2)/2$. The slopes $s_i$ satisfy the recursive relation
\begin{equation}
    s_i = \frac{x_i - b_{i-1}}{x_i - b_i} s_{i-1}, 
\end{equation}
where $b_i$ are defined as
\begin{equation}
    b_1 = 0, \quad b_2 = 1, \quad b_3 = \alpha_1, \quad b_4 = \alpha_{12}.
\end{equation}
Hence, each line is of the form $y = (\alpha - b_i) s_i$ where $b_i$ specifies the intersection with the $x$ axis and $s_i$ determines their slope. These lines are depicted in Fig.~\ref{fig:patch}.
Using these definitions, we choose the point and direction
\begin{align}
&\vec{x}=\bigg[\underbrace{d^{s_1} \delta_1,\ldots,d^{s_1} \delta_1}_{\lceil d^{b_4 s_4}\rceil },\underbrace{d^{s_2} \beta_1,\ldots,d^{s_2} \beta_1}_{\lceil d^{b_4 s_4 -b_2 s_2}\rceil},\underbrace{d^{s_3} \delta_2,\ldots,d^{s_3} \delta_2}_{\lceil d^{b_4 s_4 -b_3 s_3}\rceil},d^{s_4} \beta_2\bigg] \,,
\\
&\vec{v}=\bigg[\underbrace{0,\ldots,0}_{\lceil d^{b_4 s_4}\rceil },\underbrace{d^{s_2} ,\ldots,d^{s_2}}_{\lceil d^{b_4 s_4 -b_2 s_2}\rceil},\underbrace{0,\ldots,0}_{\lceil d^{b_4 s_4 -b_3 s_3}\rceil},-d^{s_4} \bigg]\,,
\end{align}
where $\delta_1,\delta_2$ are two fixed positive constants. Here, $\lceil \cdot \rceil$ denotes the ceiling function, ensuring that the number of entries is an integer. We have that the inner function is
\begin{equation}
\label{eq: contributions lines}
    \sum_i(x_i+\lambda v_i)^\alpha = \lceil d^{b_4 s_4}\rceil (d^{s_1}\delta_1)^\alpha  + \lceil d^{b_4 s_4 -b_2 s_2}\rceil (d^{s_2}\beta_1+\lambda d^{s_2})^\alpha + \lceil d^{b_4 s_4 -b_3 s_3}\rceil (d^{s_3}\delta_2)^\alpha + (d^{s_4}\beta_2-\lambda d^{s_4})^\alpha
\end{equation}
that multiplied by $d^{-b_4 s_4}$ gives, up to an irrelevant constant (due to the ceiling) that vanishes in the limit $d$ to infinity, 
\begin{equation}
    \sum_i(x_i+\lambda v_i)^\alpha  =  d^{y_1(\alpha)}\delta_1^\alpha  +  d^{y_2(\alpha)} (\beta_1+\lambda )^\alpha +  d^{y_3(\alpha)} \delta_2^\alpha + d^{y_4(\alpha)}(\beta_2-\lambda)^\alpha \,.
\end{equation}
In particular, for each range, there will be only one function $y(\alpha)$ that is dominant in the limit $d\rightarrow \infty$. Defining $\varepsilon = 1/d$,  the relevant inner functions are, up to terms that disappear as $\varepsilon \rightarrow 0$,
\begin{align}
    &g_1(\lambda) \approx \delta_1^\alpha \,,\;\;g_2(\lambda) \approx (\beta_1+\lambda)^\alpha  \,,\;\;g_3(\lambda) \approx \delta_2^\alpha  \,,\;\;g_4(\lambda) \approx (\beta_2-\lambda)^\alpha  \,.
\end{align}
We then split the integral into several ranges according to Fig.~\ref{fig:patch} and obtain up to positive constant factors
\begin{align}
    &f(\vec{x} + \lambda v_i) \\
    &\; = \exp\Bigg(\beta_0 \log{g_1(\lambda)}+\int_{[0,x_2)}\d \mu(\alpha)\frac{\alpha}{1-\alpha}\log{g_1(\lambda)}+ \int_{[x_2,x_3]}\d \mu(\alpha)\frac{\alpha}{1-\alpha}\log{g_2(\lambda)}\\
    &\qquad \qquad \qquad \qquad \quad \;\;+ \int_{ (x_3,x_4)}\d \mu(\alpha)\frac{\alpha}{1-\alpha}\log{g_3(\lambda)} + \int_{\alpha \geq x_4 }\d \mu(\alpha)\frac{\alpha}{1-\alpha}\log{g_4(\lambda)}\Bigg)
\end{align}
Here, $\beta_0$ is a finite constant due to Assumption~\ref{assumption:finite integrals}.
We then repeat the calculations from the previous examples and, up to constant factors and in the limit $\varepsilon \to 0$, we obtain the following second derivative
\begin{equation}
    \left.\frac{\d^2 }{\d\lambda^2}f(\vec{x}+\lambda\vec{v})\right\vert_{\lambda=0} \rightarrow \frac{-(\beta_1+\beta_2)}{\beta_1\beta_2} \,.
\end{equation}
Since both $\beta_1$ and $\beta_2$ are positive, the second derivative is negative, yielding a counterexample to convexity.
This shows that at most one point with $\alpha > 1$ is allowed, thereby concluding the proof.
\end{proof}



\subsection{Necessary conditions for derivations}
\label{sec: necessary derivations}
Next, we turn to the necessary conditions for the (extremal) derivations.
These are defined as (see equation~\eqref{eq:ExtDerivation})
 \begin{equation}
    H_{0,\alpha}(P_{XY}) = \sum_{y \in \Y} P_Y(y)H_\alpha(P_{X|Y=y}) \,.
 \end{equation} 
 Hence, it is sufficient to consider the convexity properties of the function
\begin{equation}
      \vec{x} \longmapsto \|\vec{x}\|_1 H_\alpha(\hat{x}) \,.
 \end{equation}
Next, we show that in order for the derivations to be monotone, it must be that
$\alpha>1$. we note that, by Lemma~\ref{lem: conv derivations}, monotonicity of derivations requires concavity.
\begin{proposition}
 \label{prop: prop measure deriv}
The function
 \begin{equation}
     \vec{x} \longmapsto \|\vec{x}\|_1  H_\alpha(\hat{x}) 
\end{equation}
     is concave only if $\alpha \in [0,1]$.
 \end{proposition}
 \begin{proof}
 We must show that the measure cannot have support on $\alpha > 1$. By Lemma~\ref{lem: conv derivations}, the derivations must be concave. Thus, we need to show that if the measure has support on $\alpha > 1$, we can construct a convex direction, yielding a counterexample.  

 The point $\vec{x}$ and direction $\vec{v}$ are the same as in Proposition~\ref{prop: first prop measure}. We define for $p \in (0,1)$
\begin{align}
 \vec{x}=\bigg[\underbrace{\frac{p}{d},\ldots,\frac{p}{d}}_{d\,\textrm{times}},\underbrace{\frac{1-p}{d^2},\ldots,\frac{1-p}{d^2}}_{d^2\,\textrm{times}}\bigg] \,,
 \quad 
\vec{v}=\bigg[\underbrace{-\frac{1}{d},\ldots,-\frac{1}{d}}_{d\,\textrm{times}},\underbrace{\frac{1}{d^2},\ldots,\frac{1}{d^2}}_{d^2\,\textrm{times}}\bigg]
 \end{align}
 We have that
 \begin{align}
     \sum_i (x_i+\lambda v_i)^\alpha = d^{1-\alpha}(p-\lambda)^\alpha+d^{2(1-\alpha)}(1-p+\lambda)^\alpha
 \end{align}
 Therefore, the function becomes
 \begin{align}
   &f(\vec{x}+\lambda \vec{v}) =  \frac{\alpha}{1-\alpha}\log{\left(d^{1-\alpha}(p-\lambda)^\alpha+d^{2(1-\alpha)}(1-p+\lambda)^\alpha\right)^\frac{1}{\alpha}} \,,
 \end{align}
 Furthermore, if we define the function
 \begin{align} g(\lambda)=\left((p-\lambda)^\alpha+\varepsilon^{\alpha-1}(1-p+\lambda)^\alpha\right)^\frac{1}{\alpha}
 \end{align}
 and we set $\varepsilon =1/d$, we obtain up to the additive constant factors (that do not affect the derivatives since they are constant)
 that the function becomes
 \begin{align}
 f(\vec{x}+\lambda\vec{v}) = &\frac{\alpha}{1-\alpha}\log{g(\lambda)} \,.
 \end{align}
 Next, we compute the second derivative with respect to $\lambda$ and show that if the condition on the measure in the above proposition is violated, this derivative becomes positive. This yields a convex direction, contradicting concavity. We use the identity for the derivative of the logarithm of a function, namely
 \begin{equation}
 \frac{\d^2}{\d\lambda^2} \log(g(\lambda)) = \frac{g''(\lambda)}{g(\lambda)} - \left( \frac{g'(\lambda)}{g(\lambda)} \right)^2 \,.
 \end{equation}
 The second derivative is
 \begin{align}
     &\frac{\d^2}{\d \lambda^2}f(\vec{x}+\lambda\vec{v}) =\frac{\alpha}{1-\alpha}\left(-\frac{g'(\lambda)^2}{g(\lambda)^2}+\frac{g'' (\lambda)}{g(\lambda)}\right) \,.
 \end{align}
 By taking $\lambda=0$ and the limit $\varepsilon \rightarrow 0$, and defining 
 \begin{equation}
 \beta= \frac{\alpha}{1-\alpha}  \,,
 \end{equation}
we obtain the large $d$ behavior
 \begin{align}
     &\left.\frac{\d^2 }{\d\lambda^2}f(\vec{x}+\lambda\vec{v})\right\vert_{\lambda=0}  \rightarrow -\frac{\beta}{p^2} \,.
 \end{align}
 which, since $\beta \leq 0$, is positive whenever $\beta_1 \neq 0$. This condition is equivalent to the measure having support on $\alpha > 1$, thus yielding the desired convex direction. The case $\alpha=+\infty$ is done in the same way. 
 \end{proof}

We now obtain the following result as a direct corollary of the above result combined with the sufficiency conditions of Proposition \ref{prop: sufficient derivations} and Lemma \ref{lem: conv derivations} in the appendix.

\begin{proposition}\label{prop:Derivations}
The set of derivations of the conditional majorization semiring $S$ at $\|\cdot\|$ is spanned by the quantities $H_{0,\alpha}$ of Definition \ref{def: extremal entropies} with $\alpha\in[0,1]$.
\end{proposition}

\section{Large sample and catalytic conditional majorization}
\label{sec: large sample}
In this section, we present the large-sample and catalytic results, which serve as the main tools for proving the general form of a conditional entropy in Theorem~\ref{th:barycenter conditional entropies} and will also underpin the applications discussed in Section~\ref{sec: Applications}.
In particular, these results follow directly from Theorem~\ref{thm:Fritz2022} on large-sample and catalytic transformations for general semirings, applied to the semiring of conditional majorization, together with the characterization of monotone homomorphisms and derivations established in Section~\ref{sec: semiring of conditional entropies}. 
As shown in the following section (see Theorem~\ref{th:barycenter conditional entropies}), these homomorphisms and derivations coincide, up to a constant and a logarithmic factor, with the extremal conditional entropies that generate the entire family of conditional entropies through convex combinations.
Since conditional entropies, according to Definition~\ref{def: conditional entropy}, require proper normalization and additivity under tensor products, we reformulate the monotone homomorphisms from Proposition~\ref{prop: monotone homorphisms} and the derivations from Proposition~\ref{prop: monotone derivations} as follows. 



\begin{definition}
\label{def: extremal entropies}
Let $P_{XY}$ be a probability distribution, $t\in\R\setminus\{0\}$, and $\tau:\mc B\big([0,+\infty]\big)\to\mb R$ be a probability measure. Then, we define
\begin{align}
    &H_{t,\tau}(X|Y)_P = \frac{1}{t}\log \left(\sum_{y \in \Y} P_Y(y)\exp{\left(t\int_{[0,+\infty]}H_\alpha(P_{X|Y=y})\d\tau(\alpha)\right)}\right) \,, \label{eq:TempEnt}\\
    & H_{-\infty,\tau}(X|Y)_P =  \min \limits_{\substack{y \in \mathcal{Y} \\ P_Y( y) > 0}}\int_{[0,+\infty]}H_\alpha(P_{X|Y=y})\d\tau(\alpha) \,, \label{eq:TropEnt}\\
    & H_{0,\alpha}(X|Y)_P =  \sum_{y \in \Y} P_Y(y)H_\alpha(P_{X|Y=y})\,.,\label{eq:DerivEnt} \\
    & H_{+\infty,0}(X|Y)_P =  \max \limits_{\substack{y \in \mathcal{Y} \\ P_Y( y) > 0}} H_0(P_{X|Y=y}) \,. \label{eq:TropEntno}
\end{align}
\end{definition}
Throughout this manuscript, we refer to the conditional entropy \(H_{t,\tau}\) as the bulk conditional entropies, as they are defined for finite values of the parameter \(t\). This distinguishes them from the entropies \(H_{-\infty,\tau}\), \(H_{+\infty,\tau}\), and \(H_{0,\alpha}\), which arise as limiting cases of $H_{t,\tau}$ corresponding to \(t \to -\infty\), \(t \to +\infty\), and \(t \to 0\), respectively.



As we prove later in Theorem~\ref{th:barycenter conditional entropies}, the above quantities encompass all (extremal) conditional entropies for some values of the parameters $t$ and $\tau$. While all extremal entropies necessarily take this form, as is discussed in more detail in Sections~\ref{sec: sufficient conditions measure},~\ref{sec: necessary conditions measures}, and~\ref{sec: set of conditional entropies}, not every choice of the parameters \(t\) and \(\tau\) is admissible. Indeed, for certain parameter choices the resulting conditional entropy may fail to satisfy the axioms of Definition~\ref{def: conditional entropy}, and in particular may violate monotonicity under conditionally mixing channels. See Section \ref{sec: towards complete characterization} where we summarize our current knowledge on the proper parameter sets. Note that, according to Proposition \ref{prop:Derivations}, we already know that $\alpha\in[0,1]$ is the correct parameter region for the quantities $H_{0,\alpha}$ of \eqref{eq:DerivEnt}.
Consequently, we introduce the following sets of parameters and measures.
\begin{definition}
    \label{def: true sets} Let $\D_{\rm bulk}$ denote the set of pairs $(t,\tau)$, with $t\in\mathbb{R}\setminus\{0\}$ and $\tau:\mathcal{B}([0,+\infty])\to[0,1]$ a probability measure, such that $H_{t,\tau}$ in \eqref{eq:TempEnt} is monotone under conditionally mixing channels. Similarly, let $\D_{-\infty}$ denote the set of measures $\tau$ for which $H_{-\infty,\tau}$ in \eqref{eq:TropEnt} is monotone under conditionally mixing channels.
\end{definition}

Next, we present the large-sample and catalytic results for transforming joint probability distributions under conditionally mixing operations, relabeling, and embedding, which jointly define the conditional majorization order.

\begin{theorem}
\label{th: large sample and catalytic}
    Let $P_{XY}$ and $Q_{X'Y'}$ be joint probability distributions. If
    \begin{align}
    & H_{t,\tau}(X|Y)_P < H_{t,\tau}(X'|Y')_Q, \qquad\qquad\qquad   \forall (t,\tau) \in \D_{\rm{bulk}},\label{eq:TempCond} \\
    &H_{-\infty,\tau}(X|Y)_P < H_{-\infty,\tau}(X'|Y')_Q, \; \quad\qquad \;\;\, \forall \tau\in \D_{-\infty},\label{eq:TropCond} \\
    &H_{0,\alpha}(X|Y)_P<H_{0,\alpha}(X'|Y')_Q,  \quad\,\quad\qquad\quad\;\; \forall \alpha\in[0,1]\,,\label{eq:DerivCond} \\
    &H_{+\infty,0}(X|Y)_P < H_{+\infty,0}(X'|Y')_Q\,,
    \end{align}
    then
   \begin{enumerate}[itemsep=0pt]
    \item for sufficiently large \(n \in \mathbb{N}\), it holds that $P_{XY}^{\otimes n} \succeq Q_{X'Y'}^{\otimes n}$.
    \item there exists a joint probability distribution \(R_{X''Y''}\) such that
    \begin{equation}
        P_{XY} \otimes R_{X''Y''} \succeq Q_{X'Y'} \otimes R_{X''Y''}.
    \end{equation}
\end{enumerate}
\end{theorem}
\begin{proof}
The result follows directly from an application of the general large-sample theorem for preordered semirings stated in Theorem~\ref{thm:Fritz2022} along with the specific characterization of monotone homomorphisms and derivations for the conditional majorization semiring in Proposition~\ref{prop: monotone homorphisms} and~\ref{prop: monotone derivations}. Indeed, the entropies appearing in the assumptions coincide with the homomorphisms and derivations up to additive constants and logarithmic factors, which preserve the monotonicity properties. Especially for the derivations, note that, according to Proposition \ref{prop:Derivations}, all derivations at $\|\cdot\|$ are spanned by $H_{0,\alpha}$, $\alpha\in[0,1]$, so having strict inequalities over all these derivations is equivalent to having the strict inequalities of \eqref{eq:DerivCond}. Hence, by appropriately accounting for the direction of the inequalities, the assumptions of the theorems coincide with those required in Theorem~\ref{thm:Fritz2022} on the ordering of homomorphisms and derivations.

In addition, according to Theorem~\ref{thm:Fritz2022}, we need to verify that the input distribution $Q_{X'Y'}$ is power universal. We note that here that \(Q_{X'Y'}\) is the input distribution, rather than \(P_{XY}\), since the preorder of the semiring \(\rgeq\) is reversed relative to the conditional majorization \(\succeq\) according to the definition of the preorder of the conditional majorization semiring in item~\ref{item: preorder} of Section~\ref{sec: semiring of conditional entropies}.
The fact that $Q_{X'Y'}$ is power universal is already guaranteed by the above assumptions on the strict ordering of entropies. Indeed, if $Q_{X'Y'}$ were not power universal, according to Proposition~\ref{prop: power universals}, there would exist a column of $Q_{X'Y'}$ with exactly one nonzero entry. 
In this case, the Hayashi conditional entropy, which is an element of the set \(\mathcal{D}_{-\infty}\)—a well-known fact in the literature (see, e.g., \cite{gour2024inevitability})— would be zero for $Q_{X'Y'}$. Indeed,  
\begin{equation}
    H^{\rm{H}}_{\infty}(X'|Y')_Q = \min_{y \in \mathcal{Y}'}  H_{\infty}\bigl(Q_{X'|Y'=y'}\bigr) \, = 0.
\end{equation}
Since $ H^{\rm{H}}_{\infty}(X|Y)_P\geq0$, the strict inequalities assumed in the theorem imply that $Q_{X'Y'}$ must be power universal.
\end{proof}

 \section{The set of conditional entropies}
\label{sec: set of conditional entropies}
In this section, we present the main result of our work, namely the general characterization of conditional entropies in Theorem \ref{th:barycenter conditional entropies}. Specifically, we show that any conditional entropy defined as a function satisfying the axioms in Definition~\ref{def: conditional entropy} can be expressed as a (generalized) convex combination of the entropies $H_{t,\tau}$ and its limiting forms $H_{0,\alpha}$, $H_{-\infty,\tau}$ and $H_{+\infty,0}$ defined in~\ref{def: extremal entropies}. In this sense, the conditional entropies $H_{t,\tau},H_{0,\alpha},H_{-\infty,\tau}$ and $H_{+\infty,0}$ constitute the extremal points of the set of the conditional entropies; for the extremality, see Theorem 14 of \cite{haapasalo_inprep}. As discussed in Section~\ref{sec: necessary conditions measures}, these quantities fail to be monotone under conditionally mixing channels for some parameter values. Consequently, they do not qualify as conditional entropies—according to Definition~\ref{def: conditional entropy} for arbitrary choices of the parameters. Our description of all conditional entropies has to be stated in terms of the abstract parameter sets $\mathcal{D}_{\mathrm{bulk}}$, $\mathcal{D}_{-\infty}$, and $\mathcal{D}_{0}$ whose exact identity is still unknown. However, we discuss substantial progress towards their full characterization in Section~\ref{sec: towards complete characterization}. In Section \ref{sec: literature}, we review more in depth how the conditional entropies discussed in the literature and also mentioned in the Introduction relate to our general framework.


%While one could formulate the result more broadly by allowing convex combinations over a larger family—namely, by asserting that any conditional entropy can be expressed as a convex combination of extremal conditional entropies for some choice of parameters—this description is redundant. It therefore suffices to restrict attention to the entropies $H_{t,\tau}$, $H_{0,\alpha}$, and $H_{-\infty,\tau}$ with parameters in the sets $\mathcal{D}_{\mathrm{bulk}}$, $\mathcal{D}_{0}$, and $\mathcal{D}_{-\infty}$, respectively, which precisely characterize the parameter regimes for which the corresponding extremal conditional entropies are monotone under conditionally mixing channels (see Definition~\ref{def: true sets}). Since the remaining axioms are readily verified, these entropies therefore satisfy all axioms required of a conditional entropy (see Definition~\ref{def: conditional entropy}).




%Concerning the proof strategy, the result follows directly from Theorem~\ref{th: large sample and catalytic}, which characterizes the monotones providing sufficient conditions for transformations between multiple copies of joint probability distributions under conditionally mixing channels (large-sample conditional majorization). The argument is then completed by invoking standard results that connect large-sample characterizations with barycentric decompositions. Specifically, these results imply essentially that any quantity that is both monotone and additive can be written as a convex combination of the functions that provide sufficient conditions for the corresponding large-sample transformations (see, e.g.,~\cite{mu21_renyi, haapasalo_inprep}).

\subsection{The general form of conditional entropies}

Let us note that the set of conditional entropies characterized in Definition \ref{def: conditional entropy} is convex. This means that, whenever $\mb H_1$ and $\mb H_2$ are conditional entropies (i.e., satisfy the postulates of Definition \ref{def: conditional entropy}) and $\lambda\in[0,1]$, then also $\mb H=\lambda\mb H_1+(1-\lambda)\mb H_2$,
\begin{equation}
\mb H(X|Y)_P=\lambda\mb H_1(X|Y)_P+(1-\lambda)\mb H_2(X|Y)_P\qquad\forall\, P=P_{XY}\,,
\end{equation}
is a conditional entropy. Especially all convex combinations of the conditional entropies $H_{t,\tau}$, $(t,\tau)\in\mc D_{\rm bulk}$, and $H_{-\infty,\tau}$, $\tau\in\mc D_{-\infty}$, and their pointwise limits are also conditional entropies. This fact also extends to the generalized convex combinations of these quantities, i.e.,\ barycentres over the set of these special conditional entropies with respect to any (inner and outer regular) probability measure. For the latter statement, we need the fact that the set of these special conditional entropies is a compact Hausdorff space with respect to the topology of pointwise convergence \cite[Proposition 8.5]{fritz2023abstractII}. The following theorem characterizing the total set of conditional entropies as described in Definition \ref{def: conditional entropy} tells us that any conditional entropy can be expressed as a generalized convex combination of the special conditional entropies characterized in this work.

\begin{theorem}[General form of conditional entropies]
\label{th:barycenter conditional entropies}
Let $\H$ be a conditional entropy in the sense of Definition~\ref{def: conditional entropy}. 
Then $\H$ can be represented as a generalized convex combination of the conditional entropies
$H_{t,\tau}, H_{-\infty,\tau}, H_{0,\alpha}$ and $H_{+\infty,0}$,
introduced in Definition~\ref{def: extremal entropies} with parameters $(t,\tau)\in\mathcal{D}_{\mathrm{bulk}}$, $\tau\in\mathcal{D}_{-\infty}$, and $\alpha\in[0,1]$, respectively, for the first three families. 
\end{theorem}
We do not give a full proof here, as this result is proven analogously to Proposition \ref{proposition:barycenter entropies} and it is a special case of Theorem 7 of \cite{haapasalo_inprep}. Let us however give a brief overview: Let us denote by $\mc H$ the set of conditional entropies $H_{t,\tau}$, $H_{-\infty,\tau}$, $H_{0,\alpha}$, and $H_{+\infty,0}$ as defined in \eqref{eq:TempEnt}, \eqref{eq:TropEnt}, and \eqref{eq:DerivEnt} of Definition~\ref{def: extremal entropies}, where the parameters and measures satisfy the required monotonicity conditions. According to Proposition 8.5 of \cite{fritz2023abstractII}, this set is a compact Hausdorff space when equipped with the topology of pointwise convergence. Just as in the proof of Proposition \ref{proposition:barycenter entropies}, now using Theorem \ref{th: large sample and catalytic}, we may easily prove that
\begin{equation}
\Gamma(X|Y)_P \leq \Gamma(X'|Y')_Q,
\;\; \forall \Gamma \in \mc H \;\; \implies \;\; \H(X|Y)_P \leq \H(X'|Y')_Q, \;\; \forall \H \;\text{in Definition~\ref{def: conditional entropy}} \,.
\end{equation}
Using this as a stepping stone, we may define a positive linear functional $G$ on the vector subspace of the set $C(\mc H)$ of continuous real functions on the compact Hausdorff $\mc H$ generated by the evaluation functions ${\rm ev}_{P}$, ${\rm ev}_P(\Gamma)=\Gamma(X|Y)_P$, through $G({\rm ev}_P)=\mb H(X|Y)_P$. Using Kantorovich's theorem we may extend $G$ into a positive linear functional $I:C(\mc H)\to\R$ and the Riesz-Markov-Kakutani theorem implies the existence of a finite positive measure $\nu:\mc B(\mc H)\to\R_+$ (Borel $\sigma$-algebra of the compact Hausdorff space $\mc H$) which is inner and outer regular such that
\begin{equation}
\mb H(X|Y)_P=\int_{\mc H}\Gamma(X|Y)_P\,\mathrm{d}\nu(\Gamma)\,.
\end{equation}
This measure is easily seen to be a probability measure.

We now know that all conditional entropies as defined in Definition \ref{def: conditional entropy} are spanned by the conditional entropies of Definition \ref{def: extremal entropies} with the parameter ranges of Definition \ref{def: true sets}. We also know that the conditional entropies $H_{t,\tau}$, $(t,\tau)\in\mc D_{\rm bulk}$, $H_{-\infty,\tau}$, $\tau\in\mc D_{-\infty}$, $H_{+\infty,0}$, and $H_{0,\alpha}$, $\alpha\in[0,1]$, are the extreme points of the convex set of conditional entropies. Currently, we are left with some uncertainty of the exact characterization of the parameter sets $\mc D_{\rm bulk}$ and $\mc D_{-\infty}$. However, we subsequently discuss steps towards proper identification of these sets.

% Although Theorem~\ref{th:barycenter conditional entropies} involves three different conditional entropies, its core message admits a clear interpretation. Conceptually, there is a unique conditional entropy defined as
% \begin{equation}
%     H_{t,\tau}(X|Y)_P = \frac{1}{t} \log \Biggl( \sum_{i \in \Y} P_Y(y_i) \exp\Bigl( t \int_{[0,+\infty]} H_\alpha(P_{X|Y=y_i}) \, \mathrm{d}\tau(\alpha) \Bigr) \Biggr)\,.
% \end{equation}
% The conditional entropies $H_{-\infty,\tau},\, H_{0,\alpha}$ are just limiting cases of \( H_{t,\tau} \). Specifically, the conditional entropy \( H_{-\infty,\tau} \) arises in the limit where $t\to-\infty$ 
% while the entropy \( H_{0,\alpha} \), emerges in the limit where $t\to0$ and the measure $\tau$ is the point measure at $\alpha$. Finally, since a convex combination of two conditional entropies again defines a conditional entropy, convex combinations must also be considered.

\bigskip

\subsection{Towards a complete characterization of the sets $\D_{\rm bulk}$ and $\D_{-\infty}$}
\label{sec: towards complete characterization}
In this section, we study the sets $\D_{\rm bulk}$, $\D_{-\infty}$, and $\D_0$ introduced in Definition~\ref{def: true sets} for which the conditional entropies $H_{t,\tau}$, $H_{-\infty,\tau}$, and $H_{0,\tau}$ defined in Definition~\ref{def: extremal entropies} are monotone under conditionally mixing channels. This monotonicity property constitutes the key axiom among those required for a conditional entropy according to Definition~\ref{def: conditional entropy}, as the remaining axioms are readily verified.
In particular, we provide a comprehensive discussion of the results obtained in Sections~\ref{sec: sufficient conditions measure} and~\ref{sec: necessary conditions measures} concerning the parameter regimes for which these quantities are monotone, and we highlight remaining open questions.

The results of Section~\ref{sec: sufficient conditions measure} yield sufficient conditions on the parameters under which the entropies $H_{t,\tau}$, $H_{-\infty,\tau}$, and $H_{0,\alpha}$ belong to the sets $\mathcal{D}_{\mathrm{bulk}}$, $\mathcal{D}_{-\infty}$, and $\mathcal{D}_{0}$, respectively. Accordingly, we introduce below the sets $\widetilde{\mathcal{D}}_{\mathrm{bulk}}$, $\widetilde{\mathcal{D}}_{-\infty}$, and $\widetilde{\mathcal{D}}_{0}$, defined as the collections of parameters satisfying these sufficient conditions. In addition, the results of Section~\ref{sec: necessary conditions measures} provide a collection of necessary conditions for the sets $\mathcal{D}_{\mathrm{bulk}}$, $\mathcal{D}_{-\infty}$, and $\mathcal{D}_{0}$.
These results show that, under additional assumptions—including the particularly relevant case of discrete measures with finite support—the criteria defining the set $\widetilde{\mathcal{D}}_{\mathrm{bulk}}$ are not only sufficient but also necessary, yielding
\(\widetilde{\mathcal{D}}_{\rm bulk} = \mathcal{D}_{\rm bulk}\). Moreover, for \(H_{0,\tau}\), these necessary conditions provide an exact characterization, demonstrating that \(\widetilde{\mathcal{D}}_0 = \mathcal{D}_0\) without any additional assumptions. 
Finally, we conjecture that the sufficient conditions we have established are, in fact, generally necessary.



We begin by formally defining the sets $\widetilde{\D}_{\rm bulk}$, $\widetilde{\D}_0$, and $\widetilde{\D}_{-\infty}$.
\begin{definition}
\label{def: D_bulk}
Let \(\widetilde{\D}_{\rm{bulk}}\) be the set of pairs $(t,\tau)$ consisting of $t\in\R\setminus\{0\}$ and probability measures $\tau:\mc B\big([0,+\infty]\big)\to\mb R$ satisfying one of the following conditions
\begin{enumerate}[itemsep=0pt,topsep=5pt]
    \item $t > 0$, $\;\textup{supp}(\tau) \subseteq [0,1]$, and
    \begin{equation}\label{eq:Itaas}
    \int_{[0,+\infty]}\frac{\alpha}{1-\alpha}\,\mathrm{d}\tau(\alpha)\leq\frac{1}{t}
    \end{equation}
    or
    \item $t < 0$ and
    \begin{itemize}
    \item[(a)] $\;\textup{supp}(\mu) \subseteq [0,1]$ or
    \item[(b)] $\;\operatorname{supp}(\tau) \subseteq [0,1] \cup \{\alpha_*\}$ for some $\alpha_* \in (1, +\infty]$ so that \eqref{eq:Itaas} is satisfied.
    \end{itemize}
\end{enumerate}
    % \begin{enumerate}[itemsep=0pt,topsep=5pt]
    % \item $t > 0$, $\;\textup{supp}(\tau) \subseteq [0,1]$, $\;\displaystyle\int_{[0,1]} \frac{\alpha}{1 - \alpha} \, \mathrm{d}\tau(\alpha) \leq 1/t$.
    % \item $t > 0$, $\;\textup{supp}(\mu) \subseteq [0,1]$.
    % \label{it: 2. for mu}
    % \item $t < 0$, $\;\operatorname{supp}(\tau) \subseteq [0,1] \cup \{\alpha_k\}$, $\;\alpha_k \in (1, +\infty]$, $\;\displaystyle\int_{[0,+\infty]} \frac{\alpha}{1 - \alpha} \, \mathrm{d}\mu(\alpha) \leq 1/t$.
    % \label{it: 3. for mu}
%\end{enumerate}
\end{definition}
In the integral~\eqref{eq:Itaas}, we adopt the convention that it equals to $+\infty$ if the measure $\tau$ has support at $\alpha=1$.
Hence, while in the first case of  $t<0$, the measure $\tau$ may have support at $\alpha=1$, in the second case of \(t < 0\) and for \(t > 0\) the measure cannot have support at \(\alpha = 1\). Indeed, our convention implies that the integral diverges, and consequently, the bound in \eqref{eq:IntegralCond} would be violated.

\begin{definition}
\label{def: D_infty}
 Let \(\widetilde{\D}_{-\infty}\) be the set of probability measures $\tau:\mc B\big([0,+\infty]\big)\to\mb R_+$ satisfying one of the following conditions
    \begin{enumerate}[itemsep=0pt,topsep=5pt]
    \item  $\;\textup{supp}(\tau) \subseteq [0,1]$.
    \item $\;\operatorname{supp}(\tau) \subseteq [0,1] \cup \{\alpha_*\}$ for some $\;\alpha_* \in (1, +\infty]$ so that $\;\displaystyle\int_{[0,+\infty]} \frac{\alpha}{1 - \alpha} \, \mathrm{d}\tau(\alpha) \leq 0$.
\end{enumerate}
\end{definition}

% \begin{definition}
% \label{def: D_0}
% Let \(\widetilde{\D}_{0}\) denote the set of all $\alpha \in [0,1]$.
% \end{definition}
In the following proposition, we collect the results on the necessary and sufficient conditions for monotonicity derived in Sections~\ref{sec: sufficient conditions measure} and~\ref{sec: necessary conditions measures}.
\begin{proposition}
\label{prop: necessary conditions}
The sets of parameters satisfy $\widetilde{\mathcal{D}}_{\rm bulk} \subseteq \mathcal{D}_{\rm bulk}$ and 
$\widetilde{\mathcal{D}}_{-\infty} \subseteq \mathcal{D}_{-\infty}$. 
Moreover, if we impose the additional constraints on the measure $\tau$
\begin{equation}
\label{eq: finite integrals}
\left|\int_{[0,1)} \frac{\alpha}{1-\alpha}\, \mathrm{d}\tau(\alpha)\right|\,,
\left|\int_{(1,+\infty]} \frac{\alpha}{1-\alpha}\, \mathrm{d}\tau(\alpha)\right| < +\infty,
\end{equation}
then $t$ and $\tau$ must satisfy the conditions characterizing $\widetilde{\mathcal{D}}_{\rm bulk}$ for $(t,\tau)\in\mathcal{D}_{\rm bulk}$.
\end{proposition}
\begin{proof}
    The equality \(\widetilde{\mathcal{D}}_{0} = \mathcal{D}_{0}\) follows from Propositions~\ref{prop: sufficient derivations} and~\ref{prop: prop measure deriv}, together with the equivalence between convexity and monotonicity properties discussed in Lemma~\ref{lem: conv derivations} of Appendix~\ref{app: convexity and monotonicity}.
 Similarly, the inclusions $\widetilde{\mathcal{D}}_{\rm bulk} \subseteq \mathcal{D}_{\rm bulk}$, 
$\widetilde{\mathcal{D}}_{-\infty} \subseteq \mathcal{D}_{-\infty}$ follow from the sufficient conditions derived in Propositions~\ref{prop: sufficient conditions} and~\ref{prop: sufficient tropical}. The last statement of the above proposition %about the equality for the set $\widetilde{\mathcal{D}}_{\rm bulk} = \mathcal{D}_{\rm bulk}$
is a consequence of Propositions~\ref{prop: first prop measure},~\ref{prop: second prop measure},~\ref{prop: third prop measure}, while the finiteness of the integrals in~\eqref{eq: finite integrals} is a consequence of Assumption~\ref{assumption:finite integrals} included in the propositions.
\end{proof}
In particular, we note that whenever $\tau$ is a discrete measure supported on a finite number of points, the condition in equation~\eqref{eq: finite integrals} is satisfied, and hence for this case we obtain a tight characterization of the set $\D_{\mathrm{bulk}}$.
The counterexamples constructed and analyzed in Section~\ref{sec: necessary conditions measures} strongly suggest the following conjecture.

\begin{conjecture}
\label{conj; true sets}
We propose that $\widetilde{\D}_{\rm bulk}=\D_{\rm bulk}$ and $\widetilde{\D}_{-\infty}=\D_{-\infty}$.
\end{conjecture}

We end this section by emphasizing that Theorem~\ref{th:barycenter conditional entropies} provides the explicit form of a conditional entropy. This result regarding their structure holds independently of the validity of the above conjecture and, in particular, does not rely on the assumption stated in equation~\eqref{eq: finite integrals} for the necessity of the monotonicity, which remains an important open problem.




\subsection{Conditional entropies in the literature}
\label{sec: literature}
In this section, we review in detail several conditional entropies that have been introduced in the literature and show how each of these quantities arises as a special case of the general form established in Theorem~\ref{th:barycenter conditional entropies}.



The conditional Shannon entropy $H(X|Y)$,
\begin{equation}
    H(X|Y)_P=-\sum_{x \in \X,y\in\Y} P_{XY}(x,y)\log{\frac{P_{XY}(x,y)}{P_Y(y)}} \,,
\end{equation}
is equal to \(H_{0,\alpha}\) for \(\alpha = 1\).
The Hayashi conditional entropy introduced by Hayashi and Skorić et al.~\cite{hayashi2011exponential,vskoric2011sharp}
\begin{equation}
H^\text{H}_{\alpha}(X|Y)_P=\frac{1}{1-\alpha}\log{\Big(\sum\nolimits_{x\in\X,y\in\Y} P_Y(y) P_{X|Y}(x|y)^\alpha\Big)} \,,
\end{equation}
arises as a special case of $H_{t,\tau}$ when $\tau$ is the point at $\alpha$ and $t = 1 - \alpha$. The Arimoto Rényi conditional entropy~\cite{arimoto1977information} 
\begin{align}
H^\text{A}_{\alpha}(X|Y)_P=\frac{\alpha}{1-\alpha}\log{\bigg(\sum_{y\in \Y} P_Y(y)\Big(\sum_{x \in \X} P_{X|Y}(x|y)^\alpha\Big)^\frac{1}{\alpha}\bigg)} \,,
\end{align}
is recovered similarly as a special case of $H_{t,\tau}$ when $\tau$ is the point measure at $\alpha$ and $t = \frac{1 - \alpha}{\alpha}$. A two-parameter family that interpolates between the Arimoto and Hayashi conditional entropies was introduced in~\cite{hayashi2016equivocations} and later investigated in detail in~\cite{rubboli2024quantum}
\begin{align}
\label{eq: H alpha beta}
H_{\alpha,\beta}(X|Y)_P=\frac{\alpha}{1-\alpha}\frac{1}{\beta}\log{\bigg(\sum_{y \in \Y} P_Y(y)\Big(\sum_{x \in \X} P_{X|Y}(x|y)^\alpha\Big)^\frac{\beta} {\alpha}\bigg)} \,.
\end{align}
This quantity is a special case of $H_{t,\tau}$ when $\tau$ is the point measure at $\alpha$ and $t = \frac{1 - \alpha}{\alpha} \beta$. More recently, this quantity has been studied in~\cite{li2025two}.


More exotic quantities have appeared in the literature. In~\cite{tan2018analysis}, the authors defined the conditional entropy
\begin{equation}
    H_{1+a|1+b}(X|Y)_P=-\frac{1+b}{a}\log{\bigg(\sum_{y \in \Y} P_Y(y)\Big(\sum_{x \in \X} P_{X|Y}(x|y)^{1+a}\Big)\Big(\sum_{x \in \X} P_{X|Y}(x|y)^{1+b}\Big)^{-\frac{a}{1+b}}\bigg)} \,.
\end{equation}
This corresponds to $H_{t,\tau}$ for the two-point measure $\tau=(1+b)\delta_{1+a}-b\delta_{1+b}$ and $t=-a/(1+b)$.
Another conditional entropy has been introduced in~\cite{hayashi2016uniform}
\begin{align}
    H_{1+\theta|1+\theta'}(X|Y)_P=&-\frac{1}{\theta}\log{\bigg(\sum_{y \in \Y} P_Y(y)\Big(\sum_{x \in \X} P_{X|Y}(x|y)^{1+\theta}\Big)\Big(\sum_{x \in \X} P_{X|Y}(x|y)^{1+\theta'}\Big)^{-\frac{\theta}{1+\theta'}}\bigg)} \\
    &\qquad + \frac{\theta'}{1+\theta'}H^{\text{A}}_{1+\theta'}(P_{XY}) \\
    =&\frac{1}{1+\theta'}H_{t_1,\tau_1}(X|Y)_P+\frac{\theta'}{1+\theta'}H_{t_2,\tau_2}(X|Y)_P\,,
\end{align}
Here, $\tau_1=(1-\theta')\delta_{1+\theta}-\theta'\delta_{1+\theta'}$, $t_1=-\theta/(1+\theta')$, $\tau_2$ is the point measure at $1+\theta'$, and $t_2=-\theta'/(1+\theta')$.

The conditional entropy
 \begin{align}
H_{0,\alpha}(P)=\sum_{y\in\mc Y} P_Y(y)H_\alpha(P_{X|Y=y})\,
 \end{align}
 was introduced by Cachin~\cite{cachin1997entropy} and later studied in~\cite{teixeira2012conditional}.

In~\cite{renner2005simple},
as a part of a study of conditional smooth R\'enyi entropies, the authors introduced the quantity
\begin{align}
    H^R_\alpha(X|Y) = \min _{y} H_\alpha(X|Y=y)
\end{align}
for $\alpha > 1$ and with $\min$ replaced by $\max$ for $\alpha \in (0,1)$. In the range $\alpha>1$, it coincides with $H_{-\infty,\tau}$ when $\tau$ is a point measure at $\alpha$.
For $\alpha \in (0,1)$, it coincides with the limit $t\rightarrow +\infty$ of $H_{t,\tau}$ in the case where $\tau$ is a point measure at $\alpha$. However, although this limit is true for all $\alpha \in (0,1)$, Theorem~\ref{th:barycenter conditional entropies} implies that these quantities satisfy all our axioms required for a conditional entropy  (Definition~\ref{def: conditional entropy}) only in the case $\alpha = 0$.
In this case, it recovers \(H_{+\infty,0}\).  

Finally, the quantities $ H_{+\infty,0}$ and $H_{-\infty,\tau}$ for single-point measures have been studied as limiting cases of
$H_{\alpha,\beta}$~\eqref{eq: H alpha beta} in~\cite{li2025two} and~\cite{rubboli2024quantum}.



 
\section{Applications}
\label{sec: Applications}
\subsection{Rate of transformation under conditional majorization}
As a first application of our result, we show that the conditional entropies $H_{t,\tau}$ determine the optimal transformation rate under conditionally mixing operations and isometries, thereby providing them with an operational interpretation as characterizing the number of output copies obtainable per input copy within this class of operations.

Specifically, given two joint probability distributions $P_{XY}$ and $Q_{X'Y'}$, we seek the largest $r \in \mathbb{R}$ for which $P_{XY}^{\otimes n} \succeq Q_{X'Y'}^{\otimes \lfloor r n \rfloor}$ holds for all sufficiently large $n$, i.e., we want to determine the value
\begin{equation}
    R(P_{XY} \rightarrow Q_{XY}) =\; \sup\left\{ r \ge 0 \;\middle|\; P_{XY}^{\otimes n}\succeq Q_{X'Y'}^{\otimes \lfloor rn \rfloor} \;\; \text{for large } n \right\}.
\end{equation}
It is known that this problem is equivalent to the large sample problem (see e.g.~\cite[Theorem 3.11]{Jensen_Kjaerulf_2019}). In particular, we use the necessary and sufficient conditions for large sample of Theorem~\ref{th: large sample and catalytic} to obtain
\begin{theorem}
\label{thm: rate}
    Let $P_{XY}$ and $Q_{XY}$ be two joint probability distributions. Then, 
    \begin{equation}
    \label{eq: rate}
R(P_{XY}\rightarrow Q_{X'Y'}) = \inf\left\{ \frac{H_{t,\tau}(X'|Y')_Q}{H_{t,\tau}(X|Y)_P}\Bigg| (t,\tau) \in \D_{\rm{bulk}}\right\} \cup \left\{ \frac{H_{-\infty,\tau}(X'|Y')_Q}{H_{-\infty,\tau}(X|Y)_P}\Bigg| \tau \in \D_{-\infty}\right\},
\end{equation}
\end{theorem}
\begin{proof}
   The proof follows a similar argument to that of~\cite[Theorem 3.11]{Jensen_Kjaerulf_2019}, with the additional observation that the derivations $H_{0,\alpha}$ and the tropical homomorphism $H_{+\infty,0}$ can be excluded by taking the infimum in equation~\eqref{eq: rate}, since they are limits of $H_{t,\tau}$. Indeed, these derivations are monotone only for $\alpha \in [0,1]$, a region that can itself be viewed as the limit of $\widetilde{\D}_{\rm bulk}$ (see discussion in Section~\ref{sec: towards complete characterization}). The same argument also holds for $H_{+\infty,0}$. We also refer to the very similar proof of Theorem 5 of \cite{verhagen2025}.
\end{proof}

Let us remark that if Conjecture \ref{conj; true sets} holds, we may simplify the statement of Theorem \ref{thm: rate}. Indeed, if the conjecture is true, the conditional entropies $H_{-\infty,\tau}$ with $\tau\in\mc D_{-\infty}$ are all pointwise limits of $(t,\tau)\in\mc D_{\rm bulk}$ as $t\to-\infty$. This would imply the conjectural form
\begin{equation}
R(P_{XY}\rightarrow Q_{X'Y'})=\inf_{(t,\tau)\in\mc D_{\rm bulk}}\frac{H_{t,\tau}(X'|Y')_Q}{H_{t,\tau}(X|Y)_P}
\end{equation}
for the optimal rate.


\subsection{Conditional channels under conservation conditions and an application to second laws of thermodynamics with side information}
\label{sec: second laws of thermo}

A doubly stochastic or mixing channels $S$ are exactly those (column) stochastic matrices that fix a suitably large uniform distribution $U$, i.e., $SU=U$. Conditionally mixing channels are thus conditionally uniformity-preserving channels. It is thus quite natural to consider a generalization of our results where the conserved primary distribution is not uniform $U$ but some other fixed distribution $R$. For this purpose, let us introduce the following set of operations that generalize conditionally mixing channels.

\begin{definition}
Let us fix a finite random variable $X$ together with the probability distribution $R_X$. A \textit{conditioned $R_X$-preserving channel} has the form
\begin{equation}
P_{XY}\mapsto Q_{XY'}=\sum_{i=1}^k G^{(i)}P_{XY}D^{(i)}
\end{equation}
where
\begin{enumerate}
\item the matrices $G^{(i)}$ are column-stochastic and $G^{(i)}R_X=R_X$ for $i=1,\ldots,k$,
\item the matrices $D^{(i)}$ are substochastic for $i=1,\ldots,k$, and
\item the sum $\sum_{i=1}^k D^{(i)}$ is row-stochastic.
\end{enumerate}
\end{definition}

We are interested in conditions under which \(P_{XY}\) can be transformed into \(Q_{XY'}\) catalytically by a conditionally \(R_X\)-preserving channel, with arbitrarily small error $\varepsilon$.
The error here is quantified with the total variation distance defined between joint probabilities $S_{XY}$ and $T_{XY}$ through
\begin{equation}
d(S_{XY},T_{XY})=\frac{1}{2}\sum_{x\in X}\sum_{y\in Y}\big|S_{XY}(x,y)-T_{XY}(x,y)\big|.
\end{equation}
We introduce the quantities $I_{t,\tau}$,
\begin{equation}
I_{t,\tau}(P_{XY}\|R_X)=-\frac{1}{t}\log{\left(\sum_{y\in Y}P_Y(y)\exp{\left(-t\int_{[0,+\infty]}D_\alpha(P_{X|Y=y}\|R_X)\,\mathrm{d}\tau(\alpha)\right)}\right)} \,.
\end{equation}
Here, $t\in\R\setminus\{0\}$ and $\tau$ is a Borel probability measure on $[0,\infty]$. Moreover, the R\'enyi relative entropy for $\alpha=[0,+\infty]$ is defined as
\begin{equation}
    D_\alpha(P_X\|R_X)=
\frac{1}{\alpha-1}\log{\sum_{x\in X}P_X(x)^\alpha R_X(x)^{1-\alpha}}\,,
\end{equation}
with the cases \(\alpha=1\) and \(\alpha=+\infty\) defined by the corresponding limits.
 We also define \(I_{-\infty,\tau}\) as the pointwise limit of \(I_{t,\tau}\) for \(t \to -\infty\). In the following, we provide a set of sufficient conditions under which a joint distribution can be transformed into another using conditioned preserving channels and the assistance of a catalyst.
\begin{theorem}\label{thm:ConservedClassical}
Let us denote the uniform distribution of any system $Z$ by $U_Z$. Fix a finite probability distribution $R_X$ with rational entries and assume that $P_{XY}$ and $Q_{XY'}$ are joint probability distributions such that
\begin{equation}
\max \limits_{\substack{y \in \mathcal{Y} \\ P_Y( y) > 0}} |{\rm supp}(P_{X|Y=y})|<|{\rm supp}(R_X)|,\quad Q_{X|Y'=y'}\neq R_X\ \forall\,y'\in \mc Y'.
\end{equation}
For all $\varepsilon>0$, there exist joint probability distributions $Q^\varepsilon_{XY'}$ and $S^\varepsilon_{X_\varepsilon Y_\varepsilon}$, such that $d(Q_{XY'},Q^\varepsilon_{XY'})<\varepsilon$ and there exists a conditionally $R_X\otimes U_{X_\varepsilon}$-preserving channel that maps $P_{XY}\otimes S^\varepsilon_{X\varepsilon Y\varepsilon}$ into $Q^\varepsilon_{XY'}\otimes S^\varepsilon_{X\varepsilon Y\varepsilon}$ if
\begin{align}
I_{t,\tau}(P_{XY}\|R_X)&\geq I_{t,\tau}(Q_{XY'}\|R_X)\qquad\forall(t,\tau)\in\mc D_{\rm bulk},\\
I_{-\infty,\tau}(P_{XY}\|R_X)&\geq I_{-\infty,\tau}(Q_{XY'}\|R_X)\qquad \forall\tau\in\mc D_{-\infty}.
\end{align}
\end{theorem}

\begin{proof}
   The proof follows a standard argument based on the so-called embedding map. Specifically, the conditions $I_{t,\tau}(P_{XY},R_X) \geq I_{t,\tau}(Q_{XY'},R_X)$ and $I_{-\infty,\tau} (P_{XY},R_X) \geq I_{-\infty,\tau}(Q_{XY'},R_X)$ are equivalent to the inequalities $H_{t,\tau}(\widetilde{P}_{\widetilde{X}Y}) \leq H_{t,\tau}(\widetilde{Q}_{\widetilde{X}Y'})$  and $H_{-\infty,\tau}(\widetilde{P}_{\widetilde{X}Y}) \leq H_{-\infty,\tau}(\widetilde{Q}_{\widetilde{X}Y'})$, where $\widetilde{P}_{\widetilde{X}Y}$ and $\widetilde{Q}_{\widetilde{X}Y'}$ are the embedded versions of $R_{XY}$ and $Q_{XY'}$, respectively, obtained via the application of the embedding map (see below for the formal definition).
   This entropy inequality guarantees the existence of a catalyst and a conditionally mixing channel 
   that transforms $\widetilde{R}_{XY}$ into $\widetilde{Q}_{XY'}$ according to Theorem~\ref{th: large sample and catalytic}. We then construct the desired overall channel by composing the embedding map, the conditionally mixing channel, and the (left) inverse of the embedding map. By construction, as we argue later, this overall transformation is $R_X$ preserving.

\bigskip

Let us denote the support size of $R_X$ by $d$. By assumption, $R_X$ has rational entries. Hence, there exists $D \in \mathbb{N}$ and $g_i \in \mathbb{N}$ with $D= \sum_{i=1}^d g_i$ such that
\begin{equation}
R_X = \left( \frac{g_1}{D}, \ldots, \frac{g_d}{D} \right)
\end{equation}
up to eliminating possible zero entries from $R_X$. The embedding map $\Gamma : \mathbb{R}^d \to \mathbb{R}^D$ is defined as follows
\begin{align}
\Gamma(p) := \Big( \underbrace{\frac{p_1}{g_1}, \ldots, \frac{p_1}{g_1}}_{g_1 \ \text{times}}, \ldots, \underbrace{\frac{p_d}{g_d}, \ldots, \frac{p_k}{g_d}}_{g_d \ \text{times}} \Big)\,.
\end{align}
This map has the crucial property of sending $R_X$ to the uniform distribution $U_D := \big(\tfrac{1}{D}, \ldots, \tfrac{1}{D}\big)$. This allows us to apply the results from Section~\ref{sec: large sample} on conditional majorization. 
The left inverse $\Gamma^{-1} : \mathbb{R}^D \to \mathbb{R}^d$ is 
\begin{equation}
\label{PsInvEmbed}
\Gamma^{-1}(q) := \left( \sum_{i=1}^{g_1} q_i,\ \sum_{i=g_1+1}^{g_1 + g_2} q_i,\ \ldots,\ \sum_{i = g_1 + \cdots + g_{k-1} + 1}^{D} q_i \right),
\end{equation}
It is straightforward to verify that (see e.g.,~\cite[Lemma 13]{Brandao})
\begin{align}
    D_{\alpha}(P_X \| R_X) = D_{\alpha}(\Gamma(P_X) \| U_D) = \log{D} - H_\alpha (\Gamma(P_X))\,.
\end{align}
The properties of the embedding map applied to the system $X$ imply that
\begin{align}
    I_{t,\tau}(P_{XY} \| R_X) &=  -\frac{1}{t} \log \Bigg( \sum_{y \in \mc Y} P_Y(y) \exp \Big( -t\int_{[0,+\infty]}D_\alpha(P_{X|Y=y}) \| R_{X}) \, \mathrm{d} \tau(\alpha) \Big) \Bigg) \\
    & =-\frac{1}{t} \log \Bigg( \sum_{y \in\mc Y} P_Y(y) \exp \Big( -t\int_{[0,+\infty]} D_\alpha(\Gamma_X(P_{X|Y=y}) \| U_D) \, \mathrm{d} \tau(\alpha) \Big) \Bigg) \\
    & = -\frac{1}{t} \log \Bigg( \sum_{y \in\mc Y} P_Y(y) \exp \Big( -t \log{D} +t\int_{[0,+\infty]} H_\alpha(\Gamma_X(P_{X|Y=y})) \, \mathrm{d} \tau(\alpha) \Big) \Bigg)\\
    & = -\frac{1}{t} \log \Bigg( \sum_{y \in\mc Y} P_Y(y) \exp \Big(t\int_{[0,+\infty]} H_\alpha(\Gamma_X(P_{X|Y=y})) \, \mathrm{d} \tau(\alpha) \Big) \Bigg) + \log{D}
\end{align}
Hence, we obtain that
\begin{align}
    I_{t,\tau}(P_{XY} \| R_X) = -H_{t,\tau}(X|Y)_{\Gamma_X(P_{XY})} + \log{D}
\end{align}

We note that the inequalities in the assumption of the theorem for $I_{t,\tau}$ extend naturally to the quantity $I_{0,\alpha}$, which can be defined as the limit $t \to 0$ of $F_{t,\tau}$, in direct analogy with the definition of $H_{0,\alpha}$ introduced in Section~\ref{sec: large sample}. Indeed, by a limiting argument, the condition $I_{0,\alpha}(P_{XY}\|R_X) \geq I_{0,\alpha}(Q_{XY'}\|R_X)$ in the region $\alpha \in [0,1]$ can be obtained as limits of the inequalities for $I_{t,\tau}$ in the bulk region $\mathcal{D}_{\rm bulk}$; see the discussion in Section~\ref{sec: towards complete characterization}.

To apply Theorem~\ref{th: large sample and catalytic}, we need to make the inequalities strict. Since these functions are defined in terms of convex combinations of Rényi relative entropies, it suffices to establish the strict inequality for each Rényi relative entropy individually, for all \(\alpha \ge 0\). Indeed, taking convex combinations preserves strict order, and the subsequent operations—namely summation, addition of constants, exponentiation, and application of the logarithm—are all strictly monotone. Consequently, the overall functions inherit the strict ordering. To make the inequalities strict, we construct the approximation
\begin{equation}
    Q^\varepsilon_{XY'} = \left(1-\varepsilon\right) Q_{XY'} + \varepsilon (R_X \otimes Q_{Y'}) \,.
\end{equation}
This distribution is an $\varepsilon$-approximation of $Q_{XY'}$:
\begin{align}
d(Q_{XY'},Q^\varepsilon_{XY'})&=\frac{1}{2}\sum_{x\in\mc X}\sum_{y'\in\mc Y'}\left|Q_{XY'}(x,y')-(1-\varepsilon)Q_{XY'}(x,y')-\varepsilon R_X(x)Q_{Y'}(y')\right|\\
&\leq\frac{\varepsilon}{2}\sum_{x\in\mc X}\sum_{y'\in\mc Y'}Q_{XY'}(x,y')+\frac{\varepsilon}{2}\sum_{x\in\mc X}R_X(x)\sum_{y'\in\mc Y'}Q_{Y'}(y')=\frac{\varepsilon}{2}+\frac{\varepsilon}{2}=\varepsilon.
\end{align}
The condition $Q_{X|Y'=y'}\neq R_X$ in the assumption of the claim guarantees that the inequalities between R\'{e}nyi relative entropies are strict for all $\alpha>0$. Indeed, for any $\alpha>0$ and $ y'\in \Y'$ we obtain
\begin{equation}
    D_\alpha\!\left(Q_{X|Y'=y'} \middle\| R_X \right)
    >
    D_\alpha\!\left(Q^\varepsilon_{X|Y'=y'} \middle\| R_X \right) .
\end{equation}
This follows from the joint concavity of the map $Q_\alpha(P_X\|R_X) = \exp\!\bigl((\alpha-1)D_\alpha(P_X\|R_X)\bigr)$ for $\alpha<1$ and its joint convexity for $\alpha>1$, together with the assumption that
$Q_{X|Y'=y'}\neq R_X$.
To illustrate this explicitly, consider the case $\alpha>1$. Then
\begin{align}
    Q_\alpha\!\left(Q^\varepsilon_{X|Y'=y'} \middle\| R_X \right)
    &\leq
    \left(1-\frac{\varepsilon}{2}\right)
    Q_\alpha\!\left(Q_{X|Y'=y'} \middle\| R_X \right)
    + \frac{\varepsilon}{2} <
    Q_\alpha\!\left(Q_{X|Y'=y'} \middle\| R_X \right) ,
\end{align}
where we used that the assumption $Q_{X|Y'=y'}\neq R_X$ implies $Q_\alpha\!\left(Q_{X|Y'=y'} \middle\| R_X \right)
>
Q_\alpha\!\left(R_X \middle\| R_X \right)
=1$.
By taking the logarithm, we obtain the desired strict inequality. Since the approximation decreases the value of $D_0$, and in particular satisfies $D_0(Q^{\varepsilon}_{X\mid Y'=y'} \| R_X)=0$, it follows that all $I_{t,\tau}$ are strictly ordered if the intersection of the support of $\tau$ and $(0,\infty]$ is non-empty. It therefore remains to consider the case in which the measure is supported entirely at $\alpha=0$.
In this case, the assumption that the support sizes of $P_{X|Y=y}$ are all strictly smaller than that of $R_X$ guarantees that the inequalities are also strict for $\alpha=0$. Indeed, for all $y\in\mc Y$ and $y'\in\mc Y'$ it holds that
 \begin{align}
     D_0 (P_{X|Y=y} \| R_X) \geq D_0 (Q_{X|Y'=y'} \| R_X) \geq D_0(Q^\varepsilon_{X|Y'=y'} \| R_X )  = 0 \,.
\end{align}
Here, the last inequality follows from the fact that $Q^\varepsilon_{X|Y'=y'} = (1-\varepsilon )Q_{X|Y'=y'} + \varepsilon R_X$ and hence has the same support of $R_X$. The assumption of the theorem on the support of $P_{XY}$ implies that $D_0(P_{X|Y=y}\|R_X)> 0$ for all $y\in \Y$.
 The preceding observations imply that
\begin{align}
    & H_{t,\tau}(X|Y)_{\Gamma_X(P_{XY})} < H_{t,\tau}(X|Y')_{\Gamma_X(Q^\varepsilon_{XY'})}, \qquad\qquad\qquad \forall (t,\tau) \in \D_{\rm{bulk}}, \\
    &H_{-\infty,\tau}(X|Y)_{\Gamma_X(P_{XY})} < H_{-\infty,\tau}(X|Y')_{\Gamma_X(Q^\varepsilon_{XY'})}, \quad\quad \quad  \;\;\; \forall \tau \in \D_{-\infty}, \\
    &H_{0,\alpha}(X|Y)_{\Gamma_X(P_{XY})} < H_{0,\alpha}(X|Y')_{\Gamma_X(Q^\varepsilon_{XY'})}, \quad\,\qquad \quad\quad\;\; \forall \alpha \in [0,1] \,,\\
    &H_{+\infty,0}(X|Y)_{\Gamma_X(P_{XY})} < H_{+\infty,0}(X|Y')_{\Gamma_X(Q^\varepsilon_{XY'})}.
\end{align}
Theorem~\ref{th: large sample and catalytic} implies the existence of a catalyst $S^\varepsilon_{X\varepsilon Y\varepsilon}$ and a conditionally mixing channel $\Upsilon$ such that
\begin{equation}
    \Upsilon(\Gamma_X (P_{XY}) \otimes S^\varepsilon_{X\varepsilon Y\varepsilon})= \Gamma_X (Q^\varepsilon_{XY'}) \otimes S^\varepsilon_{X\varepsilon Y\varepsilon} \,.
\end{equation}
Note that no additional embedding is required, since $P_{XY}$ and $Q^\varepsilon_{XY'}$ have the same number of rows, as discussed in Remark~\ref{remark: no isometries}.
Therefore, we construct the channel
\begin{equation}
	\Lambda := (\Gamma^{-1}_X \otimes I_{Y'}\otimes I_{X\varepsilon Y\varepsilon}) \circ \Upsilon_{XY'X\varepsilon Y\varepsilon} \circ (\Gamma_X \otimes I_{Y'}\otimes I_{X\varepsilon Y\varepsilon})
\end{equation}
so that
\begin{align}
	&\Lambda (P_{XY} \otimes S^\varepsilon_{X\varepsilon Y\varepsilon})  = Q^\varepsilon_{XY'} \otimes S^\varepsilon_{X\varepsilon Y\varepsilon}%\\
    %&\Lambda (R_{X}\otimes \eta_Y \otimes \pi_{X''Y''})  = \gamma_{X}\otimes \tau_Y \otimes \pi_{X''Y''} \,,
\end{align}
%where the last inequality holds for any state $\eta_Y$ and $\tau_Y$ is a state on $Y$. 
Finally, to show that this map is a conditionally $R_X\otimes U_{X_\varepsilon}$-preserving operation, one can verify its structure directly. Another way is to observe that, for any distribution $T_{YY_\varepsilon}$, there is a distribution $T'_{Y'Y_\varepsilon}$ such that
\begin{equation}
\Upsilon\big(U_D\otimes U_{X_\varepsilon}\otimes T_{YY_\varepsilon}\big)=U_D\otimes U_{X_\varepsilon}\otimes T'_{Y'Y_\varepsilon}
\end{equation}
which immediately follows from the fact that $\Upsilon$ is conditionally mixing and implies that
\begin{equation}
\Lambda(R_{X}\otimes U_{X_\varepsilon}\otimes T_{YY_\varepsilon})=R_X\otimes U_{X_\varepsilon
}\otimes T'_{Y'Y_\varepsilon}.
\end{equation}
Moreover, $\Lambda$ does not signal from $XX_\varepsilon$ to $Y'Y_\varepsilon$ meaning that ignoring the $XX_\varepsilon$ information after performing $\Lambda$ results into a legitimate stochastic map $YY_\varepsilon\to Y'Y_\varepsilon$ as one easily verifies using the fact that $\Upsilon$ is conditionally mixing. These two properties together are equivalent to $\Lambda$ being a conditionally $R_X\otimes U_{X_\varepsilon}$-preserving channel, by an argument analogous to that in~\cite[Section C]{gour2024inevitability}. 
\end{proof}

Let us present a direct physical application of the above result. Let us fix a finite-dimensional quantum system $A$ with the Hamiltonian
\begin{equation}
H_A=\sum_a E_a\ketbra{a}{a}
\end{equation}
where $\{|a\rangle\}_a$ is the energy eigenbasis of the Hilbert space $\mc H_A$ of the system. If all $E_a$ coincide, we say that the Hamiltonian is trivial. In the inverse temperature $\beta$ the Hamiltonian defines the partition function $Z_A$ and the Gibbs state $\gamma_A$ through
\begin{equation}
Z_A=\sum_a e^{a\beta E_a},\qquad \gamma_A=\frac{1}{Z_A}e^{-\beta H_A}.
\end{equation}

A quantum channel $G$, i.e., a completely positive trace-preserving linear map on the operators on $\mc H_A$, is Gibbs-preserving if $G(\gamma_A)=\gamma_A$. Thermal operations are defined as those that allow one to borrow an ancillary system in a Gibbs state, apply a unitary that commutes with the total Hamiltonian, and then discard a subsystem via the partial trace. For classical states, i.e.,\ those that commute with $\gamma_A$, however, thermal operations are equivalent to Gibbs-preserving maps~\cite{janzing2000thermodynamic,Oppenheim}. Accordingly, in the following, we will use the term Gibbs-preserving maps interchangeably with thermal operations.

Conditioned thermal operations generalize this concept in the setting where side information is available and were introduced in~\cite{narasimhachar2017resource}. We say that a channel $G$ from a bipartite quantum system $AB$ to another bipartite system $AB'$ is conditionally Gibbs-preserving if there exists a channel $G_{B\to B'}$ such that
\begin{equation}
\ptr{A}{G(\rho_{AB})}=G_{B\to B'}(\rho_B),\quad G(\gamma_A\otimes\rho_B)=\gamma_A\otimes G_{B\to B'}(\rho_B)
\end{equation}
for all quantum states $\rho_{AB}$ on $AB$ with the reduced state $\rho_B$ on $B$.

We say that a state $\rho_{AB}$ is classical if there is a joint probability distribution $P_{AB}$ and an orthonormal basis $\{|b\rangle\}_b$ for the Hilbert space $\mc H_B$ of $B$ such that
\begin{equation}
\rho_{AB}=\sum_{a\in\mc A}\sum_{b\in\mc B}P_{AB}(a,b)\ketbra{a}{a}\otimes\ketbra{b}{b}.
\end{equation}
It follows immediately that, if $\rho_{AB}$ and $\sigma_{AB'}$ are classical states, where $\rho_{AB}$ is associated with the joint probability $P_{AB}$ and $\sigma_{AB'}$ with the joint probability $Q_{AB'}$, and there is a conditionally Gibbs-preserving channel $G$ such that $G(\rho_{AB})=\sigma_{AB'}$, then there is a classical conditionally $R_A$-preserving channel $\Lambda$, where $R_A(a)=e^{-\beta E_a}/Z_A$, such that $\Lambda(P_{AB})=Q_{AB'}$. Naturally, with this observation, we obtain a sufficient condition for asymptotic conditionally Gibbs-preserving transformability between classical states in the catalytic regime as a direct corollary of Theorem \ref{thm:ConservedClassical}. However, before stating this result, let us introduce the free energies
\begin{equation}
F_{t,\tau}(\rho_{AB},\gamma_A)=\frac{1}{\beta}\big(I_{t,\tau}(P_{AB}\|R_A)-\log{Z_A}\big)
\end{equation}
for any classical state $\rho_{AB}$ associated with the joint probability $P_{AB}$ where $R_A(a)=e^{-\beta E_a}/Z_A$. We also define the free energies $F_{-\infty,\tau}$ as pointwise limits of the above as $t\to-\infty$.

\begin{corollary}[Second laws of thermodynamics with side information]\label{cor:secondlaws}
Let us assume that the Gibbs-spectrum is rational, i.e., $e^{-\beta E_a}\in\mb Q$ for all $a$. Suppose that
\begin{equation}
\rho_{AB}=\sum_{a,b}P_{AB}(a,b)\ketbra{a}{a}\otimes\ketbra{b}{b},\quad\sigma_{AB'}=\sum_{a,b'}Q_{AB'}(a,b')\ketbra{a}{a}\otimes\ketbra{b'}{b'}
\end{equation}
are classical states such that
\begin{equation}
\max_{b,\ P_B(b)>0}\big|{\rm supp}(P_{A|B=b})\big|<d_A,\qquad Q_{A|B'=b'}\neq\frac{e^{-\beta E_a}}{Z_A}\ \forall(a,b'),
\end{equation}
where $d_A$ is the dimension of the Hilbert space $\mc H_A$. For all $\varepsilon>0$, there exist finite quantum systems $A_\varepsilon$ and $B_\varepsilon$, a classical state $\sigma^\varepsilon_{AB'}$, a state $\omega^\varepsilon_{A_\varepsilon B_\varepsilon}$ and a conditionally Gibbs-preserving channel $G_\varepsilon:AA_\varepsilon BB_\varepsilon\to AA_\varepsilon B'B_\varepsilon$ (where $A_\varepsilon$ carries the maximally mixed state, the Gibbs state of a trivial Hamiltonian of $A_\varepsilon$) such that $\|\sigma_{AB'}-\sigma^\varepsilon_{AB'}\|_1<\varepsilon$ (approximation in trace norm) and $G_\varepsilon(\rho_{AB}\otimes\omega^\varepsilon_{A_\varepsilon B_\varepsilon})=\sigma^\varepsilon_{AB'}\otimes\omega^\varepsilon_{A_\varepsilon B_\varepsilon}$ if
\begin{align}
F_{t,\tau}(\rho_{AB},\gamma_A)&\geq F_{t,\tau}(\sigma_{AB'},\gamma_A)\qquad\forall(t,\tau)\in\mc D_{\rm bulk},\\
F_{-\infty,\tau}(\rho_{AB},\gamma_A)&\geq F_{-\infty,\tau}(\sigma_{AB'},\gamma_A)\qquad\forall\,\tau\in\mc D_{-\infty}.
\end{align}
\end{corollary}

\section{Conclusion}
We derive the general form of conditional entropy and show that any quantity satisfying operationally relevant axioms admits a representation as an exponential average of R\'enyi entropies of the conditioned distribution, parameterized by a real parameter and a probability measure. 
These results reveal an underlying structure common to the various conditional entropies introduced in the literature and clarify the relationships among them.




We derived conditions for the transformation of joint probability distributions under conditionally mixing channels, both in the many-copy regime and in the presence of a catalyst that may be borrowed during the process and returned unchanged. We showed that these conditions can be expressed in terms of conditional entropies
and that they can be used to characterize the optimal conversion rate under conditionally mixing channels and to derive a set of second laws of quantum thermodynamics with side information.

While the general form of conditional entropies can be expressed in terms of a real parameter and a probability measure, not all choices of these parameters give rise to quantities that satisfy the full set of axioms. 
We therefore provide a set of sufficient conditions on the parameters ensuring that the resulting conditional entropy is well defined and satisfies all required axioms. 
Moreover, we derive necessary conditions which, for suitably well-behaved distributions—and in particular for discrete probability measures with finite support—yield a complete characterization of the admissible parameter range.
While we conjecture that the sufficient conditions we derive are also necessary in full generality, a complete characterization of the necessary conditions for arbitrary parameter choices remains an open problem and an interesting problem for future research.




In Theorem~\ref{th: large sample and catalytic}, we provided a set of sufficient conditions for many-copy and catalytic transformations, which include the condition
\begin{align}
\label{eq: support restriction}
    H_{+\infty,0}(P_{XY}) < H_{+\infty,0}(Q_{X'Y'}) .
\end{align}
This condition requires the maximal support of the columns of $P_{XY}$ to be strictly smaller than that of $Q_{X'Y'}$ and therefore excludes the special case in which the two distributions have equal support.
As a consequence, the applicability of this result is restricted to certain classes of distributions.
 This restriction likely stems from the fact that we do not consider conditional R\'enyi entropies with negative $\alpha$. This is a well-known fact, which also holds in the absence of side information (see~\cite{Jensen_Kjaerulf_2019,farooq2024matrix}).
 Informally, this can be understood by considering the case of a single $Y$ outcome and observing that, under the support assumption~\eqref{eq: support restriction}, the R\'enyi entropies of $P_{XY}$ with negative $\alpha$ take the value $-\infty$, whereas they remain finite for $Q_{XY}$, since $P_{XY}$ must be embedded into the space of $Q_{XY}$. 
Consequently, the support condition already enforces a strict ordering for all negative values of $\alpha$, and these entropies therefore do not need to be included among the sufficient conditions.
This fact is reflected later when deriving sufficient conditions for catalytic conditionally (gibbs-)preserving operations in Section~\ref{sec: second laws of thermo}, where we imposed the support-size separation in Theorem \ref{thm:ConservedClassical} and Corollary \ref{cor:secondlaws}.
Hence, an open question is to extend these results to a more general setting without relying on the support-size assumption.
Another open question is whether the assumption that the Gibbs state has rational entries can be dropped. We believe this to be possible, similarly to what is known in the unconditional case~\cite{farooq2024matrix}. However, since the methods of this manuscript rely on the embedding map, our current construction does not cover the case of irrational entries.


In Theorem~\ref{thm: rate} we derived the optimal rate $R$ for the transformation of multiple copies of a joint probability $P_{XY}$ into $Q_{XY}$. If our conjecture regarding the characterization of $\mc D_{\rm bulk}$ and $\mc D_{-\infty}$ were correct, by a limit argument, the optimal rate could be written in the simpler form
\begin{equation}
R(P_{XY}\rightarrow Q_{X'Y'})=\inf_{(t,\tau)\in\mc D_{\rm bulk}}\frac{H_{t,\tau}(X'|Y')_Q}{H_{t,\tau}(X|Y)_P}\,.
\end{equation}
Indeed, we expect that the monotonicity region $\D_{-\infty}$ could be obtained as the limit $t\to-\infty$ of the monotonicity region $\D_{\rm bulk}$.


Finally, we expect that the conditions in Theorem~\ref{thm:ConservedClassical} and Corollary \ref{cor:secondlaws} concerning $I_{-\infty,\tau}$ and $F_{-\infty,\tau}$ could also be removed by a limiting argument.



\section{Aknowledgements}
We thank Frits Verhagen for helpful discussions. R.R.\ was supported by the National Research Foundation, Singapore, and A*STAR under its CQT Bridging Grant and acknowledges financial support from the ERC grant GIFNEQ 101163938. E.H.\ and M.T.\ acknowledge support from the National Research Foundation Investigatorship Award (NRF-NRFI10-2024-0006) and from the National Research Foundation, Singapore through the National Quantum Office, hosted in A*STAR, under its Centre for Quantum Technologies Funding Initiative (S24Q2d0009).



\newpage
\appendix

\section{Convexity and monotonicity}
\label{app: convexity and monotonicity}
In this appendix, we present several known equivalences between convexity and monotonicity under channels. These will be useful in the main text to characterize the parameter regimes for which the conditional entropy is monotone under conditionally mixing channels, and thus qualifies as a valid conditional entropy according to Definition~\ref{def: conditional entropy}. In particular, these equivalences allow us to reduce the question of monotonicity to that of convexity, which is often easier to analyze. 

In Definition~\ref{def: conditional entropy}, we introduced conditional entropies as quantities satisfying a collection of operational axioms, motivated by the requirement that they quantify conditional uncertainty. A central axiom is monotonicity under conditionally mixing channels.
Within the mathematical framework of preordered semirings, monotone homomorphisms and derivations are functions designed to satisfy closely analogous structural properties (see Section~\ref{sec: background semiring}). In particular, in the semiring associated with conditional majorization, these functions correspond—up to an additive constant and a logarithmic factor—to (extremal) conditional entropies.
Consequently, understanding their monotonicity properties is essential. In the following, we show that determining whether these functions are monotone under conditionally mixing channels is equivalent to studying their convexity properties.


We begin by recalling several notions of convexity for functions defined on
(unnormalized) joint distributions~$\mathcal V$, which—up to embedding in a larger space (padding with zeros)—
can be identified with the elements of the semiring of conditional majorization
introduced in Section~\ref{sec: semiring of conditional entropies}.
 We say that a map $F:\V\to\mb R$ is convex if, for any $P,Q\in\V$ of the same size and $t\in[0,1]$, we have
\begin{equation}
F\big(tP+(1-t)Q\big)\leq tF(P)+(1-t)F(Q).
\end{equation}
We say that $F$ is concave if the above inequality is reversed. We say that $F$ is quasi-convex if, for $P,Q\in\V$ of the same size and $t\in(0,1)$,
\begin{equation}
F\big(tP+(1-t)Q\big)\leq\max\{F(P),F(Q)\}.
\end{equation}
We say that $F$ is strongly quasi-concave if the inequality above is reversed, i.e.,
\begin{equation}
\label{eq: def strongly quasi-conc}
    F\big(tP+(1-t)Q\big)\geq\max\{F(P),F(Q)\}.
\end{equation}

 In the following, we show that the required convexity properties depend on the output space of the homomorphism. As discussed in Section~\ref{sec: sufficient conditions measure}, the key monotonicity is that under maps acting on the conditioning subsystem of a joint distribution, which form the key subclass of conditionally mixing channels. In matrix form, for a positive joint distribution $P_{XY}$, the action on the conditioning subsystem can be represented as $ P_{XY}\rightarrow P_{XY} D$, where $D$ is a row stochastic matrix.

\begin{proposition}\label{prop:convconc}
Let $S$ be the semiring of conditional majorization of Section~\ref{sec: semiring of conditional entropies}. Let
$\Phi: \cS \to \mathbb{K}$ be a monotone homomorphism, and hence of the form in Proposition~\eqref{prop: monotone homorphisms}. Here, $\mathbb{K} \in \{\mathbb{R}_+,  \mathbb{R}_+^{\mathrm{op}},\mathbb{TR}_+,\mathbb{TR}_+^{\mathrm{op}}\}$. Then, monotonicity with respect to maps acting on the conditioning system is equivalent to the following convexity properties, depending on the set $\mathbb{K}$:
\begin{itemize}[itemsep=0pt, topsep=5pt]
\item If $\mathbb{K} =\mathbb{R}_+$, $\Phi$ is monotone if and only if it is concave.
\item If $ \mathbb{K}= \mathbb{R}_+^{\mathrm{op}}$, $\Phi$ is monotone if and only if it is convex.
\item If $ \mathbb{K}=\mathbb{TR}_+$, $\Phi$ is monotone if and only if it is strongly quasi-concave.
\item If $ \mathbb{K}=\mathbb{TR}_+^{\mathrm{op}}$, $\Phi$ is monotone if and only if it is quasi-convex.
\end{itemize}
\end{proposition}

\begin{proof}
Let us first prove the `only if' part of the claim. The proof relies on associating the distributions $tP$ and $(1 - t)Q$ with distinct outcomes of the conditioning variable $Y$. Monotonicity under discarding outcomes of $Y$ then yields the desired inequality. Since monotone homomorphisms are invariant under embedding (see the discussion in
Section~\ref{sec: invariance embedding and shifting}), we may, in the following,
identify their evaluation on equivalence classes with their evaluation on
unnormalized joint distributions.


Explicitly, suppose that $P,Q\in\V$ and $t\in(0,1)$. We now have
\begin{equation}\label{eq:affinechain}
\left(\begin{array}{cc}
tP&0\\
0&(1-t)Q
\end{array}\right)\succeq\left(\begin{array}{cc}
tP&(1-t)Q\\
0&0
\end{array}\right)\succeq\left(\begin{array}{cc}
tP+(1-t)Q&0\\
0&0
\end{array}\right).
\end{equation}
The first implication follows by applying a permutation on the $X$ register conditioned on the second $Y$ outcome, while the second follows by discarding both outcomes and introducing a dummy outcome with zero probability.

In case $\mathbb{R}_+$, it now follows that
\begin{align}
\Phi\big(tP+(1-t)Q\big)&\geq\Phi\left[\left(\begin{array}{cc}
tP&0\\
0&(1-t)Q
\end{array}\right)\right]\\
&=\Phi(tP)+\Phi\big((1-t)Q\big)\\
&=\Phi(t)\Phi(P)+\Phi(1-t)\Phi(Q)\\
&=t\Phi(P)+(1-t)\Phi(Q).
\end{align}
Similarly, in case $\mathbb{R}_+^{\mathrm{op}}$, we get
\begin{align}
\Phi\big(tP+(1-t)Q\big)&\leq t\Phi(P)+(1-t)\Phi(Q) \,.
\end{align}
In case $\mathbb{TR}_+$, we get
\begin{align}
\Phi\big(tP+(1-t)Q\big)&\geq\max\{\Phi(P),\Phi(Q)\} \,.
\end{align}
In case $\mathbb{TR}^{\rm{op}}_+$, we have
\begin{align}
\Phi\big(tP+(1-t)Q\big)&\leq\Phi\left[\left(\begin{array}{cc}
tP&0\\
0&(1-t)Q
\end{array}\right)\right]\\
&=\max\{\Phi(tP),\Phi\big((1-t)Q\big)\}\\
&=\max\{\Phi(t)\Phi(P),\Phi(1-t)\Phi(Q)\}\\
&=\max\{\Phi(P),\Phi(Q)\}.
\end{align}

Next, we prove the “if” direction and assume that $\Phi$ is concave. Let $P \in \mathcal{V}$ and let $D$ be a row-stochastic matrix. We aim to show that $\Phi(PD) \leq \Phi(P)$. Let us first assume that $P=R=R_{X:YZ}$ (colon separating the primary $X$ from the conditioning $YZ$) and that $D$ is the map ignoring $Z$. %Just as in Lemma~\ref{lemma:shifting}, we see that
%\begin{equation}
%R\succeq\bigoplus_{z\in Z}R_Z(z)R_{XY|Z=z}\succeq R,
%\end{equation}
%so that the form given in \eqref{eq:FinalForm} tells us that
We have
\begin{align}
\Phi(R)&=\sum_{z\in\mc Z}R_Z(z)\Phi(R_{XY|Z=z})\leq\Phi\left(\sum_{z\in\mc Z}R_Z(z)R_{XY|Z=z}\right)=\Phi(R_{XY}),
\end{align}
where we have used the form of \eqref{eq:FinalForm} in the first equality and the assumed concavity of $\Phi$ in the inequality. This proves the first special case. Assume then the case $P_{XY}\succeq P_{XY}D=:Q_{XY'}$ where $D$ is a row-stochastic matrix which has only entries 0 or 1. This means that the matrix entries of $D$ are given by $D(y'|y)=\delta_{f(y),y'}$ for some function $f:Y\to Y'$. Let us define $R=R_{X:YY'}$ where $R(x,y,y')=P(x,y)$ whenever $f(y)=y'$ and $R(x,y,y')=0$ otherwise. It immediately follows that $P\approx R$ and $R_{XY'}=Q=PD$. Thus, using the above result, we have
\begin{align}
\Phi(PD)=\Phi(R_{XY'})\geq\Phi(R)=\Phi(P).
\end{align}
Finally, we prove the case of a general row-stochastic matrix $D$. Since $D$ can be expressed as a convex combination of row-stochastic matrices $D_k$ whose entries are either 0 or 1, there exist weights $t_k > 0$ with $\sum_k t_k = 1$
such that
\begin{equation}
    D = \sum_k t_k D_k.
\end{equation}
Our observations thus far, together with the concavity of $\Phi$ imply
\begin{align}\label{eq:apuconv}
\Phi(PD)&=\Phi\left(\sum_k t_kPD_k\right)\geq\sum_k t_k\Phi(PD_k)=\sum_k t_k\Phi(P)=\Phi(P).
\end{align}
The other cases are similar.
\end{proof}

A similar argument applies to derivations, which must be concave.

\begin{lemma}
\label{lem: conv derivations}
Let $S$ be the semiring of conditional majorization of Section~\ref{sec: semiring of conditional entropies}. Let $\Delta: \mathcal{S} \to \mathbb{R}$ be an extremal monotone derivation at $\|\cdot\|$, and hence of the form in Proposition~\ref{prop: monotone derivations}.
Then, $\Delta$ is a monotone derivation if and only if it is concave.
\end{lemma}

\begin{proof}
By assumption, $\Delta:\cS\to\mb R$ is a derivation at $\|\cdot\|$ of the form in Proposition~\ref{prop: monotone derivations}. In particular, it satisfies $\Delta(x)=0$ for all $x\in\mb R_+$. Let now $P,Q\in\V$ be of the same size and $t\in[0,1]$. Using the same argument in the proof of Lemma~\ref{prop:convconc}, and in particular the relation in equation~\eqref{eq:affinechain}, we have
\begin{align}
\Delta\big(tP+(1-t)Q\big)&\geq\Delta(tP)+\Delta\big((1-t)Q\big)\\
&=\underbrace{\Delta(t)}_{=0}\|P\|+t\Delta(P)+\underbrace{\Delta(1-t)}_{=0}\|Q\|+(1-t)\Delta(Q)\\
&=t\Delta(P)+(1-t)\Delta(Q).
\end{align}

For the reverse implication, the monotonicity of $\Delta$ can be shown in exactly the same way as in the second part of the proof of Lemma~\ref{prop:convconc}.
\end{proof}

The next lemma shows that, for homomorphisms, monotonicity under arbitrary maps acting on the conditioning system already implies monotonicity under the full class of conditionally mixing channels. This follows from the additional facts that homomorphisms are monotone under doubly stochastic maps and additive under direct sums. Together, these properties suffice to establish monotonicity under general conditionally mixing channels. 
\begin{lemma}
\label{lem: extension to general channels}
Let $S$ be the semiring of conditional majorization of Section~\ref{sec: semiring of conditional entropies}.
Let $\Phi: \cS \to \mathbb{K}$ be a monotone homomorphism, and hence of the form in Proposition~\ref{prop: monotone homorphisms}. If $\Phi$ is monotone under row stochastic maps acting on the conditioning system, then it is monotone under conditionally mixing operations.
\end{lemma}

\begin{proof}
We present the proof for the case $\mathbb{K} = \mathbb{R}_+$, as the other cases follow analogously.

We first note that, as discussed in Section~\ref{sec: monotonicty doubly stochastic},
monotonicity under doubly stochastic maps is guaranteed by the explicit form of
the homomorphisms derived in Proposition~\ref{prop: monotone homorphisms}, since
these are given by exponential expectations of R\'enyi entropies, which are
monotone under such maps.


To establish monotonicity under conditionally mixing channels, we introduce an auxiliary register that separates the individual terms appearing in the sum defining the channel, and then exploit the direct-sum property of the homomorphisms. The proof then follows by applying monotonicity under doubly stochastic maps and row-stochastic maps to each term separately.
Explicitly, let $P_{XY}\in\V$, $S^{(i)}$ be doubly stochastic, and $D^{(i)}$ have non-negative entries ($i=1,\ldots,k$) with $D^{(1)}+\cdots+D^{(k)}$ row-stochastic. Let us denote the output of the conditionally mixing channel $Q_{X'Y'}=\sum_i S^{(i)}P_{XY}D^{(i)}$. Let us define new conditioning systems $\widetilde{Y}:=Y\times\{1,\ldots,n\}$ and $\widetilde{Y}':=Y'\times\{1,\ldots,n\}$. 
Let $\pi_{Y'|\widetilde{Y}'}$ be the row-stochastic matrix which ignores the part $\{1,\ldots,n\}$ in $\widetilde{Y}'$. Let us also define the matrices
\begin{align}
P_{X\widetilde{Y}}&:=\big(P_{XY}D^{(1)}\,\cdots\,P_{XY}D^{(n)}\big)\\
Q_{X'\widetilde{Y}'}&:=\big(S^{(1)}P_{XY}D^{(1)}\,\cdots\,S^{(n)}P_{XY}D^{(n)}\big).
\end{align}
Because of the form of \eqref{eq:FinalForm}, %shifting Lemma in~\ref{lemma:shifting},
we have
\begin{equation}
\Phi(P_{X\widetilde{Y}})=\Phi\left(\bigoplus_{i=1}^n P_{XY}D^{(i)}\right),\quad\Phi(Q_{X'\widetilde{Y}'})=\Phi\left(\bigoplus_{i=1}^n S^{(i)}P_{XY}D^{(i)}\right).
\end{equation}
Putting all this together and noting that $(D^{(1)}\,\cdots D^{(n)})$ is row-stochastic and that $\bigoplus_{i=1}^n S^{(i)}$ is doubly stochastic, we have
\begin{align}
\Phi(P_{XY})&\leq\Phi\big[P_{XY}\big(D^{(1)}\,\cdots\,D^{(n)}\big)\big]\\
&=\Phi(P_{X\widetilde{Y}})=\Phi\left(\bigoplus_{i=1}^n P_{XY}D^{(i)}\right)\\
&\leq\Phi\left[\left(\bigoplus_{i=1}^n S^{(i)}\right)\left(\bigoplus_{i=1}^n P_{XY}D^{(i)}\right)\right]\\
&=\Phi\left(\bigoplus_{i=1}^n S^{(i)}P_{XY}D^{(i)}\right)\\
&=\Phi(Q_{X'\widetilde{Y}'})\\
&\leq\Phi(Q_{X'\widetilde{Y}'}\pi_{Y'|\widetilde{Y}'})\\
&=\Phi(Q_{X'Y'}),
\end{align}
showing the monotonicity in case $\mathbb{R}_+$. The other cases are essentially the same; the sums in \eqref{eq:apuconv} turn into maxima in the tropical cases $\mathbb{TR}_+$ and $\mathbb{TR}^{\mathrm{op}}_+$, and the inequalities are reversed in cases $\mathbb{R}^{\mathrm{op}}_+$ and $\mathbb{TR}^{\mathrm{op}}_+$.
\end{proof}
The same implication is also true for the derivations.
\begin{lemma}
Let $S$ be the semiring of conditional majorization of Section~\ref{sec: semiring of conditional entropies}. Let $\Delta: \mathcal{S} \to \mathbb{R}$ be an extremal monotone derivation at $\|\cdot\|$, and hence of the form in Proposition~\ref{prop: monotone derivations}.
If $\Delta$ is monotone under row stochastic maps acting on the conditioning system, then it is monotone under conditionally mixing operations.
\end{lemma}
\begin{proof}
    The proof is analogous to that of one of Lemma~\ref{lem: extension to general channels} for the homomorphisms.
\end{proof}

\bibliographystyle{ultimate}
\bibliography{bibliography}

\appendix

\end{document}
